% Leçon 9 — Vocabulaire et compréhension de texte : processus du recyclage — Chinois Tle A — Cours signé OnBuch+
% Programme officiel MINESEC. Compiler avec : tectonic ZH09-vocabulaire-processus-du-recyclage.tex
\newcommand{\DOCMATIERE}{Chinois}
\newcommand{\DOCNIVEAU}{Tle A}
\newcommand{\DOCMODULE}{Module 2 : Environnement et santé}
\newcommand{\DOCLECON}{ZH09}
\newcommand{\DOCTITRE}{Leçon 9 — Vocabulaire et compréhension de texte : processus du recyclage}
% OnBuch+ — préambule « cours signés » ENRICHI (compilé avec tectonic / XeLaTeX)
% Système visuel « Pop » d'OnBuch+ : fond crème (jamais blanc), bordures encre
% épaisses, ombres « dures » décalées, titres Archivo Black en capitales.
% Version MATHÉMATIQUES : graphes (pgfplots/tikz), boîtes pédagogiques
% (définition, propriété, méthode, attention, le savais-tu, exercice résolu,
% exercice, TP, bilan), page de garde et signature OnBuch+ ancrée en bas.
%
% Macros attendues dans main.tex AVANT \input{preamble} :
%   \DOCMATIERE  \DOCNIVEAU  \DOCMODULE  \DOCLECON (numéro)  \DOCTITRE
\documentclass[11pt]{article}

\usepackage{fontspec}
\usepackage[french]{babel} % français = langue principale ; le chinois via \zh{..} / textebox (xeCJK)
\usepackage{xeCJK}
\usepackage{pifont} % \ding{51} \ding{55} utilisés par les modèles
\usepackage[a4paper,margin=2cm,top=3.4cm,bottom=2.8cm,headheight=1.4cm,footskip=1.3cm]{geometry}
\usepackage{amsmath,amssymb,amsfonts}
\usepackage{mathtools} % \xmapsto, \coloneqq... souvent utilisés par les modèles
\usepackage{cancel} % simplifications barrées (bilans de forces)
% \abs{z} (module, valeur absolue) et \norm{u} (norme), utilisés par les modèles
\providecommand{\abs}{}\let\abs\relax\DeclarePairedDelimiter{\abs}{\lvert}{\rvert}
\providecommand{\norm}{}\let\norm\relax\DeclarePairedDelimiter{\norm}{\lVert}{\rVert}
\usepackage[table]{xcolor} % [table] : fournit aussi \rowcolors (alternance de lignes)
\usepackage{pagecolor}
\usepackage{tikz}
\usetikzlibrary{arrows.meta,positioning,calc,shapes.geometric,shapes.misc,shapes.arrows,decorations.pathmorphing,decorations.pathreplacing,decorations.markings,patterns,shadows,mindmap,trees,fit,backgrounds,matrix,intersections,angles,quotes}
\usepackage{pgfplots}
\usepackage[version=4]{mhchem} % équations nucléaires \ce{^{235}_{92}U}
\usepackage[european,siunitx]{circuitikz} % schémas électriques
\pgfplotsset{compat=1.18}
\usepgfplotslibrary{fillbetween,groupplots} % aires entre courbes (calcul intégral)
\usepackage{enumitem}
\usepackage{fancyhdr}
% [most] charge toutes les bibliothèques tcolorbox (listings, minted, raster,
% documentation...) et gonfle inutilement la mémoire TeX sur un document long
% (~300 boîtes) : on ne charge que ce qui est réellement utilisé.
\usepackage[skins,breakable]{tcolorbox}
\usepackage{graphicx}
\usepackage{colortbl}
\usepackage{booktabs}
\usepackage{tabularx}
\usepackage{diagbox} % case d'en-tête diagonale des tableaux à double entrée
% Colonnes extensibles centrée / gauche / droite (C, L, R), souvent utilisées par les modèles
\newcolumntype{C}{>{\centering\arraybackslash}X}
\newcolumntype{L}{>{\raggedright\arraybackslash}X}
\newcolumntype{R}{>{\raggedleft\arraybackslash}X}
\usepackage{array}
\usepackage{multicol}
\usepackage{multirow}
\usepackage{siunitx}
% parse-numbers=false : accepte aussi \SI{2,5 \times 10^{-3}}{...}
\sisetup{locale=FR,output-decimal-marker={,},group-minimum-digits=4,parse-numbers=false}
\DeclareSIUnit\ohm{\ensuremath{\symup{\Omega}}} % U+2126 absent de la police
\DeclareSIUnit\division{div} % réglages d'oscilloscope (ms/div, V/div)
\DeclareSIUnit\tr{tr} % tours (vitesse de rotation, ex. \SI{1500}{\tr\per\minute})
\DeclareSIUnit\molar{M} % concentration molaire (ex. \SI{3}{\molar}, \SI{1}{\micro\molar})
\DeclareSIUnit\an{an} % année (le modèle confond parfois avec \year, un compteur TeX) — ex. \SI{5}{\kilo\meter\per\an}
\DeclareSIUnit\FCFA{FCFA} % franc CFA (ex. \SI{500}{\FCFA\per\kilo\gram})
\DeclareSIUnit\degreeBrix{\textdegree Brix} % teneur en sucre d'un sirop/jus (réfractométrie)
\DeclareSIUnit\month{mois} % le modèle utilise parfois \month, un compteur TeX, comme unité de durée
\usepackage{qrcode}
% \degree : commande fréquemment utilisée par les modèles pour un angle ou une
% température en dehors de \SI{}{} (non fournie par les paquets chargés).
\providecommand{\degree}{^{\circ}}
% \qed (fin de démonstration), utilisé par les modèles sans amsthm
\providecommand{\qed}{\hfill$\square$}
% \barre : double barre des valeurs interdites dans un tableau de variations (tkz-tab)
\providecommand{\barre}{\ensuremath{\|}}
% environnement proof (amsthm non chargé) utilisé par les modèles
\ifdefined\proof\else\newenvironment{proof}[1][Démonstration]{\par\noindent\textit{#1.}\ }{\qed\par}\fi
\usepackage{lastpage}
\usepackage{hyperref}
\hypersetup{hidelinks}

% ============================================================
% Palette « Pop » OnBuch+ — mêmes hex que la classe P (plus_ui.dart)
% ============================================================
\definecolor{poporange}{HTML}{F59321}
\definecolor{popdark}{HTML}{E07A0C}
\definecolor{popcream}{HTML}{FFF6E8}
\definecolor{popink}{HTML}{1C1714}
\definecolor{poppurple}{HTML}{7A5AE0}
\definecolor{popgreen}{HTML}{2E9E6A}
\definecolor{poppink}{HTML}{E0457B}
\definecolor{popblue}{HTML}{2F7DE1}
\definecolor{popgold}{HTML}{C98A12}
\definecolor{popmuted}{HTML}{8A7F76}
\definecolor{popline}{HTML}{EADFD2}
\definecolor{popgray}{HTML}{8A7F76}
\colorlet{popgrey}{popgray}
% Alias tolérants (le modèle nomme parfois les couleurs différemment)
\colorlet{popwhite}{white}
\colorlet{popblack}{popink}
\colorlet{popred}{poppink}
\colorlet{popyellow}{popgold}
% Teintes claires (fonds de tableaux/encadrés) — le modèle nomme parfois
% les couleurs avec un suffixe L (ex. popblueL) sans les avoir définies.
\colorlet{poporangeL}{poporange!20}
\colorlet{popdarkL}{popdark!20}
\colorlet{poppurpleL}{poppurple!20}
\colorlet{popgreenL}{popgreen!20}
\colorlet{poppinkL}{poppink!20}
\colorlet{popblueL}{popblue!20}
\colorlet{popgoldL}{popgold!20}
\colorlet{popredL}{popred!20}
\colorlet{popgrayL}{popgray!20}
% Teintes claires pour fonds de tableaux / zones
\colorlet{popblueL}{popblue!10!white}
\colorlet{poporangeL}{poporange!14!white}
\colorlet{popgreenL}{popgreen!12!white}
\colorlet{poppurpleL}{poppurple!12!white}
\colorlet{poppinkL}{poppink!12!white}
\colorlet{popgoldL}{popgold!14!white}

\pagecolor{popcream}

% ============================================================
% Typographie
% ============================================================
\defaultfontfeatures{Ligatures=TeX}
\setmainfont{PlusJakartaSans-Medium.ttf}[Path=../../../fonts/,BoldFont=PlusJakartaSans-Bold.ttf]
% Symboles, API et emojis absents de la police principale : repli sur DejaVu Sans (caractères actifs)
\newfontface\symfont{DejaVuSans.ttf}[Path=../../../fonts/]
\makeatletter
\newcommand\symactive[1]{\catcode#1=13 \begingroup\lccode`\~=#1 \lowercase{\endgroup\def~}{{\symfont\char#1}}}
\newcommand\symblank[1]{\catcode#1=13 \begingroup\lccode`\~=#1 \lowercase{\endgroup\def~}{}}
\newcommand\symhyph[1]{\catcode#1=13 \begingroup\lccode`\~=#1 \lowercase{\endgroup\def~}{\mbox{-}}}
\newcount\symc
\newcommand\symrange[2]{\symc=#1 \@whilenum\symc<#2 \do{\expandafter\symactive\expandafter{\the\symc}\advance\symc by 1 }}
\newcommand\symblankrange[2]{\symc=#1 \@whilenum\symc<#2 \do{\expandafter\symblank\expandafter{\the\symc}\advance\symc by 1 }}
\makeatother
\symrange{"0250}{"02FF}\symactive{"2713}\symactive{"2714}\symactive{"2717}\symactive{"2718}\symactive{"2610}\symactive{"2611}\symactive{"2612}
\symrange{"2600}{"27BF}
\catcode"2705=13 \begingroup\lccode`\~="2705 \lowercase{\endgroup\def~}{{\symfont\char"2713}}\symblank{"26BD}\symblank{"26A1}
\symrange{"2460}{"257F}\symactive{"223C}\symactive{"2248}\symactive{"2260}
\symhyph{"2011}\symhyph{"2E1A}\symblankrange{"1F300}{"1FAFF}\symblank{"FE0F}
% Caractères chinois : WenQuanYi Zen Hei (sous-ensemble GB2312 + pinyin), gras/italique simulés
\setCJKmainfont{WenQuanYiZenHei-subset.ttf}[Path=../../../fonts/,AutoFakeBold=3,AutoFakeSlant=0.2]
\xeCJKsetup{CJKspace=false}
% Pinyin : voyelles à caron (ǎ ǐ ǒ ǔ ǖ ǘ ǚ ǜ) absentes de la police principale -> caractères actifs
\newfontface\pyfont{WenQuanYiZenHei-subset.ttf}[Path=../../../fonts/]
\newcommand\pyactive[1]{\catcode#1=13 \begingroup\lccode`\~=#1 \lowercase{\endgroup\def~}{{\pyfont\char#1}}}
\pyactive{"01CD}\pyactive{"01CE}\pyactive{"01CF}\pyactive{"01D0}\pyactive{"01D1}\pyactive{"01D2}\pyactive{"01D3}\pyactive{"01D4}\pyactive{"01D5}\pyactive{"01D6}\pyactive{"01D7}\pyactive{"01D8}\pyactive{"01D9}\pyactive{"01DA}\pyactive{"01DB}\pyactive{"01DC}
% Maths en sans-serif (Fira Math) avec chiffres et lettres droites de Plus
% Jakarta, pour que les formules se fondent dans le texte.
\usepackage[math-style=ISO,bold-style=ISO]{unicode-math}
\setmathfont{FiraMath-Regular.otf}[Path=../../../fonts/]
\setmathfont{PlusJakartaSans-Medium.ttf}[Path=../../../fonts/,range={up/num,up/{Latin,latin},\mathup}]
\usepackage{icomma} % virgule décimale sans espace en mode math
% Caractères mathématiques Unicode absents des polices de texte (Plus Jakarta,
% Archivo Black) : rendus par leur équivalent mathématique, en texte comme en
% formule — sinon ils s'affichent en carré vide (ex. « Arithmétique dans ℤ »).
\usepackage{newunicodechar}
\newunicodechar{ℝ}{\ensuremath{\mathbb{R}}}
\newunicodechar{ℤ}{\ensuremath{\mathbb{Z}}}
\newunicodechar{ℂ}{\ensuremath{\mathbb{C}}}
\newunicodechar{ℕ}{\ensuremath{\mathbb{N}}}
\newunicodechar{ℚ}{\ensuremath{\mathbb{Q}}}
\newunicodechar{𝔻}{\ensuremath{\mathbb{D}}}
\newunicodechar{α}{\ensuremath{\alpha}}
\newunicodechar{β}{\ensuremath{\beta}}
\newunicodechar{γ}{\ensuremath{\gamma}}
\newunicodechar{δ}{\ensuremath{\delta}}
\newunicodechar{ε}{\ensuremath{\varepsilon}}
\newunicodechar{θ}{\ensuremath{\theta}}
\newunicodechar{λ}{\ensuremath{\lambda}}
\newunicodechar{μ}{\ensuremath{\mu}}
\newunicodechar{σ}{\ensuremath{\sigma}}
\newunicodechar{τ}{\ensuremath{\tau}}
\newunicodechar{φ}{\ensuremath{\varphi}}
\newunicodechar{ψ}{\ensuremath{\psi}}
\newunicodechar{ω}{\ensuremath{\omega}}
\newunicodechar{Σ}{\ensuremath{\Sigma}}
% majuscules produites par \MakeUppercase dans les titres (α-aminés -> Α)
\newunicodechar{Α}{\ensuremath{\alpha}}
\newunicodechar{Β}{\ensuremath{\beta}}
\newunicodechar{Γ}{\ensuremath{\Gamma}}
\newunicodechar{Δ}{\ensuremath{\Delta}}
\newunicodechar{Ε}{\ensuremath{\varepsilon}}
\newunicodechar{Θ}{\ensuremath{\Theta}}
\newunicodechar{Λ}{\ensuremath{\Lambda}}
\newunicodechar{Μ}{\ensuremath{\mu}}
\newunicodechar{Τ}{\ensuremath{\tau}}
\newunicodechar{Φ}{\ensuremath{\Phi}}
\newunicodechar{Ψ}{\ensuremath{\Psi}}
\newunicodechar{Ω}{\ensuremath{\Omega}}
\newunicodechar{↦}{\ensuremath{\mapsto}}
\newunicodechar{∈}{\ensuremath{\in}}
\newunicodechar{∉}{\ensuremath{\notin}}
\newunicodechar{∧}{\ensuremath{\wedge}}
\newunicodechar{⇔}{\ensuremath{\Leftrightarrow}}
\newunicodechar{⟺}{\ensuremath{\Longleftrightarrow}}
\newunicodechar{⇒}{\ensuremath{\Rightarrow}}
\newunicodechar{⟹}{\ensuremath{\Longrightarrow}}
\newunicodechar{∘}{\ensuremath{\circ}}
\newunicodechar{∪}{\ensuremath{\cup}}
\newunicodechar{∩}{\ensuremath{\cap}}
\newunicodechar{≡}{\ensuremath{\equiv}}
\newunicodechar{⊂}{\ensuremath{\subset}}
\newunicodechar{⊥}{\ensuremath{\perp}}
\newunicodechar{∃}{\ensuremath{\exists}}
\newunicodechar{∀}{\ensuremath{\forall}}
\newunicodechar{∛}{\ensuremath{\sqrt[3]{\,}}}
\newunicodechar{∜}{\ensuremath{\sqrt[4]{\,}}}
\newunicodechar{⃗}{\ensuremath{{}^{\to}}}
\newunicodechar{ⁿ}{\ifmmode^{n}\else\textsuperscript{\ensuremath{n}}\fi}
\newunicodechar{ˣ}{\ifmmode^{x}\else\textsuperscript{\ensuremath{x}}\fi}
\newunicodechar{ᵇ}{\ifmmode^{b}\else\textsuperscript{\ensuremath{b}}\fi}
\newunicodechar{ʳ}{\ifmmode^{r}\else\textsuperscript{\ensuremath{r}}\fi}
\newunicodechar{ᵘ}{\ifmmode^{u}\else\textsuperscript{\ensuremath{u}}\fi}
\newunicodechar{ʸ}{\ifmmode^{y}\else\textsuperscript{\ensuremath{y}}\fi}
\newunicodechar{ᵃ}{\ifmmode^{a}\else\textsuperscript{\ensuremath{a}}\fi}
\newunicodechar{ᵉ}{\ifmmode^{e}\else\textsuperscript{\ensuremath{e}}\fi}
\newunicodechar{ᵏ}{\ifmmode^{k}\else\textsuperscript{\ensuremath{k}}\fi}
\newunicodechar{ᵖ}{\ifmmode^{p}\else\textsuperscript{\ensuremath{p}}\fi}
\newunicodechar{ᴺ}{\ifmmode^{N}\else\textsuperscript{\ensuremath{N}}\fi}
\newunicodechar{ᶜ}{\ifmmode^{c}\else\textsuperscript{\ensuremath{c}}\fi}
\newunicodechar{ʰ}{\ifmmode^{h}\else\textsuperscript{\ensuremath{h}}\fi}
\newunicodechar{ᵠ}{\ifmmode^{q}\else\textsuperscript{\ensuremath{q}}\fi}
\newunicodechar{ₙ}{\ifmmode_{n}\else\textsubscript{\ensuremath{n}}\fi}
\newunicodechar{ₐ}{\ifmmode_{a}\else\textsubscript{\ensuremath{a}}\fi}
\newunicodechar{ᵢ}{\ifmmode_{i}\else\textsubscript{\ensuremath{i}}\fi}
\newunicodechar{ᵣ}{\ifmmode_{r}\else\textsubscript{\ensuremath{r}}\fi}
\newunicodechar{ₖ}{\ifmmode_{k}\else\textsubscript{\ensuremath{k}}\fi}
\newunicodechar{ᵦ}{\ifmmode_{\beta}\else\textsubscript{\ensuremath{\beta}}\fi}
\newunicodechar{ₓ}{\ifmmode_{x}\else\textsubscript{\ensuremath{x}}\fi}
\newfontfamily\popdisplay{ArchivoBlack-Regular.ttf}[Path=../../../fonts/]
\newfontfamily\poplogo{SpaceGrotesk-Bold.ttf}[Path=../../../fonts/]
\newfontfamily\popsemi{PlusJakartaSans-SemiBold.ttf}[Path=../../../fonts/]
\newfontfamily\popxbold{PlusJakartaSans-ExtraBold.ttf}[Path=../../../fonts/]
\newfontfamily\popxboldls{PlusJakartaSans-ExtraBold.ttf}[Path=../../../fonts/,LetterSpace=3.0]

\setlist[enumerate]{leftmargin=1.7em,itemsep=5pt,topsep=4pt}
\setlist[itemize]{leftmargin=1.7em,itemsep=4pt,topsep=4pt,label={\color{poporange}\textbullet}}
\setlength{\parindent}{0pt}
\setlength{\parskip}{5pt}
\renewcommand{\arraystretch}{1.25}

% pgfplots : style Pop par défaut
\pgfplotsset{
  every axis/.append style={axis line style={popink,line width=1pt},tick style={popink},
    grid style={popline,line width=0.5pt},label style={font=\small},tick label style={font=\footnotesize}},
  popcurve/.style={poporange,line width=1.8pt,no markers,smooth},
}
% Même style disponible dans un \draw TikZ simple (hors pgfplots)
\tikzset{popcurve/.style={poporange,line width=1.8pt}}
\tikzset{popgrid/.style={popline,very thin}}
\colorlet{popcurve}{poporange} % le nom du style est parfois utilisé comme couleur

% ============================================================
% Pill / squiggle
% ============================================================
\newcommand{\poppill}[3]{%
  \tikz[baseline=-0.55ex]\node[draw=popink,line width=1.2pt,rounded corners=2.6pt,%
    fill=#2,inner xsep=6.5pt,inner ysep=3pt]{%
    \popxboldls\fontsize{7}{7}\selectfont\color{#3}\MakeUppercase{#1}};%
}
\newcommand{\popsquiggle}[1]{%
  \tikz[baseline=-0.4ex]{
    \draw[#1,line width=2.6pt,line cap=round]
      plot [smooth] coordinates {(0,0.09) (0.34,0.01) (0.68,0.21) (1.02,0.01) (1.36,0.10)};
  }%
}
% Mot-clé surligné (marqueur orange sous le texte)
\newcommand{\cle}[1]{{\popxbold\color{popdark}#1}}

% ============================================================
% En-tête / pied
% ============================================================
\pagestyle{fancy}
\fancyhf{}
\renewcommand{\headrulewidth}{0pt}
\renewcommand{\footrulewidth}{0pt}
\lhead{\raisebox{-0.15ex}{\poplogo\fontsize{13}{13}\selectfont\color{popink}OnBuch\color{poporange}+}}
\rhead{\poppill{\DOCMATIERE\ \textbullet\ \DOCNIVEAU\ \textbullet\ Leçon \DOCLECON}{white}{popink}}
\lfoot{\popxbold\fontsize{7.6}{7.6}\selectfont\color{popmuted}\MakeUppercase{Cours signé OnBuch+}\ \textbullet\ {\popsemi\DOCTITRE}}
\rfoot{\tikz[baseline=-0.6ex]\node[rounded corners=8.5pt,fill=popink,minimum height=17pt,minimum width=17pt,inner xsep=5pt,inner sep=0]{\popxbold\fontsize{8}{8}\selectfont\color{popcream}\thepage\,/\,\pageref*{LastPage}};}

% ============================================================
% Titres
% ============================================================
\newcommand{\chaptitre}[2]{%
  \begin{tcolorbox}[enhanced,colback=poporange,colframe=popink,boxrule=2.6pt,arc=11pt,
      drop shadow=popink,left=15pt,right=15pt,top=12pt,bottom=11pt,before skip=3pt,after skip=18pt]
    \poppill{#2}{popink}{popcream}\par\vspace{9pt}
    {\raggedright\hyphenpenalty=10000\exhyphenpenalty=10000\popdisplay\fontsize{22}{24}\selectfont\color{popcream}\MakeUppercase{#1}\par}
  \end{tcolorbox}
}
% Section numérotée : pastille orange avec le numéro + titre Archivo
\newcounter{popsec}
\newcommand{\coursec}[1]{%
  \stepcounter{popsec}\setcounter{popsub}{0}%
  \par\vspace{16pt}\noindent
  \begin{minipage}{\linewidth}
  \tikz[baseline=-0.7ex]\node[draw=popink,line width=1.6pt,fill=poporange,rounded corners=4pt,minimum size=22pt,inner sep=0]{\popdisplay\fontsize{12}{12}\selectfont\color{popcream}\thepopsec};%
  \hspace{8pt}{\popdisplay\fontsize{15}{17}\selectfont\color{popink}\MakeUppercase{#1}}\par
  \vspace{4pt}\noindent{\color{poporange}\rule{36pt}{3.6pt}}
  \end{minipage}\par\nopagebreak\vspace{8pt}
}
\newcounter{popsub}
\newcommand{\courssub}[1]{%
  \stepcounter{popsub}%
  \par\vspace{9pt}\noindent{\popxbold\fontsize{11.5}{13}\selectfont\color{popink}{\color{poporange}\thepopsec.\thepopsub}\hspace{6pt}#1}\par\nopagebreak\vspace{4pt}
}

% ============================================================
% Boîtes pédagogiques — même recette : fond blanc, bordure épaisse colorée,
% ombre dure encre, badge plein. Titre optionnel [#1].
% ============================================================
\tcbset{popbase/.style={breakable,enhanced,colback=white,boxrule=2.3pt,arc=9pt,
  shadow={3pt}{-3pt}{0pt}{popink},left=14pt,right=14pt,top=11pt,bottom=11pt,before skip=11pt,after skip=15pt,
  fontupper=\popsemi\fontsize{10.6}{14.5}\selectfont\color{popink},
  fontlower=\popsemi\fontsize{10.6}{14.5}\selectfont\color{popink}}}
% Coche dessinée (le glyphe ✓ n'existe pas dans les polices du document)
\newcommand{\popcheck}[1]{\tikz[baseline=-0.5ex]\draw[#1,line width=1.6pt,line cap=round,line join=round](-0.13,0.0)--(-0.04,-0.09)--(0.13,0.1);}
\providecommand{\coche}{\popcheck{popgreen}}
\providecommand{\resultat}[1]{\par\smallskip\noindent\fcolorbox{popgreen}{popgreenL}{\parbox{\dimexpr\linewidth-2\fboxsep-2\fboxrule\relax}{\textbf{Résultat.} #1}}\par\smallskip}
\newcommand{\popboxhead}[3]{\poppill{#1}{#2}{white}\ifx\relax#3\relax\else\hspace{6pt}{\popxbold\fontsize{10.6}{12}\selectfont\color{popink}#3}\fi\par\vspace{7pt}}

\newtcolorbox{aretenir}[1][]{popbase,colframe=popgreen,before upper={\popboxhead{À retenir}{popgreen}{#1}}}
% Texte en chinois (document de lecture, dialogue, extrait) : cadre
% d'étude. Un titre optionnel (source, auteur) s'affiche dans l'en-tête.
\newtcolorbox{textebox}[1][]{popbase,colframe=popblue,before upper={\popboxhead{文本 · Texte}{popblue}{#1}},after upper={}}
% Marques ✓ / ✗ (forme correcte / fautive) souvent utilisées par les modèles
\providecommand{\cross}{\textcolor{poppink}{\ensuremath{\times}}}
\providecommand{\xmark}{\cross}\providecommand{\crossmark}{\cross}\providecommand{\cmark}{\ensuremath{\checkmark}}
% Ligne de réponse / justification à compléter (inventée par les modèles)
\providecommand{\justification}[1]{\par\noindent\textit{Justification~:} \dotfill\par}
% \textipa (package tipa non chargé) : la police ne couvre pas l'API -> simple passage
\providecommand{\textipa}[1]{#1}
% Soulignement (ulem non chargé) : \uline, \ul
\providecommand{\uline}[1]{\underline{#1}}\providecommand{\ul}[1]{\underline{#1}}
% \glossaire{\item ...} inventée par les modèles : liste de glossaire
\providecommand{\glossaire}[1]{\begin{itemize}#1\end{itemize}}
% \alert (beamer) inventée par les modèles : texte mis en évidence
\providecommand{\alert}[1]{\textbf{\textcolor{poppink}{#1}}}
% variantes ASCII d'ordinaux écrites en macro
\providecommand{\iere}{\textsuperscript{re}}\providecommand{\ieme}{\textsuperscript{e}}\providecommand{\ere}{\textsuperscript{re}}\providecommand{\eme}{\textsuperscript{e}}\providecommand{\er}{\textsuperscript{er}}\providecommand{\ie}{\textsuperscript{e}}
% Ordinaux français écrits en macro par les modèles : 1\ière, 3\ième
\newcommand{\ière}{\textsuperscript{re}}\newcommand{\ième}{\textsuperscript{e}}
% \exemple (inventée par les modèles) : étiquette « Exemple : »
\providecommand{\exemple}{\textbf{Exemple~:}\ }
% \square utilisable aussi hors mode math (cases à cocher)
\AtBeginDocument{\let\squareMath\square\renewcommand{\square}{\ensuremath{\squareMath}}}
% Mot ou phrase en chinois à l'intérieur d'un paragraphe français
\newcommand{\zh}[1]{\ifmmode\text{#1}\else{#1}\fi}
\let\eng\zh \let\esp\zh \let\all\zh % alias tolérés
% flèches d'intonation inventées par les modèles
\providecommand{\textsearrow}{}\renewcommand{\textsearrow}{\ensuremath{\searrow}}\providecommand{\textnearrow}{}\renewcommand{\textnearrow}{\ensuremath{\nearrow}}
% \fr{..} : mot français (inventé par les modèles)
\providecommand{\fr}[1]{\textit{#1}}
% zhbox : encadré de texte chinois inventé par les modèles
\newenvironment{zhbox}{\begin{quote}}{\end{quote}}
% \vs : « versus » inventé par les modèles
\providecommand{\vs}{\textit{vs}\ }
% \blank : case à compléter inventée par les modèles
\providecommand{\blank}[1][]{\underline{\hspace{1.6em}}}
\newtcolorbox{exemplebox}[1][]{popbase,colframe=poporange,before upper={\popboxhead{Exemple}{poporange}{#1}}}
\newtcolorbox{definition}[1][]{popbase,colframe=popblue,before upper={\popboxhead{Définition}{popblue}{#1}}}
\newtcolorbox{propriete}[1][]{popbase,colframe=poppurple,before upper={\popboxhead{Propriété}{poppurple}{#1}}}
\newtcolorbox{methode}[1][]{popbase,colframe=popgold,before upper={\popboxhead{Méthode}{popgold}{#1}}}
\newtcolorbox{attention}[1][]{popbase,colframe=poppink,before upper={\popboxhead{Attention}{poppink}{#1}}}
\newtcolorbox{experience}[1][]{popbase,colframe=popblue,colback=popblueL,before upper={\popboxhead{Expérience}{popblue}{#1}}}
% « Le savais-tu ? » : carte violette pleine (contexte, culture, Cameroun)
\newtcolorbox{savaistu}[1][]{popbase,colback=poppurple,colframe=popink,
  fontupper=\popsemi\fontsize{10.4}{14.2}\selectfont\color{white},
  before upper={\poppill{Le savais-tu\,?}{white}{poppurple}\ifx\relax#1\relax\else\hspace{6pt}{\popxbold\color{white}#1}\fi\par\vspace{7pt}}}
% Exercice résolu : énoncé en haut, \tcblower puis la solution sur fond crème
\newcounter{popexo}\newcounter{popexr}
\newtcolorbox{exoresolu}[1][]{popbase,colframe=poporange,colbacklower=popcream,
  segmentation style={popink,line width=1.2pt,dashed},
  before upper={\stepcounter{popexr}\popboxhead{Exercice résolu \thepopexr}{poporange}{#1}},
  before lower={\poppill{Solution}{popink}{popcream}\par\vspace{6pt}}}
% Exercice (non résolu en place) — la correction est dans la partie Corrigés
\newtcolorbox{exercice}[1][]{popbase,colframe=popink,
  before upper={\stepcounter{popexo}\popboxhead{Exercice \thepopexo}{popink}{#1}}}
% Corrigé
\newtcolorbox{corrige}[1][]{popbase,colframe=popgreen,colback=popgreenL,
  before upper={\popboxhead{Corrigé}{popgreen}{#1}}}
% Objectifs / prérequis / bilan
\newtcolorbox{objectifs}{popbase,colframe=popink,colback=poporangeL,
  before upper={\popboxhead{Objectifs de la leçon}{popink}{}}}
\newtcolorbox{prerequis}{popbase,colframe=popink,colback=popblueL,
  before upper={\popboxhead{Prérequis}{popblue}{}}}
\newtcolorbox{bilan}{popbase,colframe=popink,colback=white,boxrule=2.8pt,
  before upper={\popboxhead{Fiche bilan}{poporange}{L'essentiel à connaître pour l'examen}}}
% Figure encadrée avec légende
\newtcolorbox{popfigure}[1][]{enhanced,breakable=false,colback=white,colframe=popline,boxrule=1.4pt,arc=8pt,
  left=10pt,right=10pt,top=10pt,bottom=8pt,before skip=10pt,after skip=12pt,halign=center,
  lower separated=false,halign lower=center,
  fontlower=\popsemi\fontsize{9}{11.5}\selectfont\color{popmuted},
  #1}
\newcommand{\legende}[1]{\par\vspace{6pt}{\popsemi\fontsize{9}{11.5}\selectfont\color{popmuted}#1}}

% ============================================================
% Signature OnBuch+
% ============================================================
\newcommand{\poplogoinline}[1]{{\poplogo\fontsize{#1}{#1}\selectfont\color{popink}OnBuch\color{poporange}+}}
\newcommand{\signatureonbuch}{%
  \begin{tcolorbox}[enhanced,colback=popink,colframe=popink,boxrule=2.2pt,arc=10pt,
      drop shadow=poporange,left=14pt,right=14pt,top=10pt,bottom=10pt,before skip=0pt,after skip=0pt]
    \tikz[baseline=-0.7ex]\node[circle,fill=popgreen,draw=popcream,line width=1.3pt,minimum size=22pt,inner sep=0]{\popcheck{white}};%
    \hspace{9pt}\begin{minipage}[c]{0.62\linewidth}
      {\popxbold\fontsize{12}{13}\selectfont\color{popcream}Signé\ \textbullet\ {\poplogo OnBuch\color{poporange}+}}\par\vspace{2pt}
      {\raggedright\popsemi\fontsize{8}{10.5}\selectfont\color{popline}\MakeUppercase{Équipe pédagogique OnBuch+}\\\MakeUppercase{Conforme au programme officiel MINESEC}\par}
    \end{minipage}\hfill
    {\poplogo\fontsize{22}{22}\selectfont\color{popcream}OnBuch\color{poporange}+}
  \end{tcolorbox}%
}

% ============================================================
% Page de garde
% #1 titre · #2 matière · #3 niveau · #4 module · #5 n° leçon · #6 durée ·
% #7 description / objectifs (texte ou itemize)
% ============================================================
\newcommand{\pagegarde}[7]{%
  \thispagestyle{empty}
  \newgeometry{a4paper,margin=1.7cm,top=1.6cm,bottom=1.4cm}%
  \begin{tikzpicture}[remember picture,overlay]
    \fill[poporange!18] ([xshift=-3.2cm,yshift=-3.6cm]current page.north east) circle (4.6cm);
    \fill[poppurple!14] ([xshift=2.4cm,yshift=7.6cm]current page.south west) circle (3.8cm);
  \end{tikzpicture}%
  {\poplogo\fontsize{30}{30}\selectfont\color{popink}OnBuch\color{poporange}+}\hfill
  \raisebox{0.6ex}{\poppill{Cours signé \textbullet\ Premium}{popink}{popcream}}\par
  \vspace{1pt}\popsquiggle{poppurple}\par
  \vspace{8pt}
  \begin{tcolorbox}[enhanced,colback=poporange,colframe=popink,boxrule=2.8pt,arc=13pt,
      drop shadow=popink,left=18pt,right=18pt,top=12pt,bottom=13pt,after skip=10pt]
    \poppill{#2 \textbullet\ #3}{popink}{popcream}\hspace{5pt}\poppill{#4}{white}{popink}\par\vspace{8pt}
    {\popdisplay\fontsize{14}{14}\selectfont\color{popink}LEÇON #5}\par\vspace{4pt}
    {\raggedright\hyphenpenalty=10000\exhyphenpenalty=10000\popdisplay\fontsize{25}{27}\selectfont\color{popcream}\MakeUppercase{#1}\par}
    \vspace{8pt}
    {\popxbold\fontsize{9.5}{9.5}\selectfont\color{popink}\MakeUppercase{Durée conseillée : #6}}
  \end{tcolorbox}
  \begin{tcolorbox}[enhanced,colback=white,colframe=popink,boxrule=2.2pt,arc=10pt,
      drop shadow=popink,left=16pt,right=16pt,top=10pt,bottom=10pt,after skip=8pt]
    \poppill{Dans cette leçon}{poppurple}{white}\par\vspace{5pt}
    \popsemi\fontsize{9.3}{12.6}\selectfont\color{popink}#7
  \end{tcolorbox}
  \noindent\begin{minipage}[t]{0.487\linewidth}
    \begin{tcolorbox}[enhanced,colback=popgreen,colframe=popink,boxrule=2.2pt,arc=10pt,
        drop shadow=popink,left=11pt,right=11pt,top=10pt,bottom=10pt,height=3.1cm,valign=center]
      \tcbox[colback=white,colframe=popink,boxrule=1.3pt,arc=5pt,boxsep=0pt,nobeforeafter,
        left=3pt,right=3pt,top=3pt,bottom=3pt,valign=center]{\qrcode[height=1.9cm]{https://play.google.com/store/apps/details?id=com.luvvix.onbuch&hl=fr}}%
      \hspace{8pt}\begin{minipage}[c]{0.5\linewidth}
        {\popdisplay\fontsize{11}{12.5}\selectfont\color{white}\MakeUppercase{Télécharge OnBuch}}\par\vspace{4pt}
        {\raggedright\popsemi\fontsize{8.4}{11}\selectfont\color{white}Gratuit \textbullet\ Google Play\\Tuteur IA, annales, concours, résultats\par}
      \end{minipage}
    \end{tcolorbox}
  \end{minipage}\hfill
  \begin{minipage}[t]{0.487\linewidth}
    \begin{tcolorbox}[enhanced,colback=poppurple,colframe=popink,boxrule=2.2pt,arc=10pt,
        drop shadow=popink,left=11pt,right=11pt,top=10pt,bottom=10pt,height=3.1cm,valign=center]
      \tikz[baseline=-0.7ex]\node[circle,fill=white,draw=popink,line width=1.3pt,minimum size=1.6cm,inner sep=0]{\popdisplay\fontsize{24}{24}\selectfont\color{poppurple}+};%
      \hspace{8pt}\begin{minipage}[c]{0.6\linewidth}
        {\popdisplay\fontsize{11}{12.5}\selectfont\color{white}\MakeUppercase{Passe à OnBuch+}}\par\vspace{4pt}
        {\raggedright\popsemi\fontsize{8.4}{11}\selectfont\color{white}Tous les cours signés, Tuteur IA illimité, fiches PDF — onglet OnBuch+ de l'app\par}
      \end{minipage}
    \end{tcolorbox}
  \end{minipage}
  \vspace{8pt}
  \popcontenu
  \vfill
  \signatureonbuch
  \restoregeometry
  \newpage
}

% Rangée « Ce cours contient » (page de garde)
\newcommand{\poptile}[3]{% #1 couleur #2 gros texte #3 légende
  \begin{tcolorbox}[enhanced,colback=#1,colframe=popink,boxrule=2pt,arc=9pt,shadow={2.5pt}{-2.5pt}{0pt}{popink},
      left=6pt,right=6pt,top=9pt,bottom=9pt,halign=center,valign=center,height=1.75cm,width=\linewidth,nobeforeafter]
    {\popdisplay\fontsize{12.5}{13}\selectfont\color{white}\MakeUppercase{#2}}\par\vspace{3pt}
    {\centering\popsemi\fontsize{7.8}{9.5}\selectfont\color{white}#3\par}
  \end{tcolorbox}}
\newcommand{\popcontenu}{%
  \par\noindent{\popxboldls\fontsize{7.5}{7.5}\selectfont\color{popmuted}CE COURS SIGNÉ CONTIENT}\par\vspace{7pt}
  \noindent\begin{minipage}[t]{0.235\linewidth}\poptile{popblue}{Cours}{complet, illustré, pas à pas}\end{minipage}\hfill
  \begin{minipage}[t]{0.235\linewidth}\poptile{poporange}{Méthodes}{et exercices résolus}\end{minipage}\hfill
  \begin{minipage}[t]{0.235\linewidth}\poptile{poppink}{Exercices}{type examen + corrigés}\end{minipage}\hfill
  \begin{minipage}[t]{0.235\linewidth}\poptile{popgreen}{Fiche bilan}{l'essentiel à retenir}\end{minipage}\par}

% Sommaire
\newtcolorbox{sommaire}{enhanced,colback=white,colframe=popink,boxrule=2.2pt,arc=10pt,
  drop shadow=popink,left=18pt,right=18pt,top=13pt,bottom=13pt,before skip=6pt,after skip=20pt,
  before upper={\poppill{Sommaire}{poppurple}{white}\par\vspace{9pt}\popsemi\fontsize{10.6}{14.5}\selectfont\color{popink}}}

% Signature de fin de cours — poussée tout en bas de la dernière page
\newcommand{\findecours}{%
  \par\vfill\nopagebreak
  \begin{center}\popsquiggle{poporange}\end{center}\vspace{4pt}
  \signatureonbuch
}
\AtBeginDocument{\def\llbracket{\mathopen{[\mkern-3.5mu[}}\def\rrbracket{\mathclose{]\mkern-3.5mu]}}} % intervalles d'entiers (glyphe ⟦ absent de la police)
% Glyphes absents de Fira Math : redéfinis à partir de glyphes disponibles
\AtBeginDocument{%
  \def\setminus{\mathbin{\backslash}}\let\smallsetminus\setminus
  \def\vdots{\mathord{\vbox{\baselineskip3.2pt\lineskiplimit0pt\kern4pt\hbox{.}\hbox{.}\hbox{.}}}}%
  \def\ddots{\mathinner{\mkern1mu\raise6.5pt\hbox{.}\mkern2mu\raise3.6pt\hbox{.}\mkern2mu\raise0.7pt\hbox{.}\mkern1mu}}%
  \def\triangle{\mathord{\scalebox{0.9}{$\symup{\Delta}$}}}%
  \def\nabla{\mathord{\raisebox{\depth}{\scalebox{1}[-1]{$\symup{\Delta}$}}}}%
  \def\checkmark{\popcheck{popdark}}%
  \def\star{\mathbin{\ast}}\def\bigstar{\mathord{\ast}}%
  \def\diamond{\mathbin{\scalebox{0.6}{$\Diamond$}}}%
  \def\bigcup{\mathop{\mathchoice{\raisebox{-0.3em}{\scalebox{1.7}{$\cup$}}}{\scalebox{1.25}{$\cup$}}{\cup}{\cup}}\displaylimits}%
  \def\bigcap{\mathop{\mathchoice{\raisebox{-0.3em}{\scalebox{1.7}{$\cap$}}}{\scalebox{1.25}{$\cap$}}{\cap}{\cap}}\displaylimits}%
  \def\lesseqqgtr{\mathrel{\lessgtr}}\def\bigoplus{\mathop{\oplus}\displaylimits}%
}
\providecommand{\ssi}{si et seulement si} % abréviation fréquente
\AtBeginDocument{\providecommand{\wideparen}{\overparen}} % arc de cercle

%% FIGFIX-BEGIN v5 — correctifs globaux des figures (docs/onbuch-generation/tools/figfix)
\usepackage{adjustbox}\usepackage{etoolbox}\tcbuselibrary{raster}
% toute figure TikZ/pgfplots plus large que la ligne est réduite à la largeur disponible
\BeforeBeginEnvironment{tikzpicture}{\begin{adjustbox}{max width=\linewidth}}
\AfterEndEnvironment{tikzpicture}{\end{adjustbox}}
% légendes pgfplots lisibles par-dessus les courbes
\pgfplotsset{every axis/.append style={legend style={fill=white,fill opacity=0.92,text opacity=1,draw=black!15,inner sep=3pt}}}
% étiquettes d'arêtes (syntaxe edge["..."]) avec fond blanc
\tikzset{every edge quotes/.append style={fill=white,inner sep=1.2pt}}
%% FIGFIX-END
\begin{document}

\pagegarde{Leçon 9 — Vocabulaire et compréhension de texte : processus du recyclage}{Chinois}{Tle A}{Module 2 : Environnement et santé}{ZH09}{2 h}{Cette leçon de vocabulaire et compréhension de texte fait découvrir aux élèves de Terminale A le lexique chinois du recyclage à travers un texte support original décrivant le processus complet de traitement des déchets. Elle développe la lecture, la prononciation, l'écriture des caractères clés et la réponse à des questions de compréhension.
\vspace{4pt}
\begin{multicols}{2}\small
\begin{itemize}
\item À la fin de cette leçon, je sais lire à voix haute un court texte chinois sur le recyclage avec une prononciation correcte (pinyin et tons).
\item À la fin de cette leçon, je sais reconnaître et traduire au moins 15 mots de vocabulaire liés au recyclage.
\item À la fin de cette leçon, je sais identifier les radicaux (clés) des caractères du thème et expliquer leur lien avec le sens.
\item À la fin de cette leçon, je sais écrire correctement 5 caractères du thème en respectant l'ordre des traits.
\item À la fin de cette leçon, je sais répondre par écrit à des questions de compréhension sur le texte.
\item À la fin de cette leçon, je sais employer les collocations essentielles du thème (回收利用, 保护环境, 垃圾分类) dans des phrases complètes.
\end{itemize}\end{multicols}}

\begin{sommaire}
\begin{multicols}{2}
\begin{enumerate}
\item Mise en route et découverte du texte
\item Glossaire du thème
\item Clés (radicaux) et ordre des traits
\item Compréhension détaillée du texte
\item Collocations et réemploi du lexique
\item Production guidée et évaluation
\item Activité d'intégration
\item Méthodes et exercices résolus
\item Exercices
\item Corrigés des exercices
\item Fiche bilan
\end{enumerate}
\end{multicols}
\end{sommaire}

\begin{objectifs}
\begin{itemize}
\item À la fin de cette leçon, je sais lire à voix haute un court texte chinois sur le recyclage avec une prononciation correcte (pinyin et tons).
\item À la fin de cette leçon, je sais reconnaître et traduire au moins 15 mots de vocabulaire liés au recyclage.
\item À la fin de cette leçon, je sais identifier les radicaux (clés) des caractères du thème et expliquer leur lien avec le sens.
\item À la fin de cette leçon, je sais écrire correctement 5 caractères du thème en respectant l'ordre des traits.
\item À la fin de cette leçon, je sais répondre par écrit à des questions de compréhension sur le texte.
\item À la fin de cette leçon, je sais employer les collocations essentielles du thème (回收利用, 保护环境, 垃圾分类) dans des phrases complètes.
\item À la fin de cette leçon, je sais décrire oralement et par écrit les étapes du processus de recyclage en chinois.
\end{itemize}
\end{objectifs}

\begin{prerequis}
\begin{itemize}
\item Vocabulaire de l'environnement vu en Première (环境, 保护, 污染)
\item Structures de base : phrase sujet-verbe-objet, particule 了, adverbes de fréquence
\item Lecture du pinyin et maîtrise des quatre tons
\item Notion de radicaux (clés) et ordre général des traits (de haut en bas, de gauche à droite)
\item Connecteurs simples de séquence : 先...然后...最后...
\end{itemize}
\end{prerequis}

\newpage

\coursec{Mise en route et découverte du texte}

\courssub{Situation d'accroche : la bouteille plastique, de Yaoundé à Shanghai}

Imaginez la scène : il est 17 heures à Yaoundé, la pluie vient de s'arrêter, et dans le quartier de Mokolo, des centaines de bouteilles plastiques flottent dans les caniveaux. À Douala, au marché Mboppi, c'est le même spectacle après chaque averse. Que deviennent toutes ces bouteilles ? La plupart finissent dans la nature, dans les rivières, ou brûlées à ciel ouvert.

Pourtant, de l'autre côté du monde, en Chine, le tri et le recyclage sont devenus une priorité nationale : Shanghai impose depuis 2019 un tri strict des déchets, avec des poubelles de couleurs différentes dans chaque quartier et des amendes pour les contrevenants. Une bouteille plastique jetée à Shanghai commence alors un véritable voyage : collecte, tri, nettoyage, transformation, réutilisation.

\textbf{Problématique de la leçon :} comment les Chinois décrivent-ils le parcours d'une bouteille plastique, de la poubelle à sa nouvelle vie ? Quels mots utilisent-ils pour chaque étape ? C'est ce que nous allons découvrir en lisant un court texte chinois sur le processus du recyclage.

\begin{experience}[Brainstorming : que devient une bouteille plastique jetée à Yaoundé ?]
En français, par groupes de quatre, répondez en cinq minutes aux questions suivantes :
\begin{enumerate}
\item Où jetez-vous vos bouteilles plastiques ? Existe-t-il des poubelles de tri dans votre quartier ou votre lycée ?
\item Qui ramasse les ordures dans votre ville ? Que deviennent-elles ensuite ?
\item Connaissez-vous des artisans camerounais qui transforment de vieux objets (pneus, bouteilles, boîtes de conserve) en objets neufs ?
\end{enumerate}
Le professeur note au tableau les mots-clés français proposés par la classe : \emph{jeter, ramasser, trier, laver, fondre, fabriquer, réutiliser}. Gardez-les en tête : nous allons chercher leurs équivalents chinois dans le texte.
\end{experience}

\courssub{Présentation du processus : les cinq étapes du recyclage}

Avant de lire le texte, observons le parcours d'une bouteille plastique. Le processus du recyclage se déroule en cinq étapes, que les Chinois nomment ainsi :

\begin{center}
\begin{tabularx}{\linewidth}{|X|X|X|X|}
\hline
\rowcolor{popblueL} Caractères & Pinyin & Nature & Traduction \\
\hline
收集 & shōují & verbe & collecter, ramasser \\
\hline
分类 & fēnlèi & verbe & trier, classer \\
\hline
清洗 & qīngxǐ & verbe & nettoyer, laver \\
\hline
加工 & jiāgōng & verbe & transformer, usiner \\
\hline
再利用 & zài lìyòng & verbe & réutiliser, valoriser \\
\hline
\end{tabularx}
\end{center}

Observez bien ces mots : plusieurs contiennent des indices de sens. Le caractère \zh{清} (qīng, « clair, propre ») porte la clé de l'eau \zh{氵} : logique pour un mot qui signifie « laver » ! Le caractère \zh{类} (lèi, « catégorie ») apparaît dans \zh{分类}, « séparer en catégories », c'est-à-dire « trier ».

\begin{popfigure}
\begin{tikzpicture}[
  etape/.style={rectangle, rounded corners=6pt, draw=popblue, fill=popblueL, thick, align=center, minimum width=2.6cm, minimum height=1.1cm, font=\small},
  fleche/.style={-{Stealth[length=3mm]}, very thick, poporange}
]
\node[etape, align=center] (e1) at (0,0){1. 收集\\ shōují\\ collecte};
\node[etape, align=center] (e2) at (4.2,0){2. 分类\\ fēnlèi\\ tri};
\node[etape, align=center] (e3) at (8.4,0){3. 清洗\\ qīngxǐ\\ nettoyage};
\node[etape, align=center] (e4) at (8.4,-2.6){4. 加工\\ jiāgōng\\ transformation};
\node[etape, align=center] (e5) at (4.2,-2.6){5. 再利用\\ zài lìyòng\\ réutilisation};
\draw[fleche] (e1) -- (e2);
\draw[fleche] (e2) -- (e3);
\draw[fleche] (e3) -- (e4);
\draw[fleche] (e4) -- (e5);
\draw[fleche, dashed] (e5) to[bend right=35] node[midway, below, font=\small\itshape, text=popdark] {le cycle recommence !} (e1);
\end{tikzpicture}
\legende{Le cycle du recyclage en cinq étapes : de la collecte à la réutilisation, la boucle est sans fin.}
\end{popfigure}

\begin{savaistu}[Le tri des déchets à Shanghai]
Depuis juillet 2019, les habitants de Shanghai doivent trier leurs déchets en quatre catégories : \zh{可回收物} (kě huíshōu wù, déchets recyclables), \zh{有害垃圾} (yǒuhài lājī, déchets dangereux), \zh{湿垃圾} (shī lājī, déchets humides, c'est-à-dire organiques) et \zh{干垃圾} (gān lājī, déchets secs). Des volontaires aident les habitants devant les poubelles ! Au Cameroun, des entreprises et des associations commencent aussi à collecter le plastique à Douala et à Yaoundé pour le transformer en pavés ou en objets neufs : le vocabulaire de cette leçon vous servira donc aussi pour parler de votre propre pays.
\end{savaistu}

\courssub{Première écoute globale (sans support écrit)}

Fermez vos cahiers. Le professeur lit le texte deux fois, à vitesse normale. Votre tâche : ne retenez que \textbf{l'idée générale}. Posez-vous trois questions : De quoi parle le texte ? Qui agit ? Dans quel ordre se déroulent les actions ?

Pendant l'écoute, essayez de reconnaître les mots que vous connaissez déjà : \zh{环境} (huánjìng, environnement) et \zh{保护} (bǎohù, protéger), vus dans la leçon précédente, ainsi que les mots de temps \zh{先} (xiān, d'abord), \zh{然后} (ránhòu, ensuite) et \zh{最后} (zuìhòu, enfin), qui annoncent une suite d'actions.

\begin{methode}[Lire un texte chinois en trois passes]
\begin{enumerate}
\item \textbf{Survol} (première écoute ou lecture rapide) : ne cherchez pas à tout comprendre. Identifiez le thème général et les mots déjà connus. Ignorez les mots inconnus.
\item \textbf{Lecture fine} : relisez phrase par phrase avec le glossaire. Soulignez les groupes de sens (sujet + verbe + complément) et les marqueurs de temps.
\item \textbf{Relecture de vérification} : relisez le texte entier pour vérifier vos réponses aux questions de compréhension et goûter la fluidité du texte.
\end{enumerate}
\end{methode}

\courssub{Lecture silencieuse, puis lecture chorale}

Ouvrez maintenant les yeux sur le texte. Lisez-le d'abord silencieusement en suivant le pinyin, puis nous le lirons ensemble à voix haute, en respectant les tons et les pauses.

\begin{textebox}[回收很重要 — Le recyclage est très important]
回收很重要。每天，人们把垃圾放进不同的垃圾桶。工人先把垃圾收集起来，然后送到回收站。在那里，塑料、纸和玻璃被分类、清洗，再送到工厂加工。最后，旧东西变成新产品。回收利用可以保护环境，节约资源。我们应该养成垃圾分类的好习惯。\\
\emph{Huíshōu hěn zhòngyào. Měitiān, rénmen bǎ lājī fàng jìn bùtóng de lājītǒng. Gōngrén xiān bǎ lājī shōují qǐlái, ránhòu sòng dào huíshōuzhàn. Zài nàlǐ, sùliào, zhǐ hé bōli bèi fēnlèi, qīngxǐ, zài sòng dào gōngchǎng jiāgōng. Zuìhòu, jiù dōngxi biànchéng xīn chǎnpǐn. Huíshōu lìyòng kěyǐ bǎohù huánjìng, jiéyuē zīyuán. Wǒmen yīnggāi yǎngchéng lājī fēnlèi de hǎo xíguàn.}\\
« Le recyclage est très important. Chaque jour, les gens mettent les déchets dans différentes poubelles. Les ouvriers collectent d'abord les déchets, puis les envoient au centre de recyclage. Là, le plastique, le papier et le verre sont triés, nettoyés, puis envoyés à l'usine pour être transformés. Enfin, les vieux objets deviennent de nouveaux produits. Recycler permet de protéger l'environnement et d'économiser les ressources. Nous devons prendre la bonne habitude de trier les déchets. »
\end{textebox}

\textbf{Glossaire rapide du texte} (les mots nouveaux seront étudiés en détail dans la section suivante) :

\begin{center}
\begin{tabularx}{\linewidth}{|X|X|X|X|}
\hline
\rowcolor{popblueL} Caractères & Pinyin & Nature & Traduction \\
\hline
回收 & huíshōu & verbe / nom & recycler ; recyclage \\
\hline
垃圾 & lājī & nom & déchets, ordures \\
\hline
垃圾桶 & lājītǒng & nom & poubelle \\
\hline
工人 & gōngrén & nom & ouvrier \\
\hline
回收站 & huíshōuzhàn & nom & centre de recyclage \\
\hline
塑料 & sùliào & nom & plastique \\
\hline
玻璃 & bōli & nom & verre \\
\hline
工厂 & gōngchǎng & nom & usine \\
\hline
产品 & chǎnpǐn & nom & produit \\
\hline
节约 & jiéyuē & verbe & économiser \\
\hline
资源 & zīyuán & nom & ressources \\
\hline
养成 & yǎngchéng & verbe & prendre (une habitude) \\
\hline
习惯 & xíguàn & nom & habitude \\
\hline
\end{tabularx}
\end{center}

Pour la lecture chorale, respectez les groupes de souffle : \zh{工人 / 先把垃圾 / 收集起来，/ 然后 / 送到回收站。} Marquez bien les tons : \zh{垃圾} (lājī) est en premier ton sur les deux syllabes, haut et plat ; \zh{回收} (huíshōu) commence par un deuxième ton montant.

\courssub{Repérage : mots connus et mots nouveaux}

\begin{exoresolu}[Repérer les mots connus et deviner les mots nouveaux]
Énoncé : dans le texte, (a) retrouvez deux mots déjà étudiés dans les leçons précédentes ; (b) devinez le sens de \zh{回收站} et de \zh{垃圾桶} sans dictionnaire, en vous aidant des mots que vous venez de repérer.
\tcblower
Solution détaillée :
\begin{enumerate}
\item Mots déjà connus : \zh{环境} (huánjìng, environnement) et \zh{保护} (bǎohù, protéger), étudiés dans la leçon sur la protection de l'environnement. On reconnaît aussi \zh{每天} (měitiān, chaque jour) et \zh{可以} (kěyǐ, pouvoir).
\item \zh{回收站} : on connaît \zh{回收} (recyclage) et \zh{站} (zhàn, station, arrêt, comme dans \zh{车站}, gare/station). Un « lieu-station du recyclage » est donc un \emph{centre de recyclage}.
\item \zh{垃圾桶} : \zh{垃圾} (déchets) + \zh{桶} (tǒng, seau, tonneau). Un « seau à déchets » est une \emph{poubelle}. Le contexte le confirme : \zh{人们把垃圾放进不同的垃圾桶} — « les gens mettent les déchets dans différentes poubelles ».
\end{enumerate}
Stratégie : un mot composé chinois se devine souvent par la somme de ses parties. C'est une grande force de la langue chinoise !
\end{exoresolu}

\begin{propriete}[Deviner un mot composé par ses éléments]
En chinois, la plupart des mots sont composés de deux caractères dont le sens se combine. Pour deviner un mot inconnu :
\begin{enumerate}
\item identifiez les caractères que vous connaissez déjà ;
\item combinez leurs sens ;
\item vérifiez votre hypothèse avec le contexte de la phrase.
\end{enumerate}
Exemples du texte : \zh{工人} = \zh{工} (travail) + \zh{人} (personne) = « ouvrier » ; \zh{工厂} = \zh{工} (travail) + \zh{厂} (chǎng, fabrique) = « usine » ; \zh{产品} = \zh{产} (chǎn, produire) + \zh{品} (pǐn, article) = « produit ».
\end{propriete}

\courssub{Hypothèses de sens : le contexte et les radicaux}

Deux outils vous permettent de deviner le sens d'un mot inconnu : le \textbf{contexte} de la phrase et les \textbf{clés (radicaux)} des caractères.

\textbf{1. Le contexte.} Dans la phrase \zh{塑料、纸和玻璃被分类、清洗} (sùliào, zhǐ hé bōli bèi fēnlèi, qīngxǐ), même si vous ne connaissez pas \zh{塑料} et \zh{玻璃}, vous voyez qu'ils sont énumérés avec \zh{纸} (zhǐ, papier) et qu'on les « trie » et les « lave » : ce sont donc des matériaux recyclables. Le contexte suffit à comprendre la phrase.

\textbf{2. Les radicaux.} Les clés donnent une indication sur le domaine de sens :

\begin{center}
\begin{tabularx}{\linewidth}{|X|X|X|X|}
\hline
\rowcolor{popblueL} Caractère & Clé (radical) & Domaine de sens & Traduction \\
\hline
清 (qīngxǐ) & 氵 eau & lié à l'eau, au lavage & nettoyer \\
\hline
玻 (bōli) & 王 (玉) jade/pierre & matière minérale & verre \\
\hline
桶 (lājītǒng) & 木 bois & objet (autrefois en bois) & seau, poubelle \\
\hline
环 (huánjìng) & 王 (玉) jade & objet précieux, cercle & environnement \\
\hline
\end{tabularx}
\end{center}

\begin{attention}[Ne pas traduire mot à mot]
Le piège classique du francophone : traduire caractère par caractère, dans l'ordre. La phrase \zh{工人先把垃圾收集起来} ne se traduit pas « ouvriers d'abord prendre déchets collecter se lever » ! Repérez d'abord les \textbf{groupes de sens} : \zh{工人} (sujet : les ouvriers) + \zh{先把垃圾收集起来} (bloc d'action : collectent d'abord les déchets). La construction avec \zh{把} (bǎ) place l'objet \zh{垃圾} avant le verbe : « prendre les déchets et les collecter ». Traduisez par blocs, jamais mot à mot.
\end{attention}

\begin{aretenir}[Ce qu'il faut retenir de cette première étape]
\begin{itemize}
\item Le recyclage suit cinq étapes : \zh{收集} → \zh{分类} → \zh{清洗} → \zh{加工} → \zh{再利用}.
\item On lit un texte chinois en trois passes : survol, lecture fine, relecture de vérification.
\item Les mots composés se devinent par la somme de leurs caractères (\zh{垃圾} + \zh{桶} = poubelle).
\item Les radicaux (氵 eau, 木 bois, 王 jade) orientent vers le domaine de sens.
\item On traduit par groupes de sens, jamais mot à mot.
\end{itemize}
\end{aretenir}

\begin{exercice}[À vous de jouer]
\begin{enumerate}
\item Sans regarder le glossaire, devinez le sens de : \zh{新产品} (xīn chǎnpǐn) et \zh{旧东西} (jiù dōngxi).
\item Classez ces mots du texte dans les cinq étapes du cycle : \zh{工人}, \zh{垃圾桶}, \zh{工厂}, \zh{回收站}, \zh{新产品}. À quelle étape chacun est-il associé ?
\item Prononcez à voix haute : \zh{回收利用可以保护环境} (huíshōu lìyòng kěyǐ bǎohù huánjìng), en veillant aux tons.
\end{enumerate}
\end{exercice}

\begin{corrige}[Exercice « À vous de jouer »]
\begin{enumerate}
\item \zh{新产品} = \zh{新} (nouveau) + \zh{产品} (produit) = « nouveau produit » ; \zh{旧东西} = \zh{旧} (jiù, vieux, usagé) + \zh{东西} (dōngxi, chose) = « vieil objet ».
\item \zh{垃圾桶} → étape 1 (avant la collecte, le jet) ; \zh{工人} → étape 1 (collecte) ; \zh{回收站} → étape 2 (tri) ; \zh{工厂} → étape 4 (transformation) ; \zh{新产品} → étape 5 (réutilisation).
\item Vérifiez : huíshōu (2\textsuperscript{e} + 1\textsuperscript{er} ton), lìyòng (4\textsuperscript{e} + 4\textsuperscript{e} ton), kěyǐ (3\textsuperscript{e} + 3\textsuperscript{e} ton, le premier devient 2\textsuperscript{e} ton), bǎohù (3\textsuperscript{e} + 4\textsuperscript{e} ton), huánjìng (2\textsuperscript{e} + 4\textsuperscript{e} ton).
\end{enumerate}
\end{corrige}

Nous avons maintenant une vision globale du texte et de son vocabulaire. Dans la section suivante, nous étudierons en détail le glossaire du thème : prononciation, nature des mots et emplois en contexte.

\coursec{Glossaire du thème}

Dans la section précédente, nous avons découvert le texte support sur le processus du recyclage et repéré globalement son organisation : collecte, tri, transport, transformation, nouveaux produits. Pour bien comprendre ce texte et pouvoir en parler, il faut maintenant maîtriser son vocabulaire. C'est l'objet de cette section : apprendre les mots-clés, bien les prononcer, comprendre leur formation et savoir les employer dans des phrases.

\courssub{Le tableau de vocabulaire du recyclage}

\begin{definition}[Comment lire ce tableau]
Chaque mot est présenté en quatre colonnes : les \cle{caractères} simplifiés, le \cle{pinyin} avec les marques de tons, la \cle{nature} grammaticale (n. = nom, v. = verbe, adj. = adjectif) et la \cle{traduction} française. Apprends toujours un mot avec son ton : en chinois, un mot sans ton, c'est un mot à moitié appris.
\end{definition}

\begin{center}
\begin{tabularx}{\linewidth}{|X|X|X|X|}
\hline
\rowcolor{popblueL} Caractères & Pinyin & Nature & Traduction \\
\hline
回收 & huíshōu & v. & recycler, récupérer \\
\hline
垃圾 & lājī & n. & déchets, ordures \\
\hline
分类 & fēnlèi & v. & trier, classer \\
\hline
垃圾桶 & lājītǒng & n. & poubelle \\
\hline
塑料 & sùliào & n. & plastique \\
\hline
玻璃 & bōli & n. & verre (matériau) \\
\hline
纸 & zhǐ & n. & papier \\
\hline
工厂 & gōngchǎng & n. & usine \\
\hline
加工 & jiāgōng & v. & transformer, traiter \\
\hline
产品 & chǎnpǐn & n. & produit \\
\hline
资源 & zīyuán & n. & ressource \\
\hline
节约 & jiéyuē & v. & économiser \\
\hline
保护 & bǎohù & v. & protéger \\
\hline
环境 & huánjìng & n. & environnement \\
\hline
习惯 & xíguàn & n. & habitude \\
\hline
\end{tabularx}
\end{center}

\begin{aretenir}[Mots composés : devine le sens par les parties]
Beaucoup de mots de cette liste sont des \cle{mots composés} : leur sens se devine à partir de celui des caractères.
\begin{itemize}
\item \zh{垃圾桶} lājītǒng : 垃圾 (déchets) + 桶 (seau) = « seau à déchets » = poubelle.
\item \zh{分类} fēnlèi : 分 (diviser) + 类 (catégorie) = « diviser en catégories » = trier.
\item \zh{加工} jiāgōng : 加 (ajouter) + 工 (travail) = « ajouter du travail » = transformer.
\item \zh{产品} chǎnpǐn : 产 (produire) + 品 (article) = produit.
\item \zh{工厂} gōngchǎng : 工 (travail) + 厂 (atelier) = usine.
\end{itemize}
Réflexe de lecteur : face à un mot inconnu, décompose-le mentalement caractère par caractère.
\end{aretenir}

\courssub{Le mot central : 回收 huíshōu}

\begin{definition}[回收 huíshōu]
\zh{回收} signifie « récupérer ce qui a été utilisé pour le réintroduire dans un cycle » : recycler. Le caractère \zh{回} (huí) exprime le \emph{retour}, le retour en arrière ; \zh{收} (shōu) signifie \emph{recevoir, ramasser}. Littéralement : « recevoir en retour ». L'image est celle d'un cercle : l'objet part, est utilisé, puis revient pour recommencer un cycle.
\end{definition}

\begin{exemplebox}[回收 en contexte]
\zh{我们应该回收塑料瓶。}\\
\emph{Wǒmen yīnggāi huíshōu sùliào píng.}\\
« Nous devrions recycler les bouteilles en plastique. »

\zh{这个城市每天回收很多纸和玻璃。}\\
\emph{Zhège chéngshì měitiān huíshōu hěn duō zhǐ hé bōli.}\\
« Cette ville recycle chaque jour beaucoup de papier et de verre. »

\zh{垃圾分类可以保护环境。}\\
\emph{Lājī fēnlèi kěyǐ bǎohù huánjìng.}\\
« Le tri des déchets peut protéger l'environnement. »
\end{exemplebox}

\courssub{回收 ou 利用 ? Une distinction essentielle}

Les francophones traduisent souvent « recycler » et « réutiliser » par le même mot. Le chinois distingue deux verbes :

\begin{propriete}[回收 huíshōu ≠ 利用 lìyòng]
\begin{itemize}
\item \zh{回收} huíshōu : \emph{récupérer, recycler} — on reprend un objet usagé pour le renvoyer dans le circuit (collecte, tri, usine).
\item \zh{利用} lìyòng : \emph{utiliser, exploiter} — on se sert de quelque chose (une ressource, un objet, du temps).
\item La combinaison \zh{回收利用} huíshōu lìyòng signifie « recycler et réutiliser » : c'est l'expression complète pour parler du recyclage comme filière.
\end{itemize}
\end{propriete}

\begin{center}
\begin{tabularx}{\linewidth}{|X|X|X|}
\hline
\rowcolor{popblueL} Verbe & Exemple (caractères + pinyin) & Traduction \\
\hline
回收 (récupérer) & \zh{回收旧报纸} huíshōu jiù bàozhǐ & recycler les vieux journaux \\
\hline
利用 (utiliser) & \zh{利用太阳能} lìyòng tàiyángnéng & utiliser l'énergie solaire \\
\hline
回收利用 (recycler et réutiliser) & \zh{回收利用塑料} huíshōu lìyòng sùliào & recycler et réutiliser le plastique \\
\hline
\end{tabularx}
\end{center}

\begin{exemplebox}[Compare ces deux phrases]
\zh{工厂回收玻璃，加工以后做成新产品。}\\
\emph{Gōngchǎng huíshōu bōli, jiāgōng yǐhòu zuò chéng xīn chǎnpǐn.}\\
« L'usine récupère le verre et, après transformation, en fait de nouveaux produits. »

\zh{我们应该利用每一分钟学习汉语。}\\
\emph{Wǒmen yīnggāi lìyòng měi yī fēnzhōng xuéxí Hànyǔ.}\\
« Nous devrions utiliser chaque minute pour apprendre le chinois. »
\end{exemplebox}

\begin{attention}[Ne dis pas 回收 le temps !]
On ne peut pas dire \zh{回收时间} : le temps ne « revient » pas dans un cycle, on ne le récupère pas. Pour « profiter du temps », utilise \zh{利用时间} (lìyòng shíjiān). De même, « économiser les ressources » se dit \zh{节约资源} (jiéyuē zīyuán), jamais \zh{保护资源} au sens d'économiser : \zh{保护} signifie \emph{protéger} (contre un danger), pas \emph{épargner}.
\end{attention}

\courssub{Travail de prononciation : les tons difficiles}

\begin{propriete}[Trois pièges de tons dans ce glossaire]
\begin{enumerate}
\item \zh{垃圾} lājī : \cle{1\ier{} ton + 1\ier{} ton}. Deux sons hauts et plats, sans descendre. Ne laisse pas le deuxième syllabe retomber comme en français (« la-ji » chantant).
\item \zh{塑料} sùliào : \cle{4\ieme{} ton + 4\ieme{} ton}. Deux tons descendants, secs et décidés, comme deux ordres brefs : sù ! liào !
\item \zh{玻璃} bōli : \cle{1\ier{} ton + ton léger}. La deuxième syllabe est atone (ton neutre) : courte et légère, sans mélodie propre. C'est fréquent dans les mots de deux syllabes nommant des objets courants (compare \zh{妈妈} māma).
\end{enumerate}
\end{propriete}

\begin{attention}[垃圾 se dit lājī, pas lèsè]
En Chine continentale, \zh{垃圾} se prononce \cle{lājī}, avec deux premiers tons. La prononciation \emph{lèsè} existe à Taiwan, mais elle n'est pas la norme du pǔtōnghuà (mandarin standard) enseigné au programme MINESEC. Retiens : \zh{垃圾} = lājī, deux 1\iers{} tons bien hauts et plats.
\end{attention}

\begin{experience}[Entraînement oral en chaîne]
Par groupes de deux : l'élève A prononce un mot du tableau, l'élève B doit donner la traduction et répéter le mot avec les bons tons, puis proposer une phrase courte. Inversez les rôles. Le professeur veille surtout aux trois pièges : lājī (1-1), sùliào (4-4), bōli (1-léger).
\end{experience}

\courssub{Carte mentale du lexique}

Pour mémoriser durablement, organisons le vocabulaire autour du mot central \zh{回收} en quatre familles : les \cle{matériaux}, les \cle{lieux}, les \cle{actions} et les \cle{buts}.

\begin{popfigure}
\begin{center}\tcbox[colback=popblue,colframe=popblue,coltext=white,arc=9pt,boxsep=4pt,left=6pt,right=6pt]{\bfseries\small 回收 huíshōu recycler}\end{center}
\vspace{-2pt}
\begin{tcbraster}[raster columns=2,raster equal height=rows,raster column skip=6pt,raster row skip=6pt,enhanced,blankest]
\begin{tcolorbox}[enhanced,colback=white,colframe=popgreen,boxrule=0.9pt,arc=3pt,title={matériaux 材料},fonttitle=\bfseries\scriptsize,coltitle=white,colbacktitle=popgreen,left=4pt,right=4pt,top=2pt,bottom=2pt,boxsep=1pt,toptitle=1.5pt,bottomtitle=1.5pt,halign=left]
\begin{itemize}[nosep,leftmargin=*,itemsep=1pt,font=\scriptsize]
\item 塑料 sùliào
\item 纸 zhǐ
\item 玻璃 bōli
\end{itemize}
\end{tcolorbox}
\begin{tcolorbox}[enhanced,colback=white,colframe=poporange,boxrule=0.9pt,arc=3pt,title={lieux 地方},fonttitle=\bfseries\scriptsize,coltitle=white,colbacktitle=poporange,left=4pt,right=4pt,top=2pt,bottom=2pt,boxsep=1pt,toptitle=1.5pt,bottomtitle=1.5pt,halign=left]
\begin{itemize}[nosep,leftmargin=*,itemsep=1pt,font=\scriptsize]
\item 垃圾桶 lājītǒng
\item 回收站 huíshōuzhàn
\item 工厂 gōngchǎng
\end{itemize}
\end{tcolorbox}
\begin{tcolorbox}[enhanced,colback=white,colframe=poppurple,boxrule=0.9pt,arc=3pt,title={actions 动作},fonttitle=\bfseries\scriptsize,coltitle=white,colbacktitle=poppurple,left=4pt,right=4pt,top=2pt,bottom=2pt,boxsep=1pt,toptitle=1.5pt,bottomtitle=1.5pt,halign=left]
\begin{itemize}[nosep,leftmargin=*,itemsep=1pt,font=\scriptsize]
\item 分类 fēnlèi
\item 清洗 qīngxǐ
\item 加工 jiāgōng
\end{itemize}
\end{tcolorbox}
\begin{tcolorbox}[enhanced,colback=white,colframe=poppink,boxrule=0.9pt,arc=3pt,title={buts 目的},fonttitle=\bfseries\scriptsize,coltitle=white,colbacktitle=poppink,left=4pt,right=4pt,top=2pt,bottom=2pt,boxsep=1pt,toptitle=1.5pt,bottomtitle=1.5pt,halign=left]
\begin{itemize}[nosep,leftmargin=*,itemsep=1pt,font=\scriptsize]
\item 保护环境 bǎohù huánjìng
\item 节约资源 jiéyuē zīyuán
\end{itemize}
\end{tcolorbox}
\end{tcbraster}

\legende{Carte mentale du lexique du recyclage autour de 回收 huíshōu}
\end{popfigure}

\begin{methode}[Mémoriser avec la carte mentale]
\begin{enumerate}
\item Lis le mot central et ses quatre branches à voix haute.
\item Cache la carte et récite chaque branche : trois matériaux, trois lieux, trois actions, deux buts.
\item Construis une phrase qui relie deux branches : \zh{我们把塑料放进垃圾桶。} (Wǒmen bǎ sùliào fàng jìn lājītǒng. — Nous mettons le plastique dans la poubelle.)
\item Termine par une phrase complète reliant action et but : \zh{我们分类垃圾，保护环境。} (Wǒmen fēnlèi lājī, bǎohù huánjìng. — Nous trions les déchets pour protéger l'environnement.)
\end{enumerate}
\end{methode}

\begin{savaistu}[Les poubelles de tri colorées en Chine]
Dans les grandes villes chinoises, les poubelles de tri ont des \cle{couleurs normalisées} : le bac \emph{bleu} pour les recyclables (\zh{可回收物} kě huíshōu wù), le bac \emph{rouge} pour les déchets dangereux (\zh{有害垃圾} yǒuhài lājī : piles, médicaments), le bac \emph{vert} pour les déchets de cuisine (\zh{厨余垃圾} chúyú lājī) et le bac \emph{gris} pour les autres déchets (\zh{其他垃圾} qítā lājī). Depuis 2019, des villes comme Shanghai rendent le tri obligatoire : trier est devenu une véritable \zh{习惯} (xíguàn, habitude) quotidienne, comme le marché ou le football chez nous !
\end{savaistu}

\courssub{Exercices d'assimilation}

\begin{exoresolu}[Association mot-image]
Associe chaque mot chinois à la bonne description française :
\begin{enumerate}
\item \zh{塑料} \hspace{1cm} a. endroit où l'on transforme les déchets
\item \zh{工厂} \hspace{1cm} b. matériau transparent qui se recycle bien
\item \zh{玻璃} \hspace{1cm} c. matériau léger, souvent en bouteilles
\item \zh{节约} \hspace{1cm} d. ne pas gaspiller
\end{enumerate}
\tcblower
\textbf{Solution détaillée :}
\begin{enumerate}
\item \zh{塑料} sùliào → c. Le plastique est le matériau léger des bouteilles (\zh{塑料瓶} sùliào píng).
\item \zh{工厂} gōngchǎng → a. L'usine est le lieu de la transformation (\zh{加工} jiāgōng).
\item \zh{玻璃} bōli → c'est b. Le verre est transparent et se recycle presque à l'infini.
\item \zh{节约} jiéyuē → d. Économiser, c'est-à-dire ne pas gaspiller : \zh{节约资源} (économiser les ressources).
\end{enumerate}
Astuce : pour vérifier, replace chaque mot dans une phrase du texte support ; si la phrase reste logique, l'association est bonne.
\end{exoresolu}

\begin{exercice}[À toi de jouer : complète avec le bon mot]
Complète chaque phrase avec un mot du glossaire (回收, 分类, 垃圾桶, 资源, 习惯, 环境) :
\begin{enumerate}
\item 我们应该把垃圾放进\_\_\_\_\_\_。(Nous devrions mettre les déchets dans la \_\_\_\_\_\_.)
\item 回收可以节约\_\_\_\_\_\_。(Recycler permet d'économiser les \_\_\_\_\_\_.)
\item 在雅温得，我们要养成垃圾分类的好\_\_\_\_\_\_。(À Yaoundé, nous devons prendre la bonne \_\_\_\_\_\_ de trier les déchets.)
\item 保护环境，从\_\_\_\_\_\_垃圾开始。(Protéger l'environnement commence par \_\_\_\_\_\_ les déchets.)
\end{enumerate}
\end{exercice}

\begin{corrige}[Exercice « À toi de jouer »]
\begin{enumerate}
\item \zh{垃圾桶} (lājītǒng) : \zh{我们应该把垃圾放进垃圾桶。} — « mettre dans la poubelle ».
\item \zh{资源} (zīyuán) : \zh{回收可以节约资源。} — collocation fixe \zh{节约资源}.
\item \zh{习惯} (xíguàn) : \zh{养成好习惯} (yǎngchéng hǎo xíguàn) = « prendre une bonne habitude », expression très fréquente.
\item \zh{分类} (fēnlèi) : \zh{从分类垃圾开始} — « commencer par trier les déchets » ; ici \zh{分类} est un verbe placé après la préposition \zh{从} (cóng, à partir de).
\end{enumerate}
\end{corrige}

\begin{aretenir}[Bilan du glossaire]
\begin{itemize}
\item Le mot central est \zh{回收} huíshōu (回 = retour + 收 = recevoir) ; il se distingue de \zh{利用} lìyòng (utiliser).
\item Les mots composés se devinent par leurs parties : 垃圾桶, 分类, 加工, 产品.
\item Attention aux tons : lājī (1-1), sùliào (4-4), bōli (1-léger).
\item Collocations à retenir : \zh{回收利用} (recycler et réutiliser), \zh{节约资源} (économiser les ressources), \zh{保护环境} (protéger l'environnement), \zh{养成好习惯} (prendre une bonne habitude).
\end{itemize}
Ce vocabulaire maîtrisé, tu es prêt à passer à l'étude des caractères : clés et ordre des traits, dans la section suivante.
\end{aretenir}

\coursec{Leçon 9 — Vocabulaire et compréhension de texte : processus du recyclage}

\courssub{Mise en route et découverte du texte}

\begin{experience}[Situation de communication : le tri à l'école]
Imaginez la scène : nous sommes au lycée de Buea, dans la cour de récréation. Trois poubelles de couleurs différentes viennent d'être installées : une \textbf{bleue} (papiers/cartons), une \textbf{jaune} (plastiques/métaux), une \textbf{verte} (déchets organiques). Un élève, \zh{阿明} (Ámíng), hésite devant sa bouteille d'eau vide en plastique et son paquet de biscuits en carton. Son camarade \zh{小红} (Xiǎohóng) lui explique comment trier.

Écoutez et lisez le dialogue suivant. Identifiez les mots-clés du tri sélectif.
\end{experience}

\begin{textebox}[Dialogue : Le tri sélectif à l'école]
\zh{阿明：小红，这个塑料瓶扔哪儿？}\\
\emph{Ámíng: Xiǎohóng, zhège sùliào píng rēng nǎr?}\\
« Xiaohong, où est-ce que je jette cette bouteille en plastique ? »\\[4pt]
\zh{小红：别乱扔！塑料瓶是可回收物，要扔进黄色的桶里。}\\
\emph{Xiǎohóng: Bié luàn rēng! Sùliào píng shì kě huíshōu wù, yào rēng jìn huángsè de tǒng lǐ.}\\
« Ne jette pas n'importe où ! La bouteille en plastique est un déchet recyclable, il faut la jeter dans le bac jaune. »\\[4pt]
\zh{阿明：那这个纸盒呢？}\\
\emph{Ámíng: Nà zhège zhǐ hé ne?}\\
« Et cette boîte en carton alors ? »\\[4pt]
\zh{小红：纸盒也是可回收物，扔进蓝色的桶。我们要做好垃圾分类，保护环境。}\\
\emph{Xiǎohóng: Zhǐ hé yě shì kě huíshōu wù, rēng jìn lánsè de tǒng. Wǒmen yào zuò hǎo lājī fēnlèi, bǎohù huánjìng.}\\
« La boîte en carton est aussi un déchet recyclable, jette-la dans le bac bleu. Nous devons bien faire le tri des déchets pour protéger l'environnement. »\\[4pt]
\zh{阿明：明白了。分类回收，变废为宝！}\\
\emph{Ámíng: Míngbái le. Fēnlèi huíshōu, biànfèiwéibǎo!}\\
« Compris. Tri et recyclage, transformer les déchets en trésors ! »
\end{textebox}

\begin{methode}[Comment aborder un texte chinois inconnu (3 étapes)]
\begin{enumerate}
\item \textbf{Repérage visuel} : Identifiez les caractères connus (ex: \zh{塑料}, \zh{纸}, \zh{环境}, \zh{保护}), la ponctuation chinoise (， 。 ？ ！) et la structure des phrases (Sujet + Verbe + Objet / Complément).
\item \textbf{Lecture phonétique} : Lisez le pinyin à voix haute en respectant les tons. Repérez les mots-outils (\zh{是}, \zh{要}, \zh{进}, \zh{里}, \zh{也}, \zh{我们}).
\item \textbf{Construction du sens} : Traduisez phrase par phrase. Utilisez le contexte (couleurs des poubelles, actions) pour deviner le vocabulaire nouveau (\zh{可回收物}, \zh{分类}, \zh{变废为宝}).
\end{enumerate}
\end{methode}

\courssub{Texte support : Le processus complet du recyclage}

\begin{textebox}[Article : Le voyage d'une bouteille en plastique — de la poubelle au nouveau produit]
\zh{塑料瓶的“新生”之路}\\
\emph{Sùliào píng de “xīnshēng” zhī lù}\\[4pt]
\zh{每天，雅温得的街道上产生大量垃圾。塑料瓶喝完后，如果随意丢弃，会污染土壤和河流，几百年也分解不了。}\\
\emph{Měitiān, Yǎwēndé de jiēdào shàng chǎnshēng dàliàng lājī. Sùliào píng hē wán hòu, rúguǒ suíyì diūqì, huì wūrǎn tǔrǎng hé héliú, jǐbǎi nián yě fēnjiě bù liǎo.}\\[4pt]
\zh{但是，如果我们做好垃圾分类，把塑料瓶扔进黄色回收桶，它就开始了一段“新生”之路：}\\
\emph{Dànshì, rúguǒ wǒmen zuò hǎo lājī fēnlèi, bǎ sùliào píng rēng jìn huángsè huíshōu tǒng, tā jiù kāishǐ le yī duàn “xīnshēng” zhī lù:}\\[4pt]
\zh{第一步：收集 (shōují) — 环卫工人把可回收物运送到回收站。}\\
\emph{Dì yī bù: shōují — huánwèi gōngrén bǎ kě huíshōu wù yùnsòng dào huíshōu zhàn.}\\[4pt]
\zh{第二步：分拣 (fēnjiǎn) — 机器和工人把塑料、纸、金属分开。}\\
\emph{Dì èr bù: fēnjiǎn — jīqì hé gōngrén bǎ sùliào, zhǐ, jīnshǔ fēnkāi.}\\[4pt]
\zh{第三步：加工 (jiāgōng) — 塑料瓶被清洗、粉碎、熔化，变成塑料颗粒 (kēlì)。}\\
\emph{Dì sān bù: jiāgōng — sùliào píng bèi qīngxǐ, fěnsuì, rónghuà, biàn chéng sùliào kēlì.}\\[4pt]
\zh{第四步：再制造 (zàizhìzào) — 颗粒被做成新的塑料制品，如衣服、椅子、新瓶子。}\\
\emph{Dì sì bù: zàizhìzào — kēlì bèi zuò chéng xīn de sùliào zhìpǐn, rú yīfu, yǐzi, xīn píngzi.}\\[4pt]
\zh{这就是“循环经济” (xúnhuán jīngjì) : 变废为宝，保护我们共同的家园。}\\
\emph{Zhè jiù shì “xúnhuán jīngjì” : biànfèiwéibǎo, bǎohù wǒmen gòngtóng de jiāyuán.}
\end{textebox}

\courssub{Glossaire du thème : Vocabulaire de base (HSK 3-4)}

\begin{center}
\begin{tabularx}{\linewidth}{|X|X|X|X|}\hline
\rowcolor{popblueL} \textbf{Caractères} & \textbf{Pinyin} & \textbf{Nature} & \textbf{Traduction} \\ \hline
\zh{垃圾} & lājī & nom & déchets, ordures \\ \hline
\zh{分类} & fēnlèi & verbe / nom & trier, classer / tri, classification \\ \hline
\zh{回收} & huíshōu & verbe & recycler, récupérer \\ \hline
\zh{可回收物} & kě huíshōu wù & nom & déchets recyclables \\ \hline
\zh{有害垃圾} & yǒuhài lājī & nom & déchets dangereux \\ \hline
\zh{厨余垃圾} & chúyú lājī & nom & déchets de cuisine (organiques) \\ \hline
\zh{其他垃圾} & qítā lājī & nom & autres déchets (résidus) \\ \hline
\zh{塑料} & sùliào & nom & plastique \\ \hline
\zh{纸} & zhǐ & nom & papier \\ \hline
\zh{金属} & jīnshǔ & nom & métal \\ \hline
\zh{玻璃} & bōlí & nom & verre \\ \hline
\zh{环境} & huánjìng & nom & environnement \\ \hline
\zh{污染} & wūrǎn & verbe / nom & polluer / pollution \\ \hline
\zh{保护} & bǎohù & verbe & protéger \\ \hline
\zh{环保} & huánbǎo & nom / adj & protection de l'environnement / écolo \\ \hline
\zh{循环经济} & xúnhuán jīngjì & nom & économie circulaire \\ \hline
\zh{变废为宝} & biànfèiwéibǎo & chengyu & transformer les déchets en trésors \\ \hline
\zh{收集} & shōují & verbe & collecter, ramasser \\ \hline
\zh{分拣} & fēnjiǎn & verbe & trier (séparer finement) \\ \hline
\zh{加工} & jiāgōng & verbe / nom & transformer, traiter / transformation \\ \hline
\zh{颗粒} & kēlì & nom & granulés (plastique) \\ \hline
\zh{再制造} & zàizhìzào & verbe & refabriquer, remanufacturer \\ \hline
\zh{环卫工人} & huánwèi gōngrén & nom & agent d'hygiène / éboueur \\ \hline
\zh{回收站} & huíshōu zhàn & nom & centre de tri / déchetterie \\ \hline
\zh{丢弃} & diūqì & verbe & jeter, abandonner \\ \hline
\zh{分解} & fēnjiě & verbe & se décomposer, dégrader \\ \hline
\zh{土壤} & tǔrǎng & nom & sol, terre \\ \hline
\zh{河流} & héliú & nom & rivière, fleuve \\ \hline
\zh{家园} & jiāyuán & nom & foyer, patrie (la Terre) \\ \hline
\end{tabularx}
\end{center}

\begin{aretenir}[Classificateurs utiles pour ce thème]
\begin{itemize}
\item \zh{个} (gè) : général (\zh{一个塑料瓶}).
\item \zh{吨} (dūn) : tonne (\zh{几吨垃圾}).
\item \zh{种} (zhǒng) : sorte, type (\zh{四种垃圾}).
\item \zh{步骤} (bùzhòu) : étape (\zh{四个步骤}).
\item \zh{桶} (tǒng) : bac, poubelle (\zh{黄色的桶}).
\item \zh{站} (zhàn) : station, centre (\zh{回收站}).
\end{itemize}
\end{aretenir}

\coursec{Clés (radicaux) et ordre des traits}

Dans la section précédente, nous avons découvert le glossaire du recyclage : \zh{垃圾} (lājī, déchets), \zh{环境} (huánjìng, environnement), \zh{塑料} (sùliào, plastique), \zh{纸} (zhǐ, papier), \zh{保护} (bǎohù, protéger). Ces caractères peuvent sembler difficiles à mémoriser. Bonne nouvelle : ils ne sont pas des dessins arbitraires ! Chacun est construit autour d'une \cle{clé} (ou \cle{radical}, \zh{部首} bùshǒu) qui donne un indice sur le sens, et s'écrit selon un \cle{ordre des traits} (\zh{笔顺} bǐshùn) précis. Cette section vous donne les deux clés de la réussite : comprendre la logique des caractères et les écrire correctement.

\courssub{Rappel : le rôle sémantique des radicaux}

\begin{definition}[La clé ou radical \zh{部首}]
La \cle{clé} (radical) est la partie d'un caractère chinois qui indique sa \textbf{famille de sens}. Elle sert aussi à classer les caractères dans les dictionnaires. L'autre partie du caractère donne souvent un indice sur la \textbf{prononciation} (on l'appelle la partie phonétique).
\end{definition}

Observons d'abord quelques caractères de notre glossaire :

\begin{exemplebox}[Observer avant de mémoriser]
\zh{垃、圾、境} : ces trois caractères contiennent tous \zh{土} (tǔ, la terre).\\
\emph{Lā, jī, jìng : zhè sān ge zì dōu hán yǒu “tǔ” zì páng.}\\
« 垃、圾、境 contiennent tous la clé terre : les déchets et l'environnement sont liés au sol. »\\[4pt]
\zh{保、护} : ces deux caractères contiennent des éléments humains (\zh{亻} l'homme, \zh{扌} la main).\\
\emph{Bǎo, hù : zhè liǎng ge zì dōu hán yǒu rén de yuánsù.}\\
« Protéger est une action que l'homme fait avec ses mains. »
\end{exemplebox}

\begin{savaistu}[Le pouvoir des radicaux]
Connaître les radicaux permet de deviner le sens d'environ 80 \% des caractères courants et de chercher un mot inconnu dans un dictionnaire chinois : on identifie la clé, on compte les traits restants, et on trouve le caractère dans l'index. C'est la méthode utilisée par tous les écoliers chinois — et elle marchera pour vous à Yaoundé comme à Pékin !
\end{savaistu}

\courssub{La clé \zh{土} (terre) : \zh{垃}, \zh{圾}, \zh{境}, \zh{塑}}

\begin{propriete}[La clé de la terre \zh{土}]
\zh{土} (tǔ) signifie « terre, sol ». Quand elle apparaît dans un caractère (à gauche, elle s'écrit un peu plus étroite et le dernier trait horizontal remonte légèrement), le sens du caractère est lié à la \textbf{terre, au sol, à un lieu ou à la matière}.
\end{propriete}

\textbf{1. \zh{垃} (lā) — 8 traits.} Décomposition : \zh{土} (terre) + \zh{立} (lì, se tenir debout). Ordre des traits : on écrit d'abord la clé \zh{土} à gauche (horizontal, vertical, horizontal remontant), puis \zh{立} à droite (point, horizontal, point, petit trait, horizontal final).

\textbf{2. \zh{圾} (jī) — 6 traits.} Décomposition : \zh{土} (terre) + \zh{及} (jí, atteindre — indice de prononciation). On écrit \zh{土} à gauche, puis \zh{及} à droite.

Le mot \zh{垃圾} (lājī, déchets, ordures) s'écrit donc avec deux caractères de la famille « terre » : les déchets finissent par retourner à la terre — belle image pour une leçon sur le recyclage !

\textbf{3. \zh{境} (jìng) — 14 traits.} Décomposition : \zh{土} (terre) + \zh{竟} (jìng, finalement — indice de prononciation). \zh{环境} (huánjìng, environnement) signifie littéralement « ce qui entoure la terre » : un mot très logique !

\textbf{4. \zh{塑} (sù) — 13 traits.} Ici la clé \zh{土} est placée \textbf{en bas} du caractère : \zh{朔} (indice phonétique shuò) au-dessus + \zh{土} en dessous. \zh{塑料} (sùliào, le plastique) : une matière façonnée à partir de la terre (pétrole).

\begin{attention}[Ne confondez pas \zh{土} et \zh{士}]
La clé \zh{土} (tǔ, terre) a son trait du \textbf{bas plus long} que le trait du haut. Le caractère \zh{士} (shì, lettré, soldat) a le trait du \textbf{haut plus long}. Erreur fréquente des francophones : écrire \zh{垃} avec \zh{士} à gauche. Retenez : la terre a une \emph{base large}, comme le sol sous nos pieds.
\end{attention}

\courssub{La clé \zh{王} (jade) : \zh{环}}

\begin{propriete}[La clé du jade \zh{王}]
Dans les caractères, \zh{王} à gauche n'est pas le « roi » (wáng) mais la clé du \textbf{jade} (forme abrégée de \zh{玉} yù). Les caractères qui la portent évoquent les pierres précieuses, les bijoux, les objets précieux, ou par extension la forme ronde/anneau.
\end{propriete}

\zh{环} (huán, anneau, cercle ; entourer) : \zh{王} (jade) à gauche + \zh{不} (bù — indice de prononciation approximatif) à droite, en 8 traits. À l'origine, \zh{环} désignait un anneau de jade. Aujourd'hui, dans \zh{环境} (huánjìng, environnement) et \zh{循环} (xúnhuán, cycle, recyclage en circuit), il garde l'idée de « cercle » : recycler, c'est faire tourner la matière en cercle !

\courssub{Les clés \zh{米} et \zh{纟} : \zh{料} et \zh{纸}}

\textbf{1. \zh{料} (liào) — 10 traits.} Décomposition : \zh{米} (mǐ, riz, grain) à gauche + \zh{斗} (dǒu, mesure à grain) à droite. Sens originel : « mesurer le grain », d'où « matière, matériau ». On le trouve dans \zh{塑料} (sùliào, plastique), \zh{材料} (cáiliào, matériau), \zh{饮料} (yǐnliào, boisson).

\textbf{2. \zh{纸} (zhǐ) — 7 traits.} Décomposition : \zh{纟} (sī, soie — clé du textile) à gauche + \zh{氏} (shì — indice phonétique) à droite. Pourquoi la soie ? Parce qu'avant l'invention du papier végétal, les Chinois écrivaient sur des tissus de soie, et les premiers papiers contenaient des fibres textiles. Le caractère garde la mémoire de cette histoire.

Ordre des 7 traits de \zh{纸} : (1) premier trait brisé de \zh{纟}, (2) deuxième trait brisé de \zh{纟}, (3) petit trait remontant de \zh{纟}, (4) petit trait oblique de \zh{氏}, (5) trait vertical-courbe, (6) trait horizontal, (7) trait oblique descendant final.

\courssub{Les clés de l'action humaine : \zh{保} et \zh{护}}

\textbf{1. \zh{保} (bǎo) — 9 traits.} Décomposition : \zh{亻} (rén, l'homme debout, variante de \zh{人}) à gauche + \zh{呆} (dāi — indice phonétique) à droite. Protéger est une action humaine : \zh{保护} (bǎohù, protéger), \zh{环保} (huánbǎo, protection de l'environnement, abréviation de \zh{环境保护}).

\textbf{2. \zh{护} (hù) — 7 traits.} Décomposition : \zh{扌} (shǒu, la main, variante de \zh{手}) à gauche + \zh{户} (hù, porte — indice phonétique) à droite. Protéger, c'est « tenir la porte avec la main » : une image parlante ! Autres caractères avec la main : \zh{扔} (rēng, jeter), \zh{捡} (jiǎn, ramasser), \zh{打} (dǎ, frapper).

\begin{exemplebox}[Du radical au sens : trois chaînes à retenir]
\zh{土} tǔ (terre) → \zh{垃圾} lājī (déchets) : les déchets retournent à la terre.\\
\emph{Tǔ (terre) → lājī (déchets).}\\[4pt]
\zh{纟} sī (soie) → \zh{纸} zhǐ (papier) : le papier est né des fibres textiles.\\
\emph{Sī (soie) → zhǐ (papier).}\\[4pt]
\zh{亻} rén (homme) + \zh{扌} shǒu (main) → \zh{保护} bǎohù (protéger) : l'homme agit de ses mains.\\
\emph{Rén + shǒu → bǎohù (protéger).}
\end{exemplebox}

\begin{methode}[Bien écrire un caractère : les quatre règles d'or]
Pour écrire correctement n'importe quel caractère chinois, respectez toujours l'ordre des traits :
\begin{enumerate}
\item \textbf{De haut en bas} : les traits du haut avant ceux du bas (\zh{塑} : d'abord \zh{朔}, puis \zh{土}).
\item \textbf{De gauche à droite} : la partie gauche avant la partie droite (\zh{垃} : d'abord \zh{土}, puis \zh{立} ; \zh{保} : d'abord \zh{亻}, puis \zh{呆}).
\item \textbf{L'extérieur avant l'intérieur} : pour les caractères encadrés, on trace le cadre, puis le contenu.
\item \textbf{La fermeture en dernier} : le trait qui ferme un cadre se trace en tout dernier.
\end{enumerate}
Règle supplémentaire : un trait horizontal précède généralement le vertical qui le croise (\zh{土} : horizontal, vertical, horizontal).
\end{methode}

\begin{center}
\begin{tabularx}{\linewidth}{|X|X|X|X|}\hline
\rowcolor{popblueL} Caractère & Clé (radical) & Décomposition & Lien de sens \\ \hline
\zh{垃} lā & \zh{土} terre & \zh{土} + \zh{立} & ordures au sol \\ \hline
\zh{圾} jī & \zh{土} terre & \zh{土} + \zh{及} & ordures au sol \\ \hline
\zh{环} huán & \zh{王} jade & \zh{王} + \zh{不} & anneau, cercle \\ \hline
\zh{境} jìng & \zh{土} terre & \zh{土} + \zh{竟} & territoire, lieu \\ \hline
\zh{塑} sù & \zh{土} terre (en bas) & \zh{朔} + \zh{土} & modeler la matière \\ \hline
\zh{料} liào & \zh{米} riz/grain & \zh{米} + \zh{斗} & mesurer la matière \\ \hline
\zh{纸} zhǐ & \zh{纟} soie & \zh{纟} + \zh{氏} & fibres textiles \\ \hline
\zh{保} bǎo & \zh{亻} homme & \zh{亻} + \zh{呆} & action humaine \\ \hline
\zh{护} hù & \zh{扌} main & \zh{扌} + \zh{户} & garder avec la main \\ \hline
\end{tabularx}
\end{center}

\begin{popfigure}
\begin{tikzpicture}[font=\small, node distance=0.8cm and 1.2cm]
% Décomposition 垃
\node[draw=popblue, fill=popblueL, rounded corners, align=center, minimum width=2.2cm] (la) at (0,0) {\zh{垃} lā};
\node[draw=popgreen, fill=popgreenL, rounded corners, align=center] (tu1) at (-1.6,-1.6) {\zh{土} tǔ\\ terre};
\node[draw=poporange, fill=poporangeL, rounded corners, align=center] (li) at (1.6,-1.6) {\zh{立} lì\\ se tenir};
\draw[-{Stealth}, thick, popblue] (la.south west) -- (tu1.north east);
\draw[-{Stealth}, thick, popblue] (la.south east) -- (li.north west);

% Décomposition 保
\node[draw=popblue, fill=popblueL, rounded corners, align=center, minimum width=2.2cm] (bao) at (5.5,0) {\zh{保} bǎo};
\node[draw=popgreen, fill=popgreenL, rounded corners, align=center] (ren) at (3.9,-1.6) {\zh{亻} rén\\ homme};
\node[draw=poporange, fill=poporangeL, rounded corners, align=center] (dai) at (7.1,-1.6) {\zh{呆} dāi\\ phonétique};
\draw[-{Stealth}, thick, popblue] (bao.south west) -- (ren.north east);
\draw[-{Stealth}, thick, popblue] (bao.south east) -- (dai.north west);

% Décomposition 环
\node[draw=popblue, fill=popblueL, rounded corners, align=center, minimum width=2.2cm] (huan) at (11,0) {\zh{环} huán};
\node[draw=popgreen, fill=popgreenL, rounded corners, align=center] (wang) at (9.4,-1.6) {\zh{王} yù\\ jade};
\node[draw=poporange, fill=poporangeL, rounded corners, align=center] (bu) at (12.6,-1.6) {\zh{不} bù\\ phonétique};
\draw[-{Stealth}, thick, popblue] (huan.south west) -- (wang.north east);
\draw[-{Stealth}, thick, popblue] (huan.south east) -- (bu.north west);

% Ordre des traits 纸
\node[draw=poppurple, fill=poppinkL, rounded corners, align=center, minimum width=3.5cm] (zhi) at (5.5,-4.5) {\zh{纸} zhǐ — 7 traits};
\node[draw=poppurple, fill=poppinkL, rounded corners, align=left, text width=5cm] (traits) at (5.5,-7) {
1--3 : \zh{纟} (soie) à gauche\\
4 : oblique de \zh{氏}\\
5 : vertical-courbe\\
6 : horizontal\\
7 : oblique descendant
};
\draw[-{Stealth}, thick, poppurple] (zhi.south) -- (traits.north);
\end{tikzpicture}
\legende{Décomposition de \zh{垃}, \zh{保}, \zh{环} et ordre des 7 traits de \zh{纸}}
\end{popfigure}

\begin{textebox}[Les caractères racontent le recyclage]
\zh{你看，“垃圾”两个字都有“土”字旁。垃圾最后回到土里。}\\
\emph{Nǐ kàn, “lājī” liǎng ge zì dōu yǒu “tǔ” zì páng. Lājī zuìhòu huí dào tǔ lǐ.}\\
« Regarde : les deux caractères de « déchet » portent la clé de la terre. Les déchets finissent par retourner à la terre. »\\[4pt]
\zh{“纸”字有“纟”旁，因为以前的纸和布有关系。}\\
\emph{“Zhǐ” zì yǒu “sī” páng, yīnwèi yǐqián de zhǐ hé bù yǒu guānxi.}\\
« Le caractère « papier » porte la clé de la soie, car autrefois le papier était lié au tissu. »\\[4pt]
\zh{“保护”两个字，一个有“人”，一个有“手”：保护环境是我们人的手做的事。}\\
\emph{“Bǎohù” liǎng ge zì, yī ge yǒu “rén”, yī ge yǒu “shǒu” : bǎohù huánjìng shì wǒmen rén de shǒu zuò de shì.}\\
« Dans « protéger », un caractère a « l'homme », l'autre « la main » : protéger l'environnement, c'est l'action des mains de l'homme. »\\[4pt]
\zh{学汉字的时候，先看部首，你就能猜出字的意思，也能查字典。}\\
\emph{Xué Hànzì de shíhou, xiān kàn bùshǒu, nǐ jiù néng cāi chū zì de yìsi, yě néng chá zìdiǎn.}\\
« Quand on apprend les caractères, si l'on regarde d'abord le radical, on peut deviner le sens du caractère et chercher dans le dictionnaire. »
\end{textebox}

\begin{exoresolu}[Retrouver la clé]
Énoncé : indiquez la clé (radical) de chaque caractère du glossaire : \zh{垃}, \zh{纸}, \zh{护}, \zh{料}, \zh{境}.
\tcblower
Solution détaillée :
\begin{itemize}
\item \zh{垃} : clé \zh{土} (terre) à gauche — les ordures sont sur le sol.
\item \zh{纸} : clé \zh{纟} (soie) à gauche — le papier ancien était lié aux textiles.
\item \zh{护} : clé \zh{扌} (main) à gauche — protéger avec la main.
\item \zh{料} : clé \zh{米} (riz, grain) à gauche — mesurer la matière.
\item \zh{境} : clé \zh{土} (terre) à gauche — un territoire, un lieu.
\end{itemize}
Méthode : la clé est le plus souvent à \textbf{gauche} ou \textbf{en haut} ; dans \zh{塑}, elle est exceptionnellement en bas (\zh{土}).
\end{exoresolu}

\begin{exercice}[À vous de jouer : Enquête dictionnaire]
Ouvrez un dictionnaire chinois (papier ou en ligne type Pleco, MDBG) et retrouvez la clé de ces cinq caractères du glossaire : \zh{瓶} (píng, bouteille), \zh{盒} (hé, boîte), \zh{回} (huí, retourner), \zh{类} (lèi, catégorie), \zh{收} (shōu, ramasser). Notez pour chacun : la clé, le nombre de traits restants (hors clé) et le sens de la clé.
\end{exercice}

\begin{corrige}[Exercice : Enquête dictionnaire]
\begin{itemize}
\item \zh{瓶} (píng) : clé \zh{瓦} (wǎ, tuile/poterie), 6 traits restants. Sens : récipient en terre cuite.
\item \zh{盒} (hé) : clé \zh{皿} (mǐn, récipient) \textbf{en bas}, 6 traits restants. Sens : contenant.
\item \zh{回} (huí) : clé \zh{囗} (wéi, enceinte), 4 traits restants. Sens : revenir à l'intérieur.
\item \zh{类} (lèi) : clé \zh{米} (mǐ, grain) à gauche, 8 traits restants. Sens : trier les grains par catégorie.
\item \zh{收} (shōu) : clé \zh{攵} (bō/pū, action de taper légèrement) \textbf{à droite}, 4 traits restants. Sens : action de rassembler.
\end{itemize}
Bravo : vous savez maintenant chercher n'importe quel caractère !
\end{corrige}

\begin{aretenir}[L'essentiel de la section : Radicaux et Écriture]
\begin{itemize}
\item La \cle{clé} (radical) indique la famille de sens ; l'autre partie guide souvent la prononciation.
\item Famille de la terre \zh{土} : \zh{垃、圾、境、塑} ; attention à ne pas la confondre avec \zh{士} (trait du haut long).
\item \zh{环} porte la clé du jade \zh{王} (anneau) ; \zh{料} la clé du grain \zh{米} ; \zh{纸} la clé de la soie \zh{纟}.
\item \zh{保} (homme \zh{亻}) et \zh{护} (main \zh{扌}) : protéger est une action humaine.
\item Ordre des traits : haut → bas, gauche → droite, extérieur → intérieur, fermeture en dernier.
\end{itemize}
\end{aretenir}

\coursec{Compréhension détaillée du texte}

\courssub{Questions de compréhension globale et fine}

\begin{exercice}[Compréhension de l'article « Le voyage d'une bouteille »]
Répondez en français (ou en chinois pour les questions 5-6) en vous basant strictement sur le texte.
\begin{enumerate}
\item Selon le texte, que se passe-t-il si une bouteille en plastique est jetée n'importe où dans la nature à Yaoundé ?
\item Citez les quatre étapes successives du processus de recyclage décrites dans l'article.
\item Quel est le rôle des \zh{环卫工人} (huánwèi gōngrén) à l'étape 1 ?
\item Quelle transformation subissent les bouteilles à l'étape 3 (\zh{加工}) ? Donnez le terme chinois pour le produit intermédiaire obtenu.
\item Traduisez en chinois : « L'économie circulaire transforme les déchets en trésors. »
\item Trouvez dans le texte le mot qui signifie « fondre » (verbe) et écrivez-le en caractères + pinyin.
\end{enumerate}
\end{exercice}

\begin{corrige}[Exercice : Compréhension de l'article]
\begin{enumerate}
\item Elle pollue le sol et les rivières, et met plusieurs centaines d'années à se décomposer (\zh{污染土壤和河流，几百年也分解不了}).
\item 1. Collecte (\zh{收集}) 2. Tri/Séparation (\zh{分拣}) 3. Transformation/Traitement (\zh{加工}) 4. Refabrication (\zh{再制造}).
\item Ils transportent les déchets recyclables vers le centre de tri (\zh{运送到回收站}).
\item Elles sont lavées, broyées, fondues (\zh{清洗、粉碎、熔化}) et deviennent des granulés de plastique (\zh{塑料颗粒} / \zh{颗粒} kēlì).
\item \zh{循环经济变废为宝。} (Xúnhuán jīngjì biànfèiwéibǎo.)
\item \zh{熔化} (rónghuà).
\end{enumerate}
\end{corrige}

\courssub{Exercices de vocabulaire en contexte}

\begin{exercice}[Complétez les phrases avec les mots du glossaire (forme correcte)]
Utilisez : \zh{分类}, \zh{回收}, \zh{保护}, \zh{污染}, \zh{变废为宝}, \zh{收集}, \zh{加工}, \zh{颗粒}.
\begin{enumerate}
\item 为了 \_\_\_\_\_\_ 环境，我们必须做好垃圾 \_\_\_\_\_\_。
\item 塑料瓶被 \_\_\_\_\_\_ 后，变成了塑料 \_\_\_\_\_\_。
\item \_\_\_\_\_\_ 利用是循环经济的核心理念。
\item 乱扔电池会严重 \_\_\_\_\_\_ 土壤。
\item 环卫工人每天早早就开始 \_\_\_\_\_\_ 可回收物。
\end{enumerate}
\end{exercice}

\begin{corrige}[Exercice : Complétion]
\begin{enumerate}
\item 为了 \zh{保护} (bǎohù) 环境，我们必须做好垃圾 \zh{分类} (fēnlèi)。
\item 塑料瓶被 \zh{加工} (jiāgōng) 后，变成了塑料 \zh{颗粒} (kēlì)。
\item \zh{回收} (huíshōu) 利用是循环经济的核心理念。
\item 乱扔电池会严重 \zh{污染} (wūrǎn) 土壤。
\item 环卫工人每天早早就开始 \zh{收集} (shōují) 可回收物。
\end{enumerate}
\end{corrige}

\coursec{Collocations et réemploi du lexique}

\courssub{Structures verbales fréquentes : \zh{把} (bǎ) et \zh{被} (bèi)}

Le texte support utilise deux structures grammaticales majeures pour décrire le processus (HSK 3-4) :

\begin{propriete}[Structure \zh{把} (bǎ) : Action intentionnelle sur l'objet]
\zh{Sujet + 把 + Objet + Verbe + Complément de résultat / direction}
\emph{Emploi} : Le sujet agit \textbf{activement} sur l'objet défini/connnu et le manipule (jeter, envoyer, faire).
\begin{center}
\begin{tabularx}{\linewidth}{|X|X|X|}\hline
\rowcolor{popblueL} Structure & Exemple (Caractères + Pinyin) & Traduction \\ \hline
\zh{把} + O + V + 进/到 & \zh{把塑料瓶扔进黄桶} (bǎ sùliào píng rēng jìn huáng tǒng) & Jeter la bouteille dans le bac jaune \\ \hline
\zh{把} + O + V + 好 & \zh{把垃圾分类做好} (bǎ lājī fēnlèi zuò hǎo) & Bien faire le tri des déchets \\ \hline
\zh{把} + O + 送/运 + 到 & \zh{把可回收物运送到回收站} (bǎ kě huíshōu wù yùnsòng dào huíshōu zhàn) & Transporter les recyclables au centre \\ \hline
\end{tabularx}
\end{center}
\end{propriete}

\begin{attention}[Piège \zh{把}]
L'objet après \zh{把} doit être \textbf{connu, spécifique} ou \textbf{défini} (le mien, celui-là, les bouteilles \emph{que j'ai bues}). On ne dit pas \zh{*把一个瓶子扔进去} (un bouteille quelconque) mais \zh{把这个瓶子扔进去} (cette bouteille-ci).
\end{attention}


\begin{propriete}[Structure \zh{被} (bèi) : Passif / Subir l'action]
\zh{Sujet (patient) + 被 + Agent (optionnel) + Verbe + Complément}
\emph{Emploi} : Le sujet \textbf{subit} l'action. Très utilisé pour décrire les étapes industrielles (étape 3-4).
\begin{center}
\begin{tabularx}{\linewidth}{|X|X|X|}\hline
\rowcolor{poporangeL} Structure & Exemple (Caractères + Pinyin) & Traduction \\ \hline
\zh{被} + V + 成/为 & \zh{被做成新制品} (bèi zuò chéng xīn zhìpǐn) & Être transformé en nouveaux produits \\ \hline
\zh{被} + V + 完/好 & \zh{被清洗干净} (bèi qīngxǐ gānjìng) & Être lavé proprement \\ \hline
(Sans agent) & \zh{塑料瓶被熔化} (sùliào píng bèi rónghuà) & La bouteille est fondue \\ \hline
\end{tabularx}
\end{center}
\end{propriete}

\begin{exercice}[Transformation \zh{把} $\leftrightarrow$ \zh{被}]
Transformez les phrases actives (\zh{把}) en passives (\zh{被}) et inversement.
\begin{enumerate}
\item \zh{工人把垃圾分拣好。} (Les ouvriers trient bien les déchets.)
\item \zh{塑料颗粒被做成新椅子。} (Les granulés sont faits en nouvelles chaises.)
\item \zh{我们把旧衣服送去回收站。} (Nous envoyons les vieux vêtements au centre de tri.)
\end{enumerate}
\end{exercice}

\begin{corrige}[Exercice : Transformation]
\begin{enumerate}
\item \zh{垃圾被工人分拣好。} (Les déchets sont bien triés par les ouvriers.)
\item \zh{（工人）把塑料颗粒做成新椅子。} ((Les ouvriers) font les granulés en nouvelles chaises.)
\item \zh{旧衣服被我们送去回收站。} (Les vieux vêtements sont envoyés par nous au centre de tri.)
\end{enumerate}
\end{corrige}

\courssub{Collocations lexicales (搭配 dāpèi) à mémoriser}

\begin{center}
\begin{tabularx}{\linewidth}{|X|X|X|}\hline
\rowcolor{popgreenL} \textbf{Verbe + Objet (VO)} & \textbf{Adjectif + Nom / Nom + Nom} & \textbf{Chengyu / Expressions figées} \\ \hline
\zh{做好分类} (faire le tri) & \zh{可回收物} (recyclables) & \zh{变废为宝} (déchets $\to$ trésors) \\ \hline
\zh{进行回收} (effectuer le recyclage) & \zh{有害垃圾} (déchets dangereux) & \zh{循环经济} (économie circulaire) \\ \hline
\zh{减少污染} (réduire la pollution) & \zh{厨余垃圾} (déchets cuisine) & \zh{绿水青山} (eaux vertes monts bleus = environnement préservé) \\ \hline
\zh{保护环境} (protéger l'env.) & \zh{环保意识} (conscience écolo) & \zh{可持续发展} (développement durable) \\ \hline
\zh{清洗粉碎} (laver \& broyer) & \zh{塑料颗粒} (granulés plastique) & \zh{变废为宝} \\ \hline
\end{tabularx}
\end{center}

\begin{exercice}[Traduction thème/version : Collocations]
Traduisez en chinois :
\begin{enumerate}
\item Améliorer la conscience de protection de l'environnement.
\item Les déchets de cuisine doivent être séparés.
\item C'est un projet de développement durable.
\item Nous devons réduire l'utilisation de sacs en plastique (\zh{塑料袋} sùliào dài).
\end{enumerate}
\end{exercice}

\begin{corrige}[Exercice : Traduction]
\begin{enumerate}
\item \zh{提高环保意识} (tí gāo huánbǎo yìshi).
\item \zh{厨余垃圾必须分开} (chúyú lājī bìxū fēnkāi) / \zh{做好厨余垃圾分类}.
\item \zh{这是一个可持续发展项目} (zhè shì yī gè kěchíxù fāzhǎn xiàngmù).
\item \zh{我们必须减少使用塑料袋} (wǒmen bìxū jiǎnshǎo shǐyòng sùliào dài).
\end{enumerate}
\end{corrige}

\coursec{Production guidée et évaluation}

\courssub{Expression écrite : Affiche de tri pour le lycée}

\begin{methode}[Rédiger une affiche de tri sélectif (格式 \zh{格式} géshì)]
\begin{enumerate}
\item \textbf{Titre} : Court, impératif, au centre. Ex: \zh{垃圾分类，从我做起} (Tri des déchets, ça commence par moi).
\item \textbf{Quatre catégories} : Utilisez les termes officiels chinois (couleurs) : \zh{可回收物} (蓝), \zh{有害垃圾} (红), \zh{厨余垃圾} (绿), \zh{其他垃圾} (灰).
\item \textbf{Exemples concrets} : 3-4 items par catégorie (vocabulaire de la leçon + contexte camerounais : \zh{椰子壳} yézi kécoque de noix de coco, \zh{电池} diànchí piles, \zh{快餐盒} kuàicān hé barquettes).
\item \textbf{Slogan final} : \zh{保护环境，人人有责} (Protéger l'env., responsabilité de chacun).
\end{enumerate}
\end{methode}

\begin{exercice}[Production écrite : Votre affiche]
Concevez une affiche A4 (sur papier ou numérique) pour le hall de votre lycée à Douala ou Yaoundé. Respectez la méthode ci-dessus. Écrivez \textbf{uniquement en caractères chinois} (pinyin autorisé en petit sous les mots difficiles). Utilisez des puces \zh{•} ou tirets.
\end{exercice}

\begin{exemplebox}[Modèle d'affiche (extrait)]
\zh{\textbf{垃圾分类，从我做起}}\\[4pt]
\zh{\textcolor{blue}{可回收物 (蓝桶)}：塑料瓶、纸盒、易拉罐、旧书本}\\
\zh{\textcolor{red}{有害垃圾 (红桶)}：电池、药品、灯管、油漆桶}\\
\zh{\textcolor{green}{厨余垃圾 (绿桶)}：剩饭菜、果皮、蛋壳、椰子壳、骨头}\\
\zh{\textcolor{gray}{其他垃圾 (灰桶)}：卫生纸、烟蒂、陶瓷碎片、污损塑料袋}\\[4pt]
\zh{保护环境，人人有责！变废为宝，利国利民！}
\end{exemplebox}

\courssub{Expression orale : Exposé comparatif Cameroun — Chine}

\begin{experience}[Débat guidé : « Le tri à Yaoundé et à Shanghai »]
\textbf{Consigne} : En binôme, préparez un exposé oral de 2 minutes (1 min par élève). Utilisez le vocabulaire de la leçon et les structures \zh{把}/\zh{被}.
\textbf{Plan suggéré} :
\begin{enumerate}
\item \textbf{Introduction} : \zh{大家好，今天我们谈谈垃圾分类。} (Bonjour, aujourd'hui on parle du tri.)
\item \textbf{Chine (Shanghai)} : \zh{上海2019年强制分类，四色桶，罚款。} (Shanghai 2019 tri obligatoire, 4 bacs, amendes.)
\item \textbf{Cameroun (Yaoundé/Douala)} : \zh{喀麦隆正在推广，HYSACAM收集，但分类还不普及。} (Cameroun en promotion, HYSACAM collecte, mais tri pas encore généralisé.)
\item \textbf{Comparaison} : \zh{中国做得比喀麦隆早，但喀麦隆年轻人很有热情。} (Chine plus tôt, mais jeunes Camerounais très motivés.)
\item \textbf{Proposition} : \zh{我们要从学校做起，把塑料瓶、纸分类回收。} (Commencer à l'école, trier bouteilles et papier.)
\item \textbf{Conclusion} : \zh{保护环境，变废为宝，共同努力！}
\end{enumerate}
\end{experience}

\begin{aretenir}[Grille d'auto-évaluation orale (Niveau Terminale / HSK 4)]
\begin{center}
\begin{tabularx}{\linewidth}{|X|X|X|}\hline
\rowcolor{poppurpleL} \textbf{Critère} & \textbf{Atteint (✓)} & \textbf{À travailler} \\ \hline
Vocabulaire thème (15+ mots) & & \\ \hline
Structure \zh{把} correcte & & \\ \hline
Structure \zh{被} correcte & & \\ \hline
Tons justes (surtout 3e ton) & & \\ \hline
Fluidité / Liaison phrases & & \\ \hline
Contenu culturel (CMR/CHN) & & \\ \hline
\end{tabularx}
\end{center}
\end{aretenir}

\courssub{Éléments socioculturels : Cameroun — Chine}

\begin{savaistu}[Gestion des déchets : Regards croisés]
\begin{itemize}
\item \textbf{Chine (上海/北京)} : Depuis juillet 2019, Shanghai impose le tri obligatoire en 4 catégories (secs/humides/nuisibles/recyclables) avec amendes (\zh{罚款} fákuǎn) pour les particuliers (50-200 CNY) et entreprises (50 000-500 000 CNY). Taux de tri > 90\% dans les quartiers pilotes. Applications mobiles (\zh{分类小程序}) pour scanner le code-barres du déchet.
\item \textbf{Cameroun (Yaoundé/Douala)} : La société \zh{HYSACAM} (Hygiène et Salubrité du Cameroun) gère la collecte. Le tri à la source est encore peu pratiqué par les ménages. Décharges principales : \zh{Nkolfoulou} (Yaoundé), \zh{Pk10} (Douala). Initiatives récentes : ONG, clubs environnement lycées, « Journée mondiale de l'environnement » (5 juin), projets de centres de tri pilotés par le MINEPDED.
\item \textbf{Point commun} : Les deux pays visent l'\zh{循环经济} (économie circulaire). La Chine exporte sa technologie de traitement (usines \zh{变废为宝}) vers l'Afrique via le \zh{中非合作论坛} (FOCAC).
\end{itemize}
\end{savaistu}

\begin{exercice}[Expression écrite : Courriel formel (HSK 4)]
Vous êtes délégué éco-citoyen de votre lycée. Écrivez un courriel en chinois (80-100 caractères) au Proviseur (\zh{校长} Xiàozhǎng) pour proposer l'installation de poubelles de tri. Utilisez : \zh{建议} (jiànyì, suggérer), \zh{安装} (ānzhuāng, installer), \zh{分类垃圾桶} (fēnlèi lājī tǒng), \zh{保护校园环境} (bǎohù xiàoyuán huánjìng).
\end{exercice}

\begin{corrige}[Exercice : Courriel formel]
\zh{尊敬的校长：}\\
\zh{您好！}\\
\zh{我是高三班的环保委员。建议学校在教学楼和食堂安装四色分类垃圾桶（蓝、红、绿、灰），并组织一次分类知识讲座。这能帮助同学们养成分类习惯，保护校园环境，响应国家环保号召。}\\
\zh{恳请批准指导。}\\
\zh{学生：[Votre Nom]}\\
\zh{日期：202X年X月X日}
\end{corrige}

\begin{aretenir}[Synthèse de la Leçon 9 — Module 2 Environnement]
\begin{itemize}
\item \textbf{Vocabulaire} : Maîtriser les 25 mots du glossaire (catégories, verbes processus, termes techniques \zh{颗粒}, \zh{循环经济}).
\item \textbf{Caractères} : Écrire correctement les 9 caractères clés (\zh{垃圾环境塑料纸保护} + \zh{收分加造}) en respectant l'ordre des traits et les radicaux (\zh{土, 王, 米, 纟, 亻, 扌}).
\item \textbf{Grammaire} : Utiliser \zh{把} (action active sur objet connu) et \zh{被} (passif industriel) pour décrire un processus.
\item \textbf{Compréhension} : Lire un texte explicatif structuré (étapes \zh{第一步}...\zh{第四步}) et en extraire l'information.
\item \textbf{Production} : Rédiger une affiche impérative + un courriel formel ; présenter oralement un comparatif CMR/CHN.
\item \textbf{Culture} : Connaître les systèmes de tri (Shanghai 4 couleurs vs Cameroun HYSACAM) et le vocabulaire associé (\zh{罚款}, \zh{强制}, \zh{推广}).
\end{itemize}
\end{aretenir}

\coursec{Compréhension détaillée du texte}

Dans la section précédente, nous avons appris à reconnaître les clés (radicaux) et l'ordre des traits des caractères du thème du recyclage, comme \zh{垃} (clé de la terre \zh{土}) ou \zh{璃} (clé du jade \zh{王}). Vous savez maintenant \emph{écrire} les mots. Il est temps de passer à l'étape suivante : \emph{comprendre} le texte phrase par phrase, en démontant ses mécanismes grammaticaux. C'est exactement la compétence exigée à l'épreuve de chinois du Baccalauréat : lire un texte descriptif, repérer les informations et répondre à des questions.

\courssub{Relecture du texte support}

Relisons d'abord le texte étudié dans cette leçon. Lisez-le à voix haute en suivant le pinyin, puis vérifiez votre compréhension avec la traduction.

\begin{textebox}[Le recyclage des déchets — 垃圾的回收利用]
为了保护环境，很多城市开始回收利用垃圾。\\
\emph{Wèile bǎohù huánjìng, hěn duō chéngshì kāishǐ huíshōu lìyòng lājī.}\\
« Pour protéger l'environnement, beaucoup de villes ont commencé à recycler les déchets. »\\[0.4em]
人们先把垃圾放进不同的垃圾桶，然后工人把垃圾送到回收站。\\
\emph{Rénmen xiān bǎ lājī fàng jìn bùtóng de lājī tǒng, ránhòu gōngrén bǎ lājī sòng dào huíshōuzhàn.}\\
« Les gens mettent d'abord les déchets dans des poubelles différentes, puis les ouvriers les transportent au centre de recyclage. »\\[0.4em]
在回收站，塑料、纸和玻璃被分类、清洗，再送到工厂。\\
\emph{Zài huíshōuzhàn, sùliào, zhǐ hé bōli bèi fēnlèi, qīngxǐ, zài sòng dào gōngchǎng.}\\
« Au centre de recyclage, le plastique, le papier et le verre sont triés, nettoyés, puis envoyés à l'usine. »\\[0.4em]
最后，旧东西变成新产品，比如新瓶子、新纸和新玩具。\\
\emph{Zuìhòu, jiù dōngxi biàn chéng xīn chǎnpǐn, bǐrú xīn píngzi, xīn zhǐ hé xīn wánjù.}\\
« Enfin, les vieux objets deviennent de nouveaux produits, par exemple de nouvelles bouteilles, du nouveau papier et de nouveaux jouets. »\\[0.4em]
回收利用不仅可以节约资源，还可以减少污染，让我们的生活环境更干净、更美丽。\\
\emph{Huíshōu lìyòng bùjǐn kěyǐ jiéyuē zīyuán, hái kěyǐ jiǎnshǎo wūrǎn, ràng wǒmen de shēnghuó huánjìng gèng gānjìng, gèng měilì.}\\
« Le recyclage permet non seulement d'économiser les ressources, mais aussi de réduire la pollution, et rend notre cadre de vie plus propre et plus beau. »
\end{textebox}

\begin{methode}[Relecture phrase par phrase : la bonne démarche]
\begin{enumerate}
\item \textbf{Lire} la phrase chinoise en entier, sans traduire mot à mot.
\item \textbf{Repérer} les mots connus du glossaire et les connecteurs (\zh{先}, \zh{然后}, \zh{再}, \zh{最后}).
\item \textbf{Identifier} la structure grammaticale : qui fait quoi ? Y a-t-il un \zh{把} ou un \zh{被} ?
\item \textbf{Traduire} en français naturel, pas en français mot à mot.
\item \textbf{Vérifier} : ma traduction raconte-t-elle une étape logique du processus ?
\end{enumerate}
\end{methode}

\courssub{Les connecteurs de séquence : 先…然后…再…最后…}

\courssub{Observation}

Regardez les quatre phrases centrales du texte. Soulignez mentalement les mots qui indiquent l'ordre des étapes :

\begin{exemplebox}[Les quatre étapes du texte]
\zh{人们先把垃圾放进不同的垃圾桶。} \emph{Rénmen xiān bǎ lājī fàng jìn bùtóng de lājī tǒng.} « Les gens mettent d'abord les déchets dans des poubelles différentes. »\\
\zh{然后工人把垃圾送到回收站。} \emph{Ránhòu gōngrén bǎ lājī sòng dào huíshōuzhàn.} « Ensuite, les ouvriers transportent les déchets au centre de recyclage. »\\
\zh{塑料、纸和玻璃被分类、清洗，再送到工厂。} \emph{Sùliào, zhǐ hé bōli bèi fēnlèi, qīngxǐ, zài sòng dào gōngchǎng.} « Le plastique, le papier et le verre sont triés, nettoyés, puis envoyés à l'usine. »\\
\zh{最后，旧东西变成新产品。} \emph{Zuìhòu, jiù dōngxi biàn chéng xīn chǎnpǐn.} « Enfin, les vieux objets deviennent de nouveaux produits. »
\end{exemplebox}

Chaque étape est introduite par un connecteur différent. Ces quatre mots forment une \cle{chaîne de séquence} qui ordonne le processus du début à la fin.

\begin{propriete}[La chaîne 先…然后…再…最后…]
Pour décrire un processus en plusieurs étapes, le chinois utilise la chaîne :
\begin{center}
\zh{先} (xiān, d'abord) + étape 1，\zh{然后} (ránhòu, ensuite) + étape 2，\zh{再} (zài, puis) + étape 3，\zh{最后} (zuìhòu, enfin) + étape finale。
\end{center}
\begin{itemize}
\item \zh{先} se place \textbf{avant le verbe}, à l'intérieur de la phrase : \zh{我们先分类。} « Nous trions d'abord. »
\item \zh{然后} et \zh{最后} se placent \textbf{en tête de phrase}, suivis ou non d'une virgule.
\item \zh{再} se place \textbf{avant le verbe} : \zh{再送到工厂} « puis envoyés à l'usine ».
\item Cette chaîne est \textbf{très fréquente dans les textes descriptifs au Bac} (recettes, processus, emploi du temps, récit d'une journée). La repérer, c'est découper le texte en étapes claires.
\end{itemize}
\end{propriete}

\begin{attention}[Ne confondez pas 再 et 又]
\zh{再} (zài) indique une action \textbf{suivante ou future} dans une séquence : « puis, ensuite ». \zh{又} (yòu) indique une \textbf{répétition} : « de nouveau, encore ».\\
\zh{洗干净以后，再送到工厂。} \emph{Xǐ gānjìng yǐhòu, zài sòng dào gōngchǎng.} « Après avoir nettoyé, on envoie ensuite à l'usine. » (étape suivante)\\
\zh{他又迟到了。} \emph{Tā yòu chídào le.} « Il est encore arrivé en retard. » (répétition)\\
Les francophones traduisent souvent les deux par « encore » : méfiez-vous !
\end{attention}

\begin{popfigure}
\begin{tikzpicture}[node distance=0.6cm, >=Stealth]
\node[draw=popblue, fill=popblueL, rounded corners, minimum width=2.8cm, minimum height=1.2cm, align=center] (e1) {\zh{先}\\收集 \emph{shōují}};
\node[draw=popgreen, fill=popgreenL, rounded corners, minimum width=2.8cm, minimum height=1.2cm, align=center, right=of e1] (e2) {\zh{然后}\\送到回收站 \emph{sòng dào huíshōuzhàn}};
\node[draw=poporange, fill=poporangeL, rounded corners, minimum width=2.8cm, minimum height=1.2cm, align=center, right=of e2] (e3) {\zh{再}\\分类、清洗 \emph{fēnlèi, qīngxǐ}};
\node[draw=poppurple, fill=poppurpleL, rounded corners, minimum width=2.8cm, minimum height=1.2cm, align=center, right=of e3] (e4) {\zh{最后}\\变成新产品 \emph{biàn chéng xīn chǎnpǐn}};
\draw[very thick, popdark, ->] (e1) -- (e2);
\draw[very thick, popdark, ->] (e2) -- (e3);
\draw[very thick, popdark, ->] (e3) -- (e4);
\end{tikzpicture}
\legende{Frise du processus de recyclage : les quatre connecteurs de séquence ordonnent les étapes.}
\end{popfigure}

\courssub{La structure passive avec 被}

\courssub{Observation : deux façons de dire la même chose}

Comparez ces deux phrases :

\begin{exemplebox}[Actif et passif]
\zh{工人把垃圾分类。} \emph{Gōngrén bǎ lājī fēnlèi.} « Les ouvriers trient les déchets. » (actif : l'ouvrier est le sujet)\\
\zh{垃圾被工人分类了。} \emph{Lājī bèi gōngrén fēnlèi le.} « Les déchets sont triés par les ouvriers. » (passif : les déchets sont le sujet)
\end{exemplebox}

Dans le texte, nous lisons : \zh{塑料、纸和玻璃被分类、清洗} « le plastique, le papier et le verre sont triés, nettoyés ». Pourquoi le passif ? Parce que ce qui nous intéresse, c'est le \textbf{sort des matériaux}, pas l'identité de celui qui fait l'action. Le français fait exactement le même choix : « sont triés, nettoyés ».

\begin{propriete}[Le passif avec 被 bèi]
Structure : \textbf{Objet (patient) + 被 + (agent) + Verbe + Complément (résultat, direction, 了...)}.
\begin{itemize}
\item \zh{垃圾被工人收集起来了。} \emph{Lājī bèi gōngrén shōují qǐlái le.} « Les déchets ont été collectés par les ouvriers. »
\item L'agent peut être omis s'il est inconnu ou sans importance : \zh{旧瓶子被送到工厂了。} \emph{Jiù píngzi bèi sòng dào gōngchǎng le.} « Les vieilles bouteilles ont été envoyées à l'usine. »
\item Le passif chinois s'emploie surtout pour des événements subis, souvent neutres ou défavorables (être trié, jeté, cassé, volé, réparé).
\item À l'oral, on rencontre aussi \zh{叫} (jiào) et \zh{让} (ràng) à la place de \zh{被}, mais à l'écrit et au Bac, utilisez \zh{被}.
\end{itemize}
\end{propriete}

\begin{attention}[被 exige un verbe « complet »]
Contrairement au français, le verbe qui suit \zh{被} ne peut presque jamais rester seul : il doit porter un \textbf{complément de résultat, de direction, 了 ou un objet}.
\begin{itemize}
\item Correct : \zh{垃圾被收集\cle{起来}了。} / \zh{纸被做\cle{成}了新本子。} / \zh{玻璃被清洗\cle{干净}了。}
\item Incomplet (sauf énumération) : \zh{垃圾被收集。} (sonne coupé).
\item \textbf{Exception} : dans une énumération d'actions subies par le même sujet, le dernier verbe porte la marque de complétude : \zh{被分类、清洗\cle{，再送到工厂}。} (le contexte porte la finalité).
\end{itemize}
Erreur typique des francophones : calquer « les déchets sont collectés » par \zh{垃圾被收集} sans complément. Ajoutez toujours \zh{起来}, \zh{成}, \zh{完}, \zh{好}, \zh{了} ou un complément de direction.
\end{attention}

\begin{center}
\begin{tabularx}{\linewidth}{|X|X|X|}\hline
\rowcolor{popblueL} Structure & Exemple (caractères + pinyin) & Traduction \\ \hline
Actif avec 把 & \zh{工人把垃圾送到回收站。} \emph{Gōngrén bǎ lājī sòng dào huíshōuzhàn.} & Les ouvriers transportent les déchets au centre de recyclage. \\ \hline
Passif avec agent & \zh{垃圾被工人送到回收站了。} \emph{Lājī bèi gōngrén sòng dào huíshōuzhàn le.} & Les déchets ont été transportés au centre de recyclage par les ouvriers. \\ \hline
Passif sans agent & \zh{玻璃被清洗干净了。} \emph{Bōli bèi qīngxǐ gānjìng le.} & Le verre a été nettoyé (bien propre). \\ \hline
Passif + 变成/做成 & \zh{旧东西被做成新产品。} \emph{Jiù dōngxi bèi zuò chéng xīn chǎnpǐn.} & Les vieux objets sont transformés en nouveaux produits. \\ \hline
\end{tabularx}
\end{center}

\courssub{Questions de compréhension}

\begin{methode}[Répondre à une question de compréhension]
\begin{enumerate}
\item \textbf{Repérer le mot-clé interrogatif} : \zh{哪里} (où), \zh{什么} (quoi), \zh{谁} (qui), \zh{为什么} (pourquoi), \zh{怎么} (comment).
\item \textbf{Localiser} dans le texte la phrase qui contient le même mot-clé ou son champ lexical.
\item \textbf{Citer ou reformuler} cette phrase : une bonne réponse s'appuie toujours sur le texte.
\item \textbf{Rédiger} une réponse complète en chinois, en reprenant les éléments de la question.
\end{enumerate}
\end{methode}

\begin{exoresolu}[Questions de compréhension sur le texte]
Énoncé : répondez en chinois aux cinq questions suivantes, en citant la phrase du texte qui justifie votre réponse.\\
1) \zh{人们把垃圾放进哪里？} 2) \zh{工人把垃圾送到哪里？} 3) \zh{在回收站，什么东西被分类？} 4) \zh{旧东西最后变成什么？} 5) \zh{回收利用有什么好处？}
\tcblower
\textbf{1)} Mot-clé : \zh{哪里} (où). Phrase du texte : \zh{人们先把垃圾放进不同的垃圾桶。}\\
Réponse : \zh{人们把垃圾放进不同的垃圾桶。} \emph{Rénmen bǎ lājī fàng jìn bùtóng de lājī tǒng.} « Les gens mettent les déchets dans des poubelles différentes. »\\[0.3em]
\textbf{2)} Mot-clé : \zh{哪里}. Phrase du texte : \zh{然后工人把垃圾送到回收站。}\\
Réponse : \zh{工人把垃圾送到回收站。} \emph{Gōngrén bǎ lājī sòng dào huíshōuzhàn.} « Les ouvriers transportent les déchets au centre de recyclage. »\\[0.3em]
\textbf{3)} Mot-clé : \zh{什么} (quoi). Phrase du texte : \zh{塑料、纸和玻璃被分类、清洗。}\\
Réponse : \zh{塑料、纸和玻璃被分类。} \emph{Sùliào, zhǐ hé bōli bèi fēnlèi.} « Le plastique, le papier et le verre sont triés. »\\[0.3em]
\textbf{4)} Mot-clé : \zh{什么}. Phrase du texte : \zh{最后，旧东西变成新产品。}\\
Réponse : \zh{旧东西最后变成新产品，比如新瓶子、新纸和新玩具。} \emph{Jiù dōngxi zuìhòu biàn chéng xīn chǎnpǐn, bǐrú xīn píngzi, xīn zhǐ hé xīn wánjù.} « Les vieux objets deviennent de nouveaux produits, comme de nouvelles bouteilles, du nouveau papier et de nouveaux jouets. »\\[0.3em]
\textbf{5)} Mot-clé implicite : \zh{好处} (avantages). Phrase du texte : \zh{回收利用不仅可以节约资源，还可以减少污染。}\\
Réponse : \zh{回收利用不仅可以节约资源，还可以减少污染，让环境更干净、更美丽。} \emph{Huíshōu lìyòng bùjǐn kěyǐ jiéyuē zīyuán, hái kěyǐ jiǎnshǎo wūrǎn, ràng huánjìng gèng gānjìng, gèng měilì.} « Le recyclage permet d'économiser les ressources, de réduire la pollution et de rendre l'environnement plus propre et plus beau. »
\end{exoresolu}

\begin{exoresolu}[Vrai ou Faux : justifiez en citant le texte]
Énoncé : dites si chaque affirmation est vraie (对) ou fausse (不对), puis citez la phrase du texte qui le prouve.\\
a) \zh{回收站里只有塑料。} « Il n'y a que du plastique au centre de recyclage. »\\
b) \zh{回收利用可以节约资源。} « Le recyclage peut économiser les ressources. »
\tcblower
\textbf{a) Faux (不对).} Le texte dit : \zh{塑料、纸和玻璃被分类、清洗。} \emph{Sùliào, zhǐ hé bōli bèi fēnlèi, qīngxǐ.} « Le plastique, le papier et le verre sont triés et nettoyés. » Il n'y a donc pas \emph{que} du plastique : il y a aussi du papier et du verre. Le mot \zh{只} (seulement) rend l'affirmation fausse.\\[0.3em]
\textbf{b) Vrai (对).} Le texte dit : \zh{回收利用不仅可以节约资源，还可以减少污染。} \emph{Huíshōu lìyòng bùjǐn kěyǐ jiéyuē zīyuán, hái kěyǐ jiǎnshǎo wūrǎn.} « Le recyclage permet non seulement d'économiser les ressources, mais aussi de réduire la pollution. » L'affirmation reprend exactement la première partie de cette phrase.
\end{exoresolu}

\begin{attention}[Piège des questions Vrai/Faux]
Les affirmations fausses contiennent presque toujours un \textbf{petit mot modificateur} qui trahit le texte : \zh{只} (seulement), \zh{都} (tous), \zh{没有} (ne… pas), \zh{每天} (chaque jour). Avant de répondre, comparez mot à mot l'affirmation et la phrase du texte : un seul mot peut tout changer. Et exigez-vous de \textbf{citer la phrase justificative} : au Bac, une réponse sans justification perd des points.
\end{attention}

\courssub{Production guidée : décrire un processus (Expression écrite/orale)}

\begin{methode}[Rédiger un paragraphe descriptif de processus type « recyclage / fabrication »]
\begin{enumerate}
\item \textbf{Ordonner} les étapes logiques (collecte $\rightarrow$ transport $\rightarrow$ traitement $\rightarrow$ produit final).
\item \textbf{Utiliser la chaîne} \zh{先…然后…再…最后…} pour structurer le paragraphe.
\item \textbf{Choisir la voix} : \zh{把} (actif, focus sur l'acteur) ou \zh{被} (passif, focus sur l'objet/déchet).
\item \textbf{Compléter les verbes} : ajout de compléments de résultat (\zh{干净}, \zh{好}, \zh{成}) ou \zh{了}.
\item \textbf{Conclure} par l'intérêt écologique : \zh{节约资源}, \zh{减少污染}, \zh{保护环境}.
\end{enumerate}
\end{methode}

\begin{experience}[Jeu de rôle : « Le guide du centre de recyclage »]
\begin{itemize}
\item \textbf{Binôme A (Guide) :} Accueille une classe de CM2. Explique le parcours d'une bouteille en plastique (\zh{塑料瓶}) en 4 étapes en utilisant \zh{先/然后/再/最后} et le passif \zh{被}.
\item \textbf{Binôme B (Journaliste) :} Pose 3 questions : \zh{第一步是什么？} / \zh{塑料瓶被谁清洗？} / \zh{最后变成什么？}.
\item \textbf{Échange :} Inversez les rôles. Le professeur circule et note l'usage correct de \zh{被} + complément.
\end{itemize}
\end{experience}

\begin{exercice}[Expression écrite : Le tri à l'école]
Rédigez un court paragraphe (60--80 caractères) pour décrire le processus de tri des déchets dans votre établissement (ou à la maison). Utilisez \textbf{au moins 3 connecteurs de séquence} et \textbf{une phrase passive avec 被}.
\begin{itemize}
\item Mots utiles : \zh{校园} (campus), \zh{分类桶} (poubelles de tri), \zh{纸张} (papier), \zh{回收箱} (bac de recyclage), \zh{制成} (fabriqué en).
\end{itemize}
\end{exercice}

\begin{corrige}[Exercice : Le tri à l'école]
\textbf{Modèle de réponse :}\\
\zh{为了保护校园环境，我们先把垃圾扔进不同的分类桶。然后清洁工把垃圾送到回收站。在回收站，纸张被分类、清洗，再制成新纸。最后，旧纸变成新本子，既节约资源，又减少污染。}\\
\emph{Wèile bǎohù xiàoyuán huánjìng, wǒmen xiān bǎ lājī rēng jìn bùtóng de fēnléi tǒng. Ránhòu qīngjiégōng bǎ lājī sòng dào huíshōuzhàn. Zài huíshōuzhàn, zhǐzhāng bèi fēnlèi, qīngxǐ, zài zhì chéng xīn zhǐ. Zuìhòu, jiù zhǐ biàn chéng xīn běnzi, jì jiéyuē zīyuán, yòu jiǎnshǎo wūrǎn.}
\end{corrige}

\begin{savaistu}[Le tri des déchets : Cameroun vs Chine]
\begin{itemize}
\item \textbf{Chine :} Depuis 2019, Shanghai impose un tri obligatoire en 4 catégories (\zh{干垃圾} secs, \zh{湿垃圾} humides, \zh{可回收物} recyclables, \zh{有害垃圾} dangereux). Amendes pour non-respect. Taux de recyclage en forte hausse.
\item \textbf{Cameroun :} Initiatives locales (Yaoundé, Douala) avec ONG et jeunes (\zh{Éco-citoyens}). Tri informel par les \zh{chiffonniers} (récupérateurs) pour le plastique/métal. Défi : absence de collecte sélective généralisée et de centres de traitement industriels.
\item \textbf{Point commun :} La prise de conscience grandit. L'éducation (écoles, médias) est la clé pour changer les habitudes, comme l'illustre notre texte.
\end{itemize}

\end{savaistu}

\begin{aretenir}[L'essentiel de cette section]
\begin{itemize}
\item La chaîne \zh{先…然后…再…最后…} ordonne les étapes d'un processus ; elle structure les textes descriptifs au Bac.
\item Le passif : \textbf{Objet + 被 + (agent) + Verbe + Complément} ; le verbe ne reste jamais seul (\zh{被收集起来}, \zh{被做成}, \zh{被清洗干净}).
\item Pour répondre à une question : repérer le mot interrogatif (\zh{哪里}, \zh{什么}…), localiser la phrase du texte, reformuler en citant.
\item Au Vrai/Faux, surveiller \zh{只}, \zh{都}, \zh{没有} et toujours justifier par une citation.
\item Production : utiliser la chaîne de séquence + \zh{被} + compléments pour décrire n'importe quel processus (recyclage, cuisine, fabrication).
\end{itemize}
\end{aretenir}


\coursec{Collocations et réemploi du lexique}

Après avoir analysé en détail le texte sur le processus du recyclage, nous allons maintenant \textbf{ancrer le vocabulaire} dans des combinaisons figées (collocations), construire des phrases correctes, pratiquer un dialogue de la vie courante et explorer les familles de mots. C'est l'étape indispensable pour passer de la \emph{compréhension passive} à l'\emph{expression active}, attendue à l'examen (épreuve écrite et orale).

\courssub{Les collocations incontournables du thème}

\begin{definition}[Collocation (\cle{搭配} dāpèi)]
Une collocation est une association \textbf{habituelle et naturelle} de deux mots ou plus. En chinois comme en français, on ne dit pas n'importe quel verbe avec n'importe quel nom : on \emph{trie} les déchets (pas *\emph{fait} les déchets), on \emph{protège} l'environnement (pas *\emph{garde} l'environnement). Maîtriser ces blocs tout faits évite les erreurs de traduction mot à mot.
\end{definition}

Observez d'abord ces six collocations extraites du texte et de la vie quotidienne. Elles suivent toutes la structure \textbf{Verbe + Complément d'objet} (动宾结构 dòngbīn jiégòu).

\begin{center}
\begin{tabularx}{\linewidth}{|X|X|X|X|}
\hline
\rowcolor{popblueL}
\textbf{Caractères} & \textbf{Pinyin} & \textbf{Nature} & \textbf{Traduction} \\
\hline
\cle{回收利用} & huíshōu lìyòng & Verbe composé & recycler \emph{et} réutiliser / valoriser \\
\hline
\cle{保护环境} & bǎohù huánjìng & Locution verbale & protéger l'environnement \\
\hline
\cle{节约资源} & jiéyuē zīyuán & Locution verbale & économiser les ressources \\
\hline
\cle{垃圾分类} & lājī fēnlèi & Nom verbal & tri des déchets / trier les déchets \\
\hline
\cle{养成好习惯} & yǎngchéng hǎo xíguàn & Locution verbale & prendre / acquérir de bonnes habitudes \\
\hline
\cle{变成} & biànchéng & Verbe & devenir, se transformer en \\
\hline
\end{tabularx}
\end{center}

\begin{exemplebox}[Exemples en contexte]
\begin{itemize}
\item \zh{旧塑料瓶可以回收利用。} \emph{Jiù sùliào píng kěyǐ huíshōu lìyòng.} Les vieilles bouteilles en plastique peuvent être recyclées et réutilisées.
\item \zh{保护环境是我们大家的责任。} \emph{Bǎohù huánjìng shì wǒmen dàjiā de zérèn.} Protéger l'environnement est la responsabilité de nous tous.
\item \zh{我们要节约水和电，节约资源。} \emph{Wǒmen yào jiéyuē shuǐ hé diàn, jiéyuē zīyuán.} Nous devons économiser l'eau et l'électricité, économiser les ressources.
\item \zh{请做好垃圾分类，把塑料放进蓝色垃圾桶。} \emph{Qǐng zuòhǎo lājī fēnlèi, bǎ sùliào fàng jìn lánsè lājītǒng.} Faites bien le tri des déchets, mettez le plastique dans la poubelle bleue.
\item \zh{从小养成好习惯，长大就不费力。} \emph{Cóng xiǎo yǎngchéng hǎo xíguàn, zhǎngdà jiù bù fèilì.} Acquérir de bonnes habitudes dès l'enfance rend les choses faciles une fois adulte.
\item \zh{废纸经过加工可以变成新纸。} \emph{Fèizhǐ jīngguò jiāgōng kěyǐ biànchéng xīn zhǐ.} Le papier usagé, après traitement, peut se transformer en nouveau papier.
\end{itemize}
\end{exemplebox}

\begin{propriete}[Structure \cle{变成} biànchéng — Transformation complète]
\zh{变成} indique un \textbf{changement de nature, de forme ou d'identité} : A \emph{devient} B (A n'est plus A).
\begin{center}
\begin{tabularx}{\linewidth}{|X|X|X|}
\hline
\rowcolor{popgreenL}
\textbf{Sujet (A)} & \textbf{变成} & \textbf{Objet résultant (B)} \\
\hline
旧瓶子 (vielles bouteilles) & 变成 & 新产品 (nouveaux produits) \\
\hline
厨余垃圾 (déchets de cuisine) & 变成 & 有机肥料 (engrais organique) \\
\hline
废金属 (métaux usagés) & 变成 & 原材料 (matières premières) \\
\hline
\end{tabularx}
\end{center}
\textbf{Attention :} \zh{变成} est un \textbf{verbe}. Ne le confondez pas avec le nom \zh{变化 biànhuà} « changement, variation ».
\end{propriete}

\begin{attention}[Piège fréquent : \zh{变成} vs \zh{变化}]
\begin{itemize}
\item \zh{变成} (verbe) : \zh{冰变成水了。} \emph{Bīng biànchéng shuǐ le.} La glace \textbf{est devenue} de l'eau. (Transformation achevée)
\item \zh{变化} (nom / verbe intransitif « changer ») : \zh{天气变化很快。} \emph{Tiānqì biànhuà hěn kuài.} La météo \textbf{change} vite. / \zh{发生了很大的变化。} \emph{Fāshēng le hěn dà de biànhuà.} Il y a eu de \textbf{grands changements}.
\end{itemize}
\emph{Règle mnémotechnique :} \zh{成 chéng} = « accomplir, devenir » (comme dans \zh{成功} réussir, \zh{成为} devenir). Si le résultat est un \textbf{nouvel objet}, c'est \zh{变成}.
\end{attention}

\courssub{Tableau récapitulatif : Verbe + Complément figé (blocs lexicaux)}

\begin{popfigure}
\begin{tikzpicture}[
    node distance=8mm and 12mm,
    block/.style={rectangle, draw=popblue, thick, rounded corners=4pt, fill=popblueL, text width=3.2cm, align=center, minimum height=1.6cm, font=\small},
    arrow/.style={-Stealth, thick, popdark},
    labelnode/.style={font=\footnotesize, text=popmuted, align=center}
]
% Nœuds principaux
\node[block, align=center] (v1){\textbf{保护 bǎohù}\\\emph{protéger} + \textbf{环境 huánjìng}\\\emph{environnement}};
\node[block, right=of v1, align=center] (v2){\textbf{节约 jiéyuē}\\\emph{économiser} + \textbf{资源 zīyuán}\\\emph{ressources}};
\node[block, right=of v2, align=center] (v3){\textbf{养成 yǎngchéng}\\\emph{acquérir} + \textbf{习惯 xíguàn}\\\emph{habitude}};
\node[block, below=of v1, align=center] (v4){\textbf{回收 huíshōu}\\\emph{récupérer} + \textbf{利用 lìyòng}\\\emph{utiliser}};
\node[block, right=of v4, align=center] (v5){\textbf{垃圾 lājī}\\\emph{déchet} + \textbf{分类 fēnlèi}\\\emph{classer / tri}};
\node[block, right=of v5, align=center] (v6){\textbf{减少 jiǎnshǎo}\\\emph{réduire} + \textbf{污染 wūrǎn}\\\emph{pollution}};

% Flèches de liaison (optionnel, pour montrer la cohérence thématique)
\draw[arrow, dashed] (v1) -- (v4);
\draw[arrow, dashed] (v2) -- (v5);
\draw[arrow, dashed] (v3) -- (v6);

% Légendes sous chaque bloc
\node[labelnode, below=2mm of v1] {Action citoyenne n°1};
\node[labelnode, below=2mm of v2] {Gestion des ressources};
\node[labelnode, below=2mm of v3] {Éducation / habitude};
\node[labelnode, below=2mm of v4] {Économie circulaire};
\node[labelnode, below=2mm of v5] {Geste quotidien};
\node[labelnode, below=2mm of v6] {Objectif final};
\end{tikzpicture}
\legende{Cartographie des collocations Verbe + Objet du thème « Recyclage \& Environnement ». Mémorisez-les par paires indissociables.}
\end{popfigure}

\courssub{Construction guidée : \textbf{Sujet + \cle{应该} yīnggāi + Verbe}}

\begin{propriete}[Le modal \cle{应该} yīnggāi — « devoir / avoir l'obligation morale »]
\zh{应该} exprime un \textbf{devoir, une recommandation forte, une évidence morale}. Il se place \textbf{avant le verbe principal}.
\medskip
\textbf{Schéma :} \quad \zh{Sujet + 应该 + Verbe (+ Complément)}
\medskip
\textbf{Négation :} \zh{不应该 bù yīnggāi} (ne pas devoir / il ne faut pas).
\medskip
\textbf{Question :} \zh{应该 ... 吗？} / \zh{要不要 ...？}
\end{propriete}

\begin{exemplebox}[Phrases modèles pour l'oral et l'écrit]
\begin{itemize}
\item \zh{我们应该保护环境。} \emph{Wǒmen yīnggāi bǎohù huánjìng.} Nous devons protéger l'environnement.
\item \zh{学生不应该乱扔垃圾。} \emph{Xuésheng bù yīnggāi luàn rēng lājī.} Les élèves ne doivent pas jeter les déchets n'importe où.
\item \zh{大家应该养成垃圾分类的习惯。} \emph{Dàjiā yīnggāi yǎngchéng lājī fēnlèi de xíguàn.} Tout le monde devrait prendre l'habitude du tri des déchets.
\item \zh{工厂应该减少污染。} \emph{Gōngchǎng yīnggāi jiǎnshǎo wūrǎn.} Les usines doivent réduire la pollution.
\item \zh{你应该把电池放进有害垃圾桶。} \emph{Nǐ yīnggāi bǎ diànchí fàng jìn yǒuhài lājītǒng.} Tu dois mettre les piles dans la poubelle pour déchets dangereux.
\end{itemize}
\end{exemplebox}

\begin{methode}[Transformer une idée en phrase avec \zh{应该}]
\begin{enumerate}
\item Identifier \textbf{qui} agit (Sujet : \zh{我}, \zh{我们}, \zh{大家}, \zh{工厂}, \zh{政府}...).
\item Choisir l'\textbf{action} attendue (Verbe + Objet : \zh{节约资源}, \zh{回收利用}...).
\item Insérer \zh{应该} entre le sujet et le verbe.
\item Vérifier l'ordre : \textbf{Sujet — 应该 — Verbe — Objet}. \emph{(Pas d'inversion comme en français « Devons-nous... »)}
\end{enumerate}
\end{methode}

\courssub{Mini-dialogue situationnel : Le tri à la maison / à l'école}

\begin{textebox}[Dialogue : Devant les poubelles de tri — Yaoundé, quartier Mvog-Ada]
\zh{— 你把垃圾放进哪个垃圾桶？} \\
\emph{Nǐ bǎ lājī fàng jìn nǎge lājītǒng ?} \\
— Tu mets les déchets dans quelle poubelle ?

\zh{— 我把塑料瓶放进蓝色的垃圾桶，把厨余垃圾放进绿色的垃圾桶。} \\
\emph{Wǒ bǎ sùliào píng fàng jìn lánsè de lājītǒng, bǎ chúyú lājī fàng jìn lǜsè de lājītǒng.} \\
— Je mets les bouteilles en plastique dans la poubelle bleue, et les déchets de cuisine dans la poubelle verte.

\zh{— 对！电池和旧药呢？} \\
\emph{Duì! Diànchí hé jiù yào ne?} \\
— Exact ! Et les piles et les vieux médicaments ?

\zh{— 它们是有害垃圾，要放进红色的垃圾桶。} \\
\emph{Tāmen shì yǒuhài lājī, yào fàng jìn hóngsè de lājītǒng.} \\
— Ce sont des déchets dangereux, il faut les mettre dans la poubelle rouge.

\zh{— 很好，我们都应该做好垃圾分类，保护环境。} \\
\emph{Hěn hǎo, wǒmen dōu yīnggāi zuòhǎo lājī fēnlèi, bǎohù huánjìng.} \\
— Très bien, nous devons tous bien faire le tri, protéger l'environnement.
\end{textebox}

\begin{aretenir}[Structure \cle{把} bǎ — Mise en focus sur l'objet manipulé]
Dans le dialogue : \zh{我把塑料瓶放进蓝色的垃圾桶}.
\zh{把} signale que l'objet (\zh{塑料瓶}) est \textbf{déplacé, transformé ou traité} par l'action. Ordre : \zh{Sujet + 把 + Objet + Verbe + Complément de résultat/direction}. C'est la structure standard pour dire « Je \textbf{mets} la bouteille \textbf{dans} la poubelle ».
\end{aretenir}

\begin{experience}[Jeu de rôle : « Le bon geste »]
\begin{enumerate}
\item En binôme : l'un tient des images (bouteille, épluchure, pile, papier, canette), l'autre a trois boîtes coloriées (bleu, vert, rouge).
\item A montre une image → B dit : \zh{这是……，应该放进……色的垃圾桶。} (C'est…, il faut le mettre dans la poubelle … couleur.)
\item Inverser les rôles. Chronométrer : 1 minute par binôme. Le plus grand nombre de phrases correctes gagne.
\end{enumerate}
\end{experience}

\courssub{Compléments utiles pour l'examen (niveau HSK 3-4)}

Deux expressions à très forte probabilité d'apparition à l'écrit (rédaction, traduction) et à l'oral (exposé) :

\begin{center}
\begin{tabularx}{\linewidth}{|X|X|X|X|}
\hline
\rowcolor{poporangeL}
\textbf{Caractères} & \textbf{Pinyin} & \textbf{Nature} & \textbf{Traduction} \\
\hline
\cle{减少污染} & jiǎnshǎo wūrǎn & Locution verbale & réduire la pollution \\
\hline
\cle{可回收利用的} & kě huíshōu lìyòng de & Adjectif (structure \zh{可 + V + de}) & recyclable / qui peut être recyclé \\
\hline
\end{tabularx}
\end{center}

\begin{exemplebox}[Emploi en phrase complète]
\begin{itemize}
\item \zh{垃圾分类可以减少污染。} \emph{Lājī fēnlèi kěyǐ jiǎnshǎo wūrǎn.} Le tri des déchets peut réduire la pollution.
\item \zh{请把可回收利用的纸张放进蓝色桶。} \emph{Qǐng bǎ kě huíshōu lìyòng de zhǐzhāng fàng jìn lánsè tǒng.} Mettez le papier recyclable dans la poubelle bleue.
\item \zh{这个塑料袋不可回收利用。} \emph{Zhège sùliào dài bù kě huíshōu lìyòng.} Ce sac en plastique n'est pas recyclable.
\end{itemize}
\end{exemplebox}

\begin{propriete}[Construction \cle{可} kě + Verbe + \cle{的} de = « -able / qui peut être … »]
Productif pour former des adjectifs à partir de verbes d'action :
\zh{可读} (lisible), \zh{可用} (utilisable), \zh{可回收} (recyclable), \zh{可持续} (durable — \zh{持续} soutenir/continuer).
Ajouter \zh{的} devant le nom : \zh{可回收利用的材料} (matériaux recyclables).
\end{propriete}

\courssub{Jeu des familles de mots : Étendre son lexique par les radicaux}

Travailler par \textbf{familles de caractères} (même clé / même composant phonétique) est la méthode la plus efficace pour mémoriser durablement.

\begin{center}
\begin{tabularx}{\linewidth}{|X|X|X|X|X|}
\hline
\rowcolor{poppurpleL}
\textbf{Clé / Composant} & \textbf{Caractère} & \textbf{Pinyin} & \textbf{Sens} & \textbf{Exemple de mot} \\
\hline
\zh{囗} (enceinte) & \zh{回} & huí & revenir, retourner & \zh{回收 huíshōu} (recycler), \zh{回来 huílái} (revenir), \zh{回答 huídá} (répondre), \zh{回家 huíjiā} (rentrer chez soi) \\
\hline
\zh{工} (travail, outil) & \zh{工} & gōng & travail, ouvrier & \zh{工人 gōngrén} (ouvrier), \zh{工厂 gōngchǎng} (usine), \zh{工作 gōngzuò} (travail/emploi), \zh{工程师 gōngchéngshī} (ingénieur) \\
\hline
\end{tabularx}
\end{center}

\begin{savaistu}[Le savais-tu ? — Le recyclage au Cameroun : Namé Recycling]
À \textbf{Douala}, l'entreprise \textbf{Namé Recycling} (créée en 2016) collecte les bouteilles en plastique (PET) et les transforme en \textbf{pavés écologiques} et en \textbf{granulats} pour l'industrie. C'est un exemple concret de \zh{变废为宝 biàn fèi wéi bǎo} « transformer le déchet en trésor ». En 2023, l'usine traite plusieurs tonnes par jour et emploie des dizaines de \zh{工人 gōngrén}. C'est un sujet idéal pour votre exposé oral : reliez le vocabulaire du cours (\zh{回收利用}, \zh{保护环境}, \zh{减少污染}) à une réalité locale vérifiable.
\end{savaistu}

\courssub{Texte d'entraînement intégré : « Une journée verte à l'école »}

\begin{textebox}[Texte original — Niveau HSK 3 / Terminale]
\zh{我们学校上周举办了“绿色校园”活动。星期一早上，校长在升旗仪式上说：“同学们，保护环境不仅是口号，更要行动。我们应该从垃圾分类做起，养成好习惯。”} \\
\emph{Wǒmen xuéxiào shàngzhōu jǔbàn le “Lǜsè xiàoyuán” huódòng. Xīngqīyī zǎoshang, xiàozhǎng zài shēngqí yíshì shàng shuō: “Tóngxuémen, bǎohù huánjìng bùjǐn shì kǒuhào, gèng yào xíngdòng. Wǒmen yīnggāi cóng lājī fēnlèi zuòqǐ, yǎngchéng hǎo xíguàn.”}

\zh{活动第一天，每个班级派两名同学当“环保小卫士”，站在教学楼下的垃圾桶旁，指导大家把垃圾放对地方。蓝色桶装可回收利用的纸、塑料、金属；绿色桶装厨余垃圾；红色桶装电池、旧药等有害垃圾。} \\
\emph{Huódòng dì yī tiān, měige bānjí pài liǎng míng tóngxué dāng “huánbǎo xiǎo wèishi”, zhàn zài jiàoxuélóu xià de lājītǒng páng, zhǐdǎo dàjiā bǎ lājī fàng duì dìfang. Lánsè tǒng zhuāng kě huíshōu lìyòng de zhǐ, sùliào, jīnshǔ; lǜsè tǒng zhuāng chúyú lājī; hóngsè tǒng zhuāng diànchí, jiù yào děng yǒuhài lājī.}

\zh{第三天，环保社团组织参观了城市郊区的回收利用工厂。工人们把废塑料清洗、粉碎、熔化，最后变成塑料颗粒，再做成新产品。同学们看得很仔细，纷纷说：“原来垃圾也能变成宝贝！以后我一定要节约资源，减少污染。”} \\
\emph{Dì sān tiān, huánbǎo shètuán zǔzhī cānguān le chéngshì jiāoqū de huíshōu lìyòng gōngchǎng. Gōngrénmen bǎ fèi sùliào qīngxǐ, fěnsuì, rónghuà, zuìhòu biànchéng sùliào kēlì, zài zuòchéng xīn chǎnpǐn. Tóngxuémen kàn de hěn zǐxì, fēnfēn shuō: “Yuánlái lājī yě néng biànchéng bǎobèi! Yǐhòu wǒ yídìng yào jiéyuē zīyuán, jiǎnshǎo wūrǎn.”}
\end{textebox}

\begin{center}
\begin{tabularx}{\linewidth}{|X|X|X|}
\hline
\rowcolor{popblueL}
\textbf{Vocabulaire clé du texte} & \textbf{Pinyin} & \textbf{Traduction} \\
\hline
举办 & jǔbàn & organiser (événement) \\
\hline
升旗仪式 & shēngqí yíshì & cérémonie de lever le drapeau \\
\hline
口号 & kǒuhào & slogan \\
\hline
行动 & xíngdòng & action, agir \\
\hline
派 & pài & envoyer, désigner \\
\hline
环保小卫士 & huánbǎo xiǎo wèishi & petit garde de l'environnement (élève bénévole) \\
\hline
指导 & zhǐdǎo & guider, encadrer \\
\hline
放对 & fàng duì & mettre au bon endroit \\
\hline
组织 & zǔzhī & organiser \\
\hline
参观 & cānguān & visiter \\
\hline
清洗 & qīngxǐ & laver, nettoyer \\
\hline
粉碎 & fěnsuì & broyer, pulvériser \\
\hline
熔化 & rónghuà & fondre \\
\hline
颗粒 & kēlì & granulé \\
\hline
纷纷 & fēnfēn & les uns après les autres, en foule \\
\hline
原来 & yuánlái & à l'origine / « ah bon ! » (surprise) \\
\hline
宝贝 & bǎobèi & trésor, chose précieuse \\
\hline
\end{tabularx}
\end{center}

\begin{exercice}[Questions de compréhension et réemploi]
\begin{enumerate}
\item D'après le texte, quelles sont les trois couleurs de poubelles et quel type de déchets va dans chacune ? Répondez en chinois (caractères + pinyin).
\item Relevez dans le texte \textbf{trois phrases} contenant la structure \zh{应该 + V} ou \zh{要 + V} (obligation / nécessité). Copiez-les.
\item Transformez les phrases suivantes en utilisant \zh{把} :
\begin{itemize}
\item 我扔垃圾。 → 我把垃圾扔进垃圾桶。
\item 他们回收塑料。 → \_\_\_\_\_\_\_\_\_\_
\item 我们保护环境。 → \_\_\_\_\_\_\_\_\_\_
\end{itemize}
\item Complétez avec \zh{可回收利用的} ou \zh{减少污染} :
\begin{itemize}
\item 垃圾分类可以 \_\_\_\_\_\_\_\_\_\_。
\item 请把 \_\_\_\_\_\_\_\_\_\_ 瓶子放进蓝色桶。
\end{itemize}
\item Expression orale (1 minute) : Présentez l'activité « École verte » de votre établissement (réel ou imaginé) en utilisant \textbf{au moins 5 collocations} de cette leçon.
\end{enumerate}
\end{exercice}

\begin{corrige}[Exercice 1]
\begin{enumerate}
\item \zh{蓝色桶：可回收利用的纸、塑料、金属 (lán sè tǒng: kě huíshōu lìyòng de zhǐ, sùliào, jīnshǔ) ; 绿色桶：厨余垃圾 (lǜ sè tǒng: chúyú lājī) ; 红色桶：电池、旧药等有害垃圾 (hóng sè tǒng: diànchí, jiù yào děng yǒuhài lājī).}
\item \zh{1. “我们应该从垃圾分类做起，养成好习惯。” 2. “以后我一定要节约资源，减少污染.” (含“要”表强烈意愿) 3. “它们是有害垃圾，要放进红色的垃圾桶.”}
\item \zh{他们把塑料回收利用。 (Tāmen bǎ sùliào huíshōu lìyòng.) ; 我们把环境保护好。 (Wǒmen bǎ huánjìng bǎohù hǎo.) — \emph{Note :} \zh{保护好} = protéger correctement (complément de résultat).}
\item \zh{1. 减少污染 (jiǎnshǎo wūrǎn) ; 2. 可回收利用的 (kě huíshōu lìyòng de).}
\item \emph{Critères d'évaluation oral :} fluidité, tons corrects, utilisation naturelle de \zh{回收利用 / 保护环境 / 节约资源 / 垃圾分类 / 养成好习惯 / 变成 / 减少污染 / 可回收利用的}, structure \zh{应该/把} correcte.
\end{enumerate}
\end{corrige}

\courssub{Synthèse visuelle : Du déchet au produit — Le cycle en 4 étapes}

\begin{popfigure}
\begin{tikzpicture}[
    node distance=10mm and 8mm,
    stage/.style={rectangle, rounded corners=6pt, draw=popgreen, thick, fill=popgreenL, text width=2.8cm, align=center, minimum height=2cm, font=\small},
    arrow/.style={-Stealth, thick, popgreen!70!popdark},
    label/.style={font=\footnotesize, text=popmuted, align=center}
]
\node[stage, align=center] (s1){\textbf{1. 分类 fēnlèi}\\\zh{垃圾分类}\\\emph{Tri à la source}};
\node[stage, right=of s1, align=center] (s2){\textbf{2. 收集 shōují}\\\zh{回收收集}\\\emph{Collecte sélective}};
\node[stage, right=of s2, align=center] (s3){\textbf{3. 加工 jiāgōng}\\\zh{清洗、粉碎、熔化}\\\emph{Traitement usine}};
\node[stage, right=of s3, align=center] (s4){\textbf{4. 重生 chóngshēng}\\\zh{变成新产品}\\\emph{Nouveaux objets}};

\draw[arrow] (s1) -- node[above, label] {把垃圾放对桶} (s2);
\draw[arrow] (s2) -- node[above, label] {送往工厂} (s3);
\draw[arrow] (s3) -- node[above, label] {变成颗粒/原料} (s4);

% Boucle de retour (économie circulaire)
\draw[arrow, dashed, bend left=40] (s4) to node[above, label] {再次使用 → 分类} (s1);
\end{tikzpicture}
\legende{Le cycle du recyclage : vocabulaire clé à chaque étape. La flèche en pointillés illustre l'économie circulaire (循环经济 xúnhuán jīngjì).}
\end{popfigure}

\begin{aretenir}[Check-list de révision — Section 5/6]
\begin{itemize}
\item \textbf{6 collocations de base} : \zh{回收利用, 保护环境, 节约资源, 垃圾分类, 养成好习惯, 变成} — savoir les écrire, les prononcer, les employer.
\item \textbf{2 compléments examen} : \zh{减少污染, 可回收利用的}.
\item \textbf{Structure \zh{应该 + V}} pour exprimer le devoir citoyen.
\item \textbf{Structure \zh{把 + O + V + Compl}} pour décrire le geste de tri.
\item \textbf{Familles de mots} : \zh{回} (retour) → \zh{回收/回来/回答} ; \zh{工} (travail) → \zh{工人/工厂/工作}.
\item \textbf{Distinction} : \zh{变成} (verbe, devenir) ≠ \zh{变化} (nom/verbe, changement).
\item \textbf{Contexte Cameroun} : Namé Recycling (Douala) = exemple vivant de \zh{变废为宝}.
\end{itemize}
\end{aretenir}

\coursec{Production guidée et évaluation}

Après avoir consolidé le lexique et les collocations du thème « processus du recyclage » dans la section précédente, nous passons maintenant à la mise en œuvre active : produire un texte cohérent, l'évaluer entre pairs, en discuter à l'oral et faire le bilan des acquis. Cette dernière étape transforme la compréhension passive en compétence communicative, conformément aux attendus du programme MINESEC (expression écrite, expression orale, auto-évaluation).

\courssub{Production écrite guidée : décrire le cycle du recyclage}

\begin{methode}[Plan d'écriture en 3 temps]
\textbf{Objectif :} rédiger un paragraphe de 4 à 5 phrases qui décrit le parcours d'une bouteille en plastique, du tri à la réutilisation, en utilisant la structure temporelle \cle{先…然后…再…最后…} et la banque de mots imposée.

\begin{enumerate}
\item \textbf{Lister les étapes en chinois} (mots-clés + verbes).\\
Exemple : 收集 (shōují) → 送到回收站 (sòng dào huíshōuzhàn) → 分类、清洗 (fēnlèi, qīngxǐ) → 送到工厂加工 (sòng dào gōngchǎng jiāgōng) → 变成新产品 (biànchéng xīn chǎnpǐn).
\item \textbf{Relier les étapes avec les connecteurs} \cle{先} (xiān), \cle{然后} (ránhòu), \cle{再} (zài), \cle{最后} (zuìhòu). Attention : \cle{然后} et \cle{再} s'utilisent tous deux pour « ensuite », mais \cle{然后} marque souvent une transition plus nette entre deux grandes phases, tandis que \cle{再} enchaîne des actions rapprochées.
\item \textbf{Conclure avec un bénéfice} : \cle{保护环境} (bǎohù huánjìng) ou \cle{节约资源} (jiéyuē zīyuán).
\end{enumerate}
\end{methode}

\begin{exemplebox}[Modèle attendu]
\begin{textebox}[Production modèle — Le cycle d'une bouteille en plastique]
我们先收集塑料瓶，然后把它们送到回收站。\\
在那里，瓶子被分类和清洗，再送到工厂加工。\\
最后，旧瓶子变成新产品，这样可以保护环境。
\end{textebox}
\textbf{Pinyin :} Wǒmen xiān shōují sùliàopíng, ránhòu bǎ tāmen sòng dào huíshōuzhàn. Zài nàlǐ, píngzi bèi fēnlèi hé qīngxǐ, zài sòng dào gōngchǎng jiāgōng. Zuìhòu, jiù píngzi biànchéng xīn chǎnpǐn, zhèyàng kěyǐ bǎohù huánjìng.\\
\textbf{Traduction :} Nous collectons d'abord les bouteilles en plastique, puis nous les envoyons au centre de recyclage. Là-bas, les bouteilles sont triées et lavées, puis envoyées à l'usine pour être transformées. Finalement, les vieilles bouteilles deviennent de nouveaux produits, ce qui permet de protéger l'environnement.
\end{exemplebox}

\begin{propriete}[Structure \cle{先…然后…再…最后…} — ordre chronologique]
\begin{center}
\begin{tabularx}{\linewidth}{|X|X|X|}
\hline
\rowcolor{popblueL} \textbf{Connecteur} & \textbf{Exemple (caractères + pinyin)} & \textbf{Traduction} \\ \hline
先 (xiān) & 我们先分类垃圾。 (Wǒmen xiān fēnlèi lājī.) & Nous trions d'abord les déchets. \\ \hline
然后 (ránhòu) & 然后把纸和塑料分开。 (Ránhòu bǎ zhǐ hé sùliào fēnkāi.) & Puis on sépare le papier et le plastique. \\ \hline
再 (zài) & 再送到工厂处理。 (Zài sòng dào gōngchǎng chǔlǐ.) & Ensuite on les envoie à l'usine pour traitement. \\ \hline
最后 (zuìhòu) & 最后变成新资源。 (Zuìhòu biànchéng xīn zīyuán.) & Finalement, ils deviennent de nouvelles ressources. \\ \hline
\end{tabularx}
\end{center}
\textbf{Remarque :} Le sujet peut être omis après le premier connecteur si le même agent effectue toutes les actions. La préposition \cle{把} (bǎ) est fréquente après \cle{然后} / \cle{再} quand l'objet est déplacé ou transformé. La voix passive avec \cle{被} (bèi) est naturelle quand l'objet devient sujet (ex. : \zh{瓶子被分类}).
\end{propriete}

\begin{exoresolu}[Écrire son paragraphe — Étape par étape]
\textbf{Consigne :} En utilisant la banque de mots imposée ci-dessous, écris 4 à 5 phrases qui décrivent le recyclage d'une bouteille en plastique. Respecte l'ordre \cle{先…然后…再…最后…} et termine par un bénéfice écologique.

\textbf{Banque de mots imposée (étendue) :}
\begin{center}
\begin{tabularx}{\linewidth}{|X|X|X|X|}
\hline
\rowcolor{popblueL} Caractères & Pinyin & Nature & Traduction \\ \hline
收集 & shōují & verbe & collecter \\ \hline
塑料瓶 & sùliàopíng & nom & bouteille en plastique \\ \hline
回收站 & huíshōuzhàn & nom & centre de recyclage \\ \hline
分类 & fēnlèi & verbe/nom & trier / tri \\ \hline
清洗 & qīngxǐ & verbe & laver, nettoyer \\ \hline
工厂 & gōngchǎng & nom & usine \\ \hline
加工 & jiāgōng & verbe & transformer, traiter \\ \hline
变成 & biànchéng & verbe & devenir, se transformer en \\ \hline
旧 & jiù & adj. & vieux, usagé \\ \hline
产品 & chǎnpǐn & nom & produit \\ \hline
保护环境 & bǎohù huánjìng & V-O & protéger l'environnement \\ \hline
\end{tabularx}
\end{center}

\textbf{Solution détaillée :}

\textit{Étape 1 — Liste des actions (verbes + compléments) :}
\begin{itemize}
\item 收集塑料瓶 (shōují sùliàopíng) — collecter les bouteilles en plastique
\item 送到回收站 (sòng dào huíshōuzhàn) — envoyer au centre de recyclage
\item 分类、清洗 (fēnlèi, qīngxǐ) — trier et laver
\item 送到工厂加工 (sòng dào gōngchǎng jiāgōng) — envoyer à l'usine pour transformation
\item 变成新产品 (biànchéng xīn chǎnpǐn) — devenir de nouveaux produits
\end{itemize}

\textit{Étape 2 — Insertion des connecteurs et structures grammaticales (把, 被) :}
\begin{zh}
我们先收集塑料瓶，然后把它们送到回收站。在那里，瓶子被分类和清洗，再送到工厂加工。最后，旧瓶子变成新产品，这样可以保护环境。
\end{zh}
\textit{Pinyin :} Wǒmen xiān shōují sùliàopíng, ránhòu bǎ tāmen sòng dào huíshōuzhàn. Zài nàlǐ, píngzi bèi fēnlèi hé qīngxǐ, zài sòng dào gōngchǎng jiāgōng. Zuìhòu, jiù píngzi biànchéng xīn chǎnpǐn, zhèyàng kěyǐ bǎohù huánjìng.

\textit{Étape 3 — Vérification :}
\begin{itemize}
\item[✓] 5 segments reliés par 先 / 然后 / 再 / 最后
\item[✓] 10 mots de la banque utilisés correctement
\item[✓] Structures \cle{把} (objet déplacé) et \cle{被} (passif) employées à bon escient
\item[✓] Caractères corrects, ponctuation chinoise pleine chasse (， 。)
\item[✓] Conclusion avec bénéfice écologique \cle{保护环境}
\end{itemize}
\end{exoresolu}

\begin{exercice}[À vous d'écrire — Variante : le recyclage du papier]
Rédigez un paragraphe similaire (4-5 phrases) pour décrire le recyclage du \cle{纸} (zhǐ, papier). Utilisez la même structure \cle{先…然后…再…最后…} et intégrez au moins 5 mots de la banque ci-dessous.
\begin{center}
\begin{tabularx}{\linewidth}{|X|X|X|X|}
\hline
\rowcolor{popblueL} Caractères & Pinyin & Nature & Traduction \\ \hline
收集 & shōují & verbe & collecter \\ \hline
废纸 & fèizhǐ & nom & papier usagé, déchets de papier \\ \hline
回收站 & huíshōuzhàn & nom & centre de recyclage \\ \hline
浆 & jiāng & nom & pâte (à papier) \\ \hline
制成 & zhìchéng & verbe & transformer en, fabriquer en \\ \hline
造纸 & zàozhǐ & verbe & fabriquer du papier \\ \hline
节约资源 & jiéyuē zīyuán & V-O & économiser les ressources \\ \hline
\end{tabularx}
\end{center}
\textit{Indication :} Pensez à l'ordre : collecter → envoyer au centre → faire de la pâte → fabriquer du nouveau papier → économiser les ressources. Utilisez \cle{被} pour le passif (le papier est transformé).
\end{exercice}

\courssub{Relecture par les pairs et grille d'évaluation}

\begin{methode}[Relecture croisée — Grille à 5 critères]
Chaque élève échange sa production avec un camarade. Le correcteur coche \textbf{OUI} ou \textbf{NON} et note un commentaire si nécessaire.

\begin{center}
\begin{tabularx}{\linewidth}{|X|X|X|X|}
\hline
\rowcolor{popblueL} \textbf{Critère} & \textbf{Indicateur observable} & \textbf{OUI} & \textbf{NON} \\ \hline
Vocabulaire imposé & Au moins 5 mots de la banque utilisés correctement & \square & \square \\ \hline
Ordre des étapes & Structure \cle{先…然后…再…最后…} respectée, logique chronologique & \square & \square \\ \hline
Structures grammaticales & \cle{把} et/ou \cle{被} utilisés correctement selon le contexte & \square & \square \\ \hline
Caractères & Écriture lisible, ordre des traits respecté pour 纸, 保, 环, 分, 类, 废, 浆 & \square & \square \\ \hline
Ponctuation & Ponctuation chinoise pleine chasse (， 。) ; pas d'espace entre caractères & \square & \square \\ \hline
Conclusion & Bénéfice écologique exprimé (保护环境 / 节约资源) & \square & \square \\ \hline
\end{tabularx}
\end{center}

\textbf{Conseil :} Si un critère est « NON », l'auteur corrige immédiatement avant remise au professeur. Cette relecture développe l'esprit critique et l'autonomie, compétences transversales du programme.
\end{methode}

\courssub{Questions orales de synthèse}

\begin{experience}[Débat guidé — Pourquoi et comment recycler ?]
\textbf{Durée :} 10 minutes (binômes, puis mise en commun).\\
\textbf{Consigne :} En binôme, répondez aux deux questions ci-dessous en chinois. Préparez 2-3 phrases par question. Le professeur passe écouter et note la prononciation des tons (surtout 回收 huíshōu, 分类 fēnlèi, 环境 huánjìng, 节约 jiéyuē).

\begin{enumerate}
\item \textbf{为什么回收很重要？} (Wèishéme huíshōu hěn zhòngyào?) — Pourquoi le recyclage est-il important ?
\item \textbf{我们应该怎么做？} (Wǒmen yīnggāi zěnme zuò?) — Que devrions-nous faire (au quotidien) ?
\end{enumerate}

\textbf{Pistes de réponse (vocabulaire d'appui) :}
\begin{itemize}
\item 回收可以减少垃圾，保护环境，节约资源。 (Huíshōu kěyǐ jiǎnshǎo lājī, bǎohù huánjìng, jiéyuē zīyuán.)
\item 我们应该分类垃圾：纸、塑料、金属、厨余垃圾分开放。 (Wǒmen yīnggāi fēnlèi lājī: zhǐ, sùliào, jīnshǔ, chúyú lājī fēnkāi fàng.)
\item 少用一次性塑料袋，多用布袋。 (Shǎo yòng yīcìxìng sùliàodài, duō yòng bùdài.)
\item 到回收站送旧电池、旧衣服、旧书。 (Dào huíshōuzhàn sòng jiù diànchí, jiù yīfu, jiù shū.)
\end{itemize}

\textbf{Mise en commun :} 3-4 binômes présentent leurs réponses au tableau. Le professeur corrige la prononciation (tons 3-3 → 2-3 pour 回收 huíshōu, 环境 huánjìng ; attention au ton neutre de 资源 zīyuán) et souligne les structures utiles : \cle{可以} (kěyǐ), \cle{应该} (yīnggāi), \cle{分开放} (fēnkāi fàng), \cle{少用…多用…} (shǎo yòng… duō yòng…).
\end{experience}

\courssub{Correction collective et affichage des meilleures productions}

\begin{aretenir}[Affichage et valorisation]
\begin{itemize}
\item Le professeur sélectionne 3 productions variées (papier, plastique, métal) et les affiche au mur de la classe ou sur le tableau numérique.
\item Chaque production affichée est annotée collectivement : \textbf{✓ connecteurs bien utilisés}, \textbf{✓ caractères soignés}, \textbf{✓ conclusion pertinente}, \textbf{✓ structures 把/被 correctes}.
\item Les élèves photographient l'affichage pour leur portfolio numérique OnBuch+ (preuve de compétence « Expression écrite »).
\item Cette trace écrite servira de base pour la révision du Module 2 (évaluation sommative).
\end{itemize}
\end{aretenir}

\courssub{Bilan et auto-évaluation}

\begin{aretenir}[Auto-évaluation des objectifs de la leçon ZH09]
Complète le tableau ci-dessous en toute honnêteté. Coche la case correspondante : \textbf{Acquis}, \textbf{En cours}, \textbf{À revoir}.

\begin{center}
\begin{tabularx}{\linewidth}{|X|X|X|X|X|}
\hline
\rowcolor{popblueL} \textbf{Objectif (compétence visée)} & \textbf{Acquis} & \textbf{En cours} & \textbf{À revoir} & \textbf{Action si En cours / À revoir} \\ \hline
Je reconnais et prononce le vocabulaire du recyclage (15 mots) & \square & \square & \square & Refaire les flashcards OnBuch+ \\ \hline
Je lis et comprends un texte de 200 mots sur le processus & \square & \square & \square & Relire le texte support + glossaire \\ \hline
J'écris un paragraphe avec \cle{先…然后…再…最后…} & \square & \square & \square & Réécrire l'exercice « Variante papier » \\ \hline
J'utilise correctement \cle{把} pour le déplacement/transformation & \square & \square & \square & Réviser la structure \cle{把} (leçon ZH07) \\ \hline
J'utilise correctement \cle{被} pour le passif & \square & \square & \square & Réviser la structure \cle{被} (leçon ZH08) \\ \hline
Je réponds oralement à \cle{为什么} et \cle{怎么做} & \square & \square & \square & S'enregistrer sur OnBuch+ (oral) \\ \hline
J'écris les caractères 纸, 保, 环, 分, 类, 废, 浆 dans l'ordre des traits & \square & \square & \square & S'entraîner sur cahier d'écriture (5× chaque) \\ \hline
\end{tabularx}
\end{center}

\textbf{Bilan personnel (1 phrase en français) :} \underline{\hspace{12cm}}
\end{aretenir}

\begin{popfigure}
\begin{tikzpicture}[
  box/.style={draw=popblue, thick, rounded corners=4pt, fill=popblueL, minimum width=2.8cm, minimum height=1.4cm, align=center, font=\small},
  arrow/.style={->, >=Stealth, thick, popdark},
  node distance=0.6cm and 0.4cm
]
% 5 boxes in a horizontal cycle
\node[box, align=center] (c1){收集\\(shōují)};
\node[box, right=of c1, align=center] (c2){\textbf{?}\\Transport};
\node[box, right=of c2, align=center] (c3){\textbf{?}\\Tri \& Lavage};
\node[box, right=of c3, align=center] (c4){\textbf{?}\\Usine};
\node[box, right=of c4, align=center] (c5){再利用\\(zàilìyòng)};

% Arrows
\draw[arrow] (c1.east) -- (c2.west);
\draw[arrow] (c2.east) -- (c3.west);
\draw[arrow] (c3.east) -- (c4.west);
\draw[arrow] (c4.east) -- (c5.west);
% Return arrow (curved)
\draw[arrow, bend left=40] (c5.west) to node[above, font=\tiny, text=popdark] {cycle} (c1.west);

% Labels for students to fill
\node[below=0.3cm of c2, font=\tiny, text=popmuted] {À compléter};
\node[below=0.3cm of c3, font=\tiny, text=popmuted] {À compléter};
\node[below=0.3cm of c4, font=\tiny, text=popmuted] {À compléter};
\end{tikzpicture}
\legende{Schéma du cycle de recyclage à compléter par les élèves (activité de révision ou évaluation formative). Les cases 2, 3, 4 attendent : 运送/送到回收站, 分类清洗, 工厂加工/变成原料.}
\end{popfigure}

\begin{attention}[Vérification finale des caractères — Pièges d'examen]
\begin{itemize}
\item \cle{纸} (zhǐ, papier) : 7 traits. Clé \cle{纟} (sī, soie) à gauche (3 traits : 丶, ㇆, ㇀) + \cle{氏} (shì) à droite (4 traits : 丶, 一, 丿, 乀). Ordre : clé d'abord, puis le reste haut→bas. \textbf{Piège :} ne pas confondre avec \cle{只} (zhǐ, seul/classificateur) qui a la clé \cle{口} (kǒu).
\item \cle{保} (bǎo, protéger) : 9 traits. Clé \cle{亻} (rén, homme) à gauche (2 traits : 丿, 丨) + \cle{呆} (dāi) à droite (6 traits). Ordre : 人 d'abord, puis 呆 : 口 (丨, ㇆, 一) → 木 (一, 丨, 丿, 乀). \textbf{Piège :} le trait final de 呆 est le \cle{乀} (nà, trait oblique droit) du 木, qui « ferme » le caractère ; ne pas le remplacer par un trait vertical.
\item \cle{环} (huán, anneau, environnement) : 8 traits. Structure gauche-droite : la clé \cle{王} (wáng, jade) à gauche, 4 traits, + la partie phonétique \cle{不} (bù) à droite, 4 traits. Ordre : on écrit d'abord la partie gauche \cle{王} (horizontal, horizontal, vertical, horizontal montant), puis la partie droite \cle{不}. \textbf{Piège :} la clé \cle{王} (jade) prend un trait horizontal montant à la fin quand elle est à gauche d'un caractère.
\item \cle{分} (fēn, diviser) : 4 traits. \cle{八} (bā, huit : 丿, 乀) + \cle{刀} (dāo, couteau : 丿, 乀). Ordre : 八 d'abord (haut→bas), puis 刀. \textbf{Piège :} ne pas écrire 刀 avant 八.
\item \cle{类} (lèi, catégorie) : 18 traits. Clé \cle{米} (mǐ, riz) en haut à gauche (6 traits) + \cle{大} (dà, grand : 3 traits) + \cle{页} (yè, page : 9 traits). Ordre : 米 → 大 → 页. \textbf{Piège :} oublier le trait horizontal final de 页 (le 9e trait de 页, le 18e du caractère) qui « ferme » le bas.
\item \cle{废} (fèi, déchet/usagé) : 8 traits. Clé \cle{广} (guǎng, toit) + \cle{己} (jǐ). Ordre : 广 (点, 横, 撇, 竖) → 己 (横, 撇, 折, 横). \textbf{Piège :} le 2e trait de 广 est un horizontal court, pas un point.
\item \cle{浆} (jiāng, pâte) : 10 traits. Clé \cle{氵} (shuǐ, eau) 3 traits + \cle{将} (jiāng) 7 traits. Ordre : 氵(点, 点, 提) → 将 (横, 竖, 撇, 点, 横, 横, 竖钩). \textbf{Piège :} le dernier trait de 将 est un \cle{竖钩} (shù gōu), pas un trait vertical simple.
\end{itemize}
\textbf{Conseil d'examen :} Un caractère mal tracé (ordre inversé, trait manquant, trait surnuméraire) peut être considéré comme incorrect au baccalauréat. Entraînez-vous sur quadrillage 田字格 (tiánzìgé) en respectant scrupuleusement l'ordre des traits standard (GB13000.1).
\end{attention}

\begin{savaistu}[Recyclage au Cameroun et en Chine — Regard croisé]
\begin{itemize}
\item \textbf{Cameroun :} Yaoundé et Douala ont mis en place des points de tri sélectif (papier, plastique, métal) gérés par des ONG comme \emph{Green Cameroon} ou l'\emph{Association pour la Protection de l'Environnement}. Le \cle{ndolé} et les déchets de cuisine (épluchures) sont souvent compostés dans les jardins familiaux — une forme de recyclage organique traditionnel. Des startups locales (ex. \emph{Kemit Ecology}) transforment les déchets plastiques en pavés écologiques.
\item \textbf{Chine :} Depuis 2019, Shanghai impose le tri obligatoire en 4 catégories : \cle{可回收物} (kě huíshōu wù, recyclables), \cle{有害垃圾} (yǒuhài lājī, déchets dangereux), \cle{厨余垃圾} (chúyú lājī, déchets de cuisine/humides), \cle{其他垃圾} (qítā lājī, autres/résidus). Des poubelles intelligentes à reconnaissance faciale récompensent le tri par des points échangeables. La Chine atteint un taux de recyclage des bouteilles PET supérieur à 90\%, leader mondial.
\item \textbf{Point commun :} Les deux pays développent l'économie circulaire (\cle{循环经济} xúnhuán jīngjì). Au Cameroun, le plastique devient matériau de construction ; en Chine, les \cle{无废城市} (wú fèi chéngshì, villes sans déchets) pilotes intègrent le recyclage dans l'urbanisme et l'industrie.
\end{itemize}
\end{savaistu}

\begin{corrige}[Exercice « Variante : le recyclage du papier »]
\textbf{Production attendue (exemple) :}
\begin{zh}
我们先收集废纸，然后送到回收站。\\
在那里，纸被分类和制成浆，再送到工厂造纸。\\
最后，旧纸变成新纸，这样可以节约资源。
\end{zh}
\textbf{Pinyin :} Wǒmen xiān shōují fèizhǐ, ránhòu sòng dào huíshōuzhàn. Zài nàlǐ, zhǐ bèi fēnlèi hé zhìchéng jiāng, zài sòng dào gōngchǎng zàozhǐ. Zuìhòu, jiù zhǐ biànchéng xīn zhǐ, zhèyàng kěyǐ jiéyuē zīyuán.\\
\textbf{Traduction :} Nous collectons d'abord le papier usagé, puis l'envoyons au centre de recyclage. Là-bas, le papier est trié et transformé en pâte, puis envoyé à l'usine pour fabriquer du papier. Finalement, le vieux papier devient du nouveau papier, ce qui permet d'économiser les ressources.

\textbf{Points de vigilance :}
\begin{itemize}
\item \cle{废纸} (fèizhǐ) = papier usagé / déchets de papier (vocabulaire HSK 4, très utile).
\item \cle{制成浆} (zhìchéng jiāng) = « être transformé en pâte » (passif par \cle{被} + verbe résultat).
\item \cle{造纸} (zàozhǐ) = fabriquer du papier (verbe composé V-O).
\item Structure \cle{先…然后…再…最后…} respectée, conclusion avec \cle{节约资源}.
\item Éviter la répétition \cle{再在} (zài zài) : préférer \cle{再送到} (zài sòng dào) ou \cle{在工厂} (zài gōngchǎng) seul.
\end{itemize}
\end{corrige}

\newpage

\coursec{Activité d'intégration}

\coursec{Leçon 9 — Vocabulaire et compréhension de texte : processus du recyclage}

\courssub{Mise en route et découverte du texte}
\noindent Le thème du recyclage (\zh{循环利用}) s'inscrit dans le Module 2 « Environnement et santé ». En Terminale A, vous devez être capables de lire un texte descriptif et explicatif sur un processus technique, d'en extraire le vocabulaire spécialisé et de réemployer les structures séquentielles pour décrire une procédure. Nous commençons par un texte support original, rédigé en chinois standard (niveau HSK 3-4), qui modélise le parcours d'un déchet ménager du tri à la renaissance.

\begin{textebox}[Texte support : Le parcours d'une bouteille en plastique — 从垃圾到资源]
\zh{塑料瓶的“新生”记} \\
\zh{每天，我们喝完饮料后，随手扔进垃圾桶的塑料瓶，究竟去哪儿了？如果分类正确，它将开启一段“变废为宝”的旅程。} \\
\zh{第一步是\cle{分类投放}。居民把空瓶子投放进\cle{蓝色}可回收物垃圾桶。这是\cle{源头分类}，最关键的一环。} \\
\zh{第二步是\cle{收集运输}。\cle{环卫车}定时收集垃圾，运送到\cle{处理厂}。途中不得混装混运。} \\
\zh{第三步是\cle{处理加工}。在工厂，瓶子被\cle{清洗}、\cle{粉碎}、\cle{熔融}，变成塑料\cle{颗粒}（\cle{原料}）。} \\
\zh{最后一步是\cle{循环利用}。塑料颗粒被制成新\cle{产品}，如纤维衣物、公园长椅、新瓶子，重新回到市场。} \\
\zh{如果分类错误，瓶子会被当作\cle{其他垃圾}，送往\cle{焚烧厂}或\cle{填埋场}，造成资源浪费和环境污染。所以，\cle{垃圾分类，从我做起}，保护环境，人人有责。} \\[0.5em]
\textit{Měitiān, wǒmen hē wán yǐnliào hòu, suíshǒu rēng jìn lājītǒng de sùliào píng, jiūjìng qù nǎr le? Rúguǒ fēnlèi zhèngquè, tā jiāng kāiqǐ yī duàn ``biànfèiwéibǎo'' de lǚchéng. -- Dì yī bù shì \textbf{fēnlèi tóufàng}. Jūmín bǎ kōng píngzi tóufàng jìn \textbf{lán sè} kě huíshōu wù lājītǒng. Zhè shì \textbf{yuántóu fēnlèi}, zuì guānjiàn de yī huán. -- Dì èr bù shì \textbf{shōují yùnshū}. \textbf{Huánwèi chē} dìngshí shōují lājī, yùnsòng dào \textbf{chǔlǐ chǎng}. Túzhōng bùdé hùnzhuāng hùnyùn. -- Dì sān bù shì \textbf{chǔlǐ jiāgōng}. Zài gōngchǎng, píngzi bèi \textbf{qīngxǐ}、\textbf{fěnsuì}、\textbf{róngróng}, biànchéng sùliào \textbf{kēlì} (\textbf{yuánliào}). -- Zuìhòu yī bù shì \textbf{xúnhuán lìyòng}. Sùliào kēlì bèi zhìchéng xīn \textbf{chǎnpǐn}, rú xiānwéi yīwù、gōngyuán chángyǐ、xīn píngzi, chóngxīn huí dào shìchǎng. -- Rúguǒ fēnlèi cuòwù, píngzi huì bèi dāngzuò \textbf{qítā lājī}, sòngwǎng \textbf{fénshāo chǎng} huò \textbf{tiánmái chǎng}, zàochéng zīyuán lànfèi hé huánjìng wūrǎn. Suǒyǐ, \textbf{lājī fēnlèi, cóng wǒ zuòqǐ}, bǎohù huánjìng, rénrén yǒu zé.}

\noindent \textbf{Traduction~:} La « nouvelle vie » de la bouteille en plastique. Chaque jour, où finissent les bouteilles en plastique que nous jetons négligemment après avoir bu une boisson ? Si le tri est correct, elles entament un voyage de « transformation des déchets en trésors ». Étape 1 : Tri et dépôt. Les habitants jettent les bouteilles vides dans la poubelle bleue des recyclables. C'est le \textbf{tri à la source}, le maillon le plus crucial. Étape 2 : Collecte et transport. Les camions de voirie collectent les déchets à heures fixes et les acheminent vers l'\textbf{usine de traitement}. En route, pas de mélange. Étape 3 : Traitement et transformation. En usine, les bouteilles sont \textbf{lavées}, \textbf{broyées}, \textbf{fondues}, pour devenir des \textbf{granulés} de plastique (\textbf{matière première}). Dernière étape : \textbf{Recyclage en boucle}. Les granulés servent à fabriquer de nouveaux \textbf{produits} : vêtements en fibres, bancs publics, nouvelles bouteilles, de retour sur le marché. En cas d'erreur de tri, la bouteille finit en \textbf{autres déchets}, envoyée à l'\textbf{incinérateur} ou à la \textbf{décharge}, causant gaspillage et pollution. Donc : \textbf{Tri des déchets, ça commence par moi}, protéger l'environnement est le devoir de tous.
\end{textebox}

\noindent \textbf{Observation guidée~:} Repérez dans le texte les marqueurs chronologiques (\zh{第一步}, \zh{第二步}, \zh{最后一步}), la structure \cle{把} (\zh{居民把空瓶子投放进\ldots}, \zh{瓶子被清洗\ldots}), le vocabulaire technique (\zh{颗粒}, \zh{熔融}, \zh{焚烧厂}) et l'expression d'engagement \cle{从\ldots 做起}.

\begin{savaistu}[Contexte socioculturel : Chine vs Cameroun]
En Chine, la politique de \zh{垃圾分类} (tri obligatoire des déchets) est devenue loi à Shanghai en juillet 2019, puis étendue à 46 villes clés. Quatre catégories standardisées : \zh{可回收物} (bleu), \zh{有害垃圾} (rouge), \zh{厨余垃圾} (vert), \zh{其他垃圾} (gris/noir). Des amendes sanctionnent les contrevenants (individus : 200 CNY max ; entreprises : 50 000 CNY max). L'objectif : « Zéro déchet direct en décharge » (\zh{零填埋}).

Au Cameroun, la gestion des déchets est confiée à des sociétés comme \textbf{HYSACAM} (Hygiène et Salubrité du Cameroun) pour Douala et Yaoundé. Le tri à la source reste embryonnaire, mais des initiatives scolaires (ex. clubs environnement, poubelles sélectives dans certains lycées comme Nkolbisson, Bilingue de Buea) et le projet « Yaoundé Ville Propre » visent à changer les comportements. Le défi majeur : la filière industrielle de recyclage (\zh{循环利用产业链}) est peu développée ; beaucoup de « recyclables » finissent encore en décharge (\zh{填埋场}) ou sont brûlés à l'air libre.
\end{savaistu}

\courssub{Glossaire du thème}
\noindent Voici le lexique essentiel de la leçon, classé par catégories grammaticales et thématiques. Maîtrisez la prononciation (tons), l'écriture (ordre des traits) et l'emploi contextuel.

\begin{definition}[Vocabulaire complet — Unité 9 : Environnement \& Recyclage]
\begin{center}
\begin{tabularx}{\linewidth}{|X|X|X|X|}
\hline
\rowcolor{popblueL} \textbf{Caractères} & \textbf{Pinyin} & \textbf{Nature} & \textbf{Traduction} \\ \hline
\zh{垃圾} & lājī & nom & Déchets, ordures \\ \hline
\zh{分类} & fēnlèi & verbe / nom & Trier, classer / classification \\ \hline
\zh{投放} & tóufàng & verbe & Jeter (dans une poubelle), déposer \\ \hline
\zh{收集} & shōují & verbe & Collecter, ramasser \\ \hline
\zh{运输} & yùnshū & verbe & Transporter, acheminer \\ \hline
\zh{处理} & chǔlǐ & verbe & Traiter (déchets, affaires) \\ \hline
\zh{循环利用} & xúnhuán lìyòng & verbe & Recycler, valoriser en boucle \\ \hline
\zh{变废为宝} & biànfèiwéibǎo & chengyu & Transformer les déchets en trésors (valorisation) \\ \hline
\zh{可回收物} & kě huíshōu wù & nom & Déchets recyclables (poubelle bleue) \\ \hline
\zh{有害垃圾} & yǒuhài lājī & nom & Déchets dangereux/toxiques (poubelle rouge) \\ \hline
\zh{厨余垃圾} & chúyú lājī & nom & Déchets de cuisine/organiques (poubelle verte) \\ \hline
\zh{其他垃圾} & qítā lājī & nom & Autres déchets / résidus (poubelle grise/noire) \\ \hline
\zh{源头分类} & yuántóu fēnlèi & nom & Tri à la source \\ \hline
\zh{环卫车} & huánwèi chē & nom & Camion poubelle / véhicule de voirie \\ \hline
\zh{处理厂} & chǔlǐ chǎng & nom & Usine de traitement \\ \hline
\zh{焚烧厂} & fénshāo chǎng & nom & Usine d'incinération \\ \hline
\zh{填埋场} & tiánmái chǎng & nom & Décharge (enfouissement) \\ \hline
\zh{清洗} & qīngxǐ & verbe & Laver, nettoyer \\ \hline
\zh{粉碎} & fěnsuì & verbe & Broyer, pulvériser \\ \hline
\zh{熔融} & róngróng & verbe & Fondre (métal, plastique) \\ \hline
\zh{颗粒} & kēlì & nom & Granulés, paillettes (matière première plastique) \\ \hline
\zh{原料} & yuánliào & nom & Matière première \\ \hline
\zh{产品} & chǎnpǐn & nom & Produit fini \\ \hline
\zh{流程} & liúchéng & nom & Processus, flux, procédure \\ \hline
\zh{阶段} & jiēduàn & nom & Étape, phase \\ \hline
\zh{关键} & guānjiàn & adj / nom & Crucial, clé / point clé \\ \hline
\zh{混装混运} & hùnzhuāng hùnyùn & verbe & Mélanger (collecte/transport non sélectif) \\ \hline
\zh{居民} & jūmín & nom & Habitants, résidents \\ \hline
\zh{随手} & suíshǒu & adv & Machinalement, sans réfléchir, « d'un geste de la main » \\ \hline
\zh{究竟} & jiūjìng & adv & Au juste,到底 (question rhétorique) \\ \hline
\zh{不得} & bùdé & verbe modal & Ne pas avoir le droit de (formel, règlement) \\ \hline
\zh{造成} & zàochéng & verbe & Causer, entraîner (conséquence négative) \\ \hline
\zh{浪费} & lànfèi & verbe / nom & Gaspiller / gaspillage \\ \hline
\zh{污染} & wūrǎn & verbe / nom & Polluer / pollution \\ \hline
\zh{保护} & bǎohù & verbe & Protéger \\ \hline
\zh{环境} & huánjìng & nom & Environnement \\ \hline
\zh{责任} & zérèn & nom & Responsabilité, devoir \\ \hline
\zh{行动} & xíngdòng & verbe / nom & Agir, passer à l'action / action \\ \hline
\end{tabularx}
\end{center}
\end{definition}

\courssub{Clés (radicaux) et ordre des traits}
\noindent L'analyse des radicaux (\zh{部首} \emph{bùshǒu}) aide à mémoriser le sens et à classer les caractères dans le dictionnaire. Respectez l'ordre des traits (\zh{笔顺} \emph{bǐshùn}) : horizontal avant vertical, gauche avant droite, haut avant bas, extérieur avant intérieur, trait central avant les ailes.

\begin{propriete}[Analyse de caractères clés du thème]
\begin{center}
\begin{tabularx}{\linewidth}{|X|X|X|X|X|}
\hline
\rowcolor{popgreenL} \textbf{Caractère} & \textbf{Radical (Clé)} & \textbf{Nb traits} & \textbf{Sens du radical} & \textbf{Aide-mémoire / Étymologie} \\ \hline
\zh{分} & \zh{刀} (dāo, couteau) & 4 & Séparer, couper & Un couteau \zh{刀} coupe en \zh{八} (deux parts) $\to$ diviser. \\ \hline
\zh{类} & \zh{米} (mǐ, riz) & 18 & Grains, catégories & Trier le bon grain (\zh{米}) de l'ivraie (\zh{页} tête/page) $\to$ catégorie. \\ \hline
\zh{回} & \zh{囗} (wéi, enclos) & 6 & Enceinte, retour & Revenir \zh{口} à l'intérieur de l'enceinte \zh{囗} $\to$ retourner. \\ \hline
\zh{害} & \zh{宀} (mián, toit) & 10 & Maison, abri & Sous le toit \zh{宀}, un \zh{口} et un \zh{戈} (arme) $\to$ danger intérieur. \\ \hline
\zh{厨} & \zh{厂} (chǎng, falaise/atelier) & 12 & Atelier, cuisine & Atelier \zh{厂} où l'on range (\zh{寸}) les ustensiles $\to$ cuisine. \\ \hline
\zh{运} & \zh{辶} (chuò, marche) & 12 & Mouvement, route & Marcher \zh{辶} avec une charge (\zh{军} armée) $\to$ transporter. \\ \hline
\zh{处} & \zh{夂} (zhǐ, pas lent) & 5 & Traiter, lieu & Un tigre \zh{虎} qui s'arrête (\zh{夂}) pour gérer sa proie $\to$ traiter, lieu. \\ \hline
\zh{循} & \zh{彳} (chì, pas court) & 12 & Suivre, circuler & Marcher \zh{彳} selon l'ordre (\zh{旬} cycle de 10 jours) $\to$ suivre un cycle. \\ \hline
\zh{变} & \zh{又} (yòu, main droite) & 8 & Changer, devenir & Deux mains \zh{又} \zh{又} qui retournent quelque chose $\to$ transformation. \\ \hline
\end{tabularx}
\end{center}
\end{propriete}

\begin{experience}[Atelier écriture : Ordre des traits sur ardoise]
\noindent L'enseignant projette l'animation de l'ordre des traits (ou dessine au tableau) pour \zh{分}, \zh{类}, \zh{回}, \zh{运}, \zh{处}, \zh{循}. Les élèves reproduisent sur ardoise magique ou papier quadrillé (style \zh{田字格}). Points de vigilance~:
\begin{itemize}
    \item \zh{类} : Écrire \zh{米} en premier (6 traits), puis \zh{页} (6 traits) : haut $\to$ bas.
    \item \zh{运} : Écrire \zh{云} (4 traits) \textbf{avant} le radical \zh{辶} (3 traits). Règle : composant principal avant radical enveloppant \zh{辶/走/廴}.
    \item \zh{循} : Radical \zh{彳} (3 traits) en premier, puis \zh{旬} (9 traits).
    \item \zh{处} : Haut \zh{虎} (6 traits) sans le trait final horizontal, puis \zh{夂} (3 traits).
\end{itemize}
\end{experience}

\courssub{Compréhension détaillée du texte}
\noindent Relisez le texte support « \zh{塑料瓶的“新生”记} » et répondez aux questions en français (compréhension) ou en chinois (expression guidée).

\begin{exercice}[Questions de compréhension]
\begin{enumerate}
    \item \textbf{Global~:} Quel est le sort final de la bouteille en plastique \textbf{si le tri est correct} ? Et \textbf{si le tri est incorrect} ? Citez les termes chinois correspondants.
    \item \textbf{Détaillé~:} Quelles sont les quatre étapes (\zh{第一步} à \zh{最后一步}) ? Résumez chacune en un verbe chinois + son objet (ex. \zh{分类投放 + 空瓶子}).
    \item \textbf{Lexique en contexte~:} Expliquez le sens de \zh{源头分类} et \zh{混装混运} à partir du texte. Pourquoi le texte dit-il \zh{途中不得混装混运} ?
    \item \textbf{Grammaire~:} Identifiez deux phrases à la voix passive avec \zh{被} (\emph{bèi}). Transformez-les en voix active avec \cle{把} (si possible) ou sujet agent.
    \item \textbf{Inférence~:} Selon le texte, pourquoi le \zh{源头分类} est-il le \zh{最关键的一环} (maillon le plus crucial) ? Justifiez par la chaîne logique du processus.
    \item \textbf{Expression~:} Réécrivez le dernier paragraphe (\zh{如果分类错误\ldots}) en utilisant la structure \cle{先\ldots 然后\ldots 最后\ldots} pour décrire le parcours négatif du déchet mal trié.
\end{enumerate}
\end{exercice}

\begin{corrige}[Exercice Compréhension]
\begin{enumerate}
    \item Correct : \zh{循环利用} $\to$ \zh{变成新产品} (ex. \zh{纤维衣物、公园长椅}). Incorrect : \zh{当作其他垃圾} $\to$ \zh{送往焚烧厂或填埋场}.
    \item 1. \zh{分类投放} (空瓶子 / 可回收物) ; 2. \zh{收集运输} (垃圾 / 到处理厂) ; 3. \zh{处理加工} (瓶子 $\to$ 清洗/粉碎/熔融 $\to$ 颗粒/原料) ; 4. \zh{循环利用} (颗粒 $\to$ 制成新产品).
    \item \zh{源头分类} = tri effectué par le producteur (ménage) dès la poubelle domestique. \zh{混装混运} = mélanger les catégories lors de la collecte/transport, annulant le tri. Interdit (\zh{不得}) car cela rend le recyclage industriel impossible ou trop coûteux.
    \item Passif : \zh{瓶子被清洗、粉碎、熔融} ; \zh{塑料颗粒被制成新产品}. Actif (把) : \zh{工人把瓶子清洗、粉碎、熔融} ; \zh{工厂把塑料颗粒制成新产品}.
    \item Si le tri à la source échoue, les recyclables sont souillés par les déchets humides (\zh{厨余}) ou dangereux (\zh{有害}), devenant impropres au recyclage (\zh{降解} qualité matière). Tout le processus aval (collecte, transport, usine) perd son efficacité.
    \item \zh{如果分类错误，垃圾先被当作其他垃圾，然后被运送到焚烧厂或填埋场，最后造成资源浪费和环境污染。}
\end{enumerate}
\end{corrige}

\courssub{Collocations et réemploi du lexique}
\noindent Le chinois fonctionne par « blocs » lexicaux (collocations). Mémorisez ces associations Verbe + Objet / Sujet + Verbe / Adjectif + Nom pour gagner en fluidité et justesse.

\begin{propriete}[Collocations verbales du processus de recyclage]
\begin{center}
\begin{tabularx}{\linewidth}{|X|X|X|}
\hline
\rowcolor{poporangeL} \textbf{Structure} & \textbf{Collocations types (Caractères + Pinyin)} & \textbf{Traduction} \\ \hline
Verbe + Objet (Action citoyenne) & \zh{分类垃圾} (fēnlèi lājī), \zh{投放垃圾} (tóufàng lājī), \zh{把垃圾分好} (bǎ lājī fēn hǎo) & Trier les déchets, jeter les déchets, bien trier les déchets \\ \hline
Verbe + Objet (Action service public) & \zh{收集垃圾} (shōují lājī), \zh{运输垃圾} (yùnshū lājī), \zh{清运垃圾} (qīngyùn lājī) & Collecter, transporter, enlever les déchets \\ \hline
Verbe + Objet (Action industrielle) & \zh{处理垃圾} (chǔlǐ lājī), \zh{粉碎瓶子} (fěnsuì píngzi), \zh{熔融塑料} (róngróng sùliào), \zh{生产颗粒} (shēngchǎn kēlì) & Traiter, broyer, fondre, produire des granulés \\ \hline
Verbe + Complément de résultat & \zh{分类好} (fēnlèi hǎo), \zh{处理成} (chǔlǐ chéng), \zh{变成} (biànchéng), \zh{制成} (zhìchéng) & Bien trié, transformé en, devenu, fabriqué en \\ \hline
Sujet + Verbe (Acteurs) & \zh{居民分类} (jūmín fēnlèi), \zh{环卫车收运} (huánwèi chē shōuyùn), \zh{工厂处理} (gōngchǎng chǔlǐ) & Les habitants trient, le camion collecte/transporte, l'usine traite \\ \hline
Adjectif + Nom (Catégories) & \zh{可回收物} (kě huíshōu wù), \zh{有害垃圾} (yǒuhài lājī), \zh{厨余垃圾} (chúyú lājī), \zh{其他垃圾} (qítā lājī) & Recyclables, Dangereux, Organiques, Résidus \\ \hline
\end{tabularx}
\end{center}
\end{propriete}

\begin{propriete}[Structures grammaticales récurrentes du thème]
\textbf{1. Séquence chronologique : \cle{先\ldots 然后\ldots 最后\ldots}} \\
Schéma : \textbf{Sujet + 先 + V1 (+ Compl) , 然后 + V2 (+ Compl) , 最后 + V3 (+ Compl).} \\
\emph{Exemple~:} \zh{我们 先 把垃圾 分类投放 ， 然后 环卫车 收集运输 ， 最后 工厂 循环利用 。} \\
\emph{Variante passive~:} \zh{垃圾 先 被 投放 ， 然后 被 运输 ， 最后 变成 产品 。}

\textbf{2. Structure \cle{把} (Disposition / Manipulation d'objet connu)} \\
Schéma : \textbf{Sujet + 把 + Objet (défini) + Verbe + Complément (résultat / direction / lieu).} \\
\emph{Condition~:} L'objet est connu, spécifique ou générique défini. Le verbe \textbf{doit} avoir un complément (particule \zh{了}, \zh{好}, \zh{完}, \zh{成}, \zh{到\ldots}, \zh{在\ldots}). \\
\emph{Exemples~:} \zh{请把电池投放进红桶。} (Direction) ; \zh{我已经把垃圾分类好了。} (Résultat \zh{好}) ; \zh{工厂把塑料处理成颗粒了。} (Résultat \zh{成}).

\textbf{3. Engagement personnel : \cle{从\ldots 做起} / \cle{从\ldots 开始}} \\
\zh{从我做起} (C'est parti de moi / Ça commence par moi) — Slogan civique figé. \\
\zh{从现在开始} (À partir de maintenant). \\
\zh{从小事做起} (En commençant par les petites choses).
\end{propriete}

\begin{exemplebox}[Entraînement : Transformation de phrases]
\begin{enumerate}
    \item \textbf{Actif $\to$ \cle{把} :} \zh{居民投放空瓶子。} $\to$ \zh{居民\cle{把}空瓶子\cle{投放进蓝桶里}。} (Ajout complément lieu obligatoire).
    \item \textbf{Séquence~:} Reliez : \zh{清洗} -- \zh{粉碎} -- \zh{熔融} -- \zh{变成颗粒}. $\to$ \zh{工厂\cle{先}清洗瓶子，\cle{然后}粉碎，\cle{最后}熔融\cle{成}颗粒。}
    \item \textbf{Engagement~:} « Protégeons l'environnement en commençant par le tri. » $\to$ \zh{保护环境，\cle{从分类做起}。} (Ou \zh{从我做起}).
\end{enumerate}
\end{exemplebox}

\courssub{Production guidée et évaluation : Concours d'affiche « Campus Vert »}
\noindent Cette activité finale mobilise l'ensemble des acquis (vocabulaire, grammaire, caractères, oral, culture) dans une tâche actionnelle authentique.

\courssub{Objectifs de l'activité}
Cette activité d'intégration vise à mobiliser l'ensemble des acquis de la leçon 9 (vocabulaire du recyclage, structures temporelles \cle{先\ldots 然后\ldots 最后}, structure \cle{把} pour la manipulation d'objets, expression de l'obligation ou de l'initiative avec \cle{从\ldots 做起}) dans une tâche finale authentique, collaborative et valorisante. Elle s'inscrit dans l'approche actionnelle préconisée par le programme MINESEC : l'élève devient acteur de sa communication en produisant un artefact réel (une affiche) destiné à un public réel (la communauté du lycée), dans un contexte camerounais concret (la Semaine de l'Environnement au lycée).

\noindent Les objectifs spécifiques sont les suivants~:
\begin{itemize}
    \item \textbf{Compétence lexicale~:} réemployer avec exactitude le vocabulaire des quatre catégories de poubelles (\zh{可回收物}, \zh{有害垃圾}, \zh{厨余垃圾}, \zh{其他垃圾}), les verbes d'action (\zh{分类}, \zh{投放}, \zh{收集}, \zh{运输}, \zh{处理}, \zh{循环利用}) et les noms du processus (\zh{流程}, \zh{阶段}, \zh{原料}, \zh{产品}).
    \item \textbf{Compétence grammaticale~:} construire des phrases correctes en utilisant la structure séquentielle \cle{先\ldots 然后\ldots 最后\ldots} pour décrire le processus, la structure \cle{把\ldots 动词\ldots 结果/方向} pour indiquer le tri actif (\zh{把垃圾分类投放}), et l'expression \cle{从\ldots 做起} pour l'engagement personnel.
    \item \textbf{Compétence scripturale~:} écrire les caractères clés de l'affiche avec un ordre des traits correct, une proportion équilibrée et une lisibilité adaptée à l'affichage public (taille, espacement).
    \item \textbf{Compétence orale~:} présenter l'affiche en 1 minute, avec une prononciation standard (tons, intonation), une fluidité suffisante et un regard porté vers l'auditoire.
    \item \textbf{Compétence socioculturelle et civique~:} établir un parallèle entre la politique chinoise de \zh{垃圾分类} (obligatoire à Shanghai depuis 2019) et les initiatives camerounaises (ex. \emph{HYSACAM} à Douala/Yaoundé, projets de tri à la source dans les établissements scolaires), pour ancrer l'apprentissage dans la réalité locale.
\end{itemize}

\begin{experience}[Situation déclenchante : Avis officiel -- Semaine de l'Environnement]
\begin{textebox}[Avis officiel -- Lycée Bilingue de Nkolbisson (Yaoundé)]
\zh{关于举办“绿色校园，从我做起”环保周活动的通知} \\
\zh{全体师生：} \\
\zh{为响应喀麦隆环境部号召，提高全校师生环保意识，我校决定于下周一至周五（3月18日--22日）举办“绿色校园，从我做起”环保周活动。} \\
\zh{主要活动包括：} \\
\zh{1. 环保知识讲座（周一上午，大礼堂）；} \\
\zh{2. 校园清洁大扫除（周三下午，全校参与）；} \\
\zh{3. “垃圾分类，从我做起”海报设计大赛（全周，各班提交一份作品，周五评选）。} \\
\zh{海报要求：主题鲜明，内容科学，汉法双语或纯汉语，标注拼音。优胜作品将在教学楼主走廊长期展示。} \\
\zh{请各班班长周五中午前将海报交至学生会办公室。} \\
\zh{学生会  二〇二五年三月十二日} \\[0.5em]
\textit{Guānyú jǔbàn ``Lǜsè xiàoyuán, cóng wǒ zuòqǐ'' huánbǎo zhōu huódòng de tōngzhī -- Quántǐ shīshēng: Wèi xiǎngyìng Kāmàilóng Huánjìngbù hàozhào, tígāo quánxiào shīshēng huánbǎo yìshí, wǒxiào juédìng yú xiàzhōu yī zhì zhōuwǔ (3 yuè 18 rì--22 rì) jǔbàn ``Lǜsè xiàoyuán, cóng wǒ zuòqǐ'' huánbǎo zhōu huódòng. Zhǔyào huódòng bāokuò: 1. Huánbǎo zhīshì jiǎngzuò (zhōuyī shàngwǔ, dà lǐtáng); 2. Xiàoyuán qīngjiē dàsǎochú (zhōusān xiàwǔ, quánxiào cānyù); 3. ``Lājī fēnlèi, cóng wǒ zuòqǐ'' hǎibào shèjì dàsài (quánzhōu, gè bān tíjiāo yī fèn zuòpǐn, zhōuwǔ píngxuǎn). Hǎibào yāoqiú: Zhǔtí xiānmíng, nèiróng kēxué, Hàn-Fǎ shuāngyǔ huò chún Hànyǔ, biāozhù pīnyīn. Yōushèng zuòpǐn jiāng zài jiàoxué lóu zhǔ zǒuláng chángqí zhǎnshì. Qǐng gè bān bānzhǎng zhōuwǔ zhōngwǔ qián jiāng hǎibào jiāo zhì xuéshēnghuì bàngōngshì. Xuéshēnghuì Èrlíng'èrwǔ nián sān yuè shí'èr rì.}

\noindent \textbf{Traduction~:} Avis concernant l'organisation de la Semaine de l'Environnement « Campus Vert, ça commence par moi ». Chers enseignants et élèves, pour répondre à l'appel du Ministère camerounais de l'Environnement et renforcer la conscience écologique, notre lycée organise du lundi 18 au vendredi 22 mars la Semaine de l'Environnement. Activités principales~: 1. Conférence (lundi matin, grand amphithéâtre) ; 2. Grand nettoyage du campus (mercredi après-midi, tous participants) ; 3. Concours d'affiches « Tri des déchets, ça commence par moi » (toute la semaine, une affiche par classe, vote vendredi). Exigences~: thème clair, contenu scientifique, bilingue chinois-français ou chinois seul avec pinyin. Les affiches gagnantes seront exposées durablement dans le hall principal. Remise au bureau des élèves vendredi midi. Bureau des élèves, 12 mars 2025.
\end{textebox}

\noindent Votre classe de Terminale A (LV2 Chinois) participe au concours. Le professeur vous divise en groupes de 4 élèves. Chaque groupe doit concevoir \textbf{une seule affiche} respectant le cahier des charges ci-dessus. L'affiche sera évaluée sur la correction linguistique, la richesse du vocabulaire de la leçon, la clarté visuelle et l'originalité du slogan.
\end{experience}

\courssub{Consignes de réalisation}
\noindent \textbf{Étape 1 -- Analyse de la situation et répartition des rôles (10 min)}
\begin{enumerate}
    \item Lisez l'avis officiel ci-dessus. Identifiez les informations clés : date limite, lieu d'exposition, exigences linguistiques (chinois + pinyin).
    \item Au sein du groupe, désignez~: un \textbf{rédacteur principal} (écrit les caractères au propre), un \textbf{recherchiste lexical} (vérifie le vocabulaire, le pinyin, les tons dans le glossaire de la leçon), un \textbf{maquettiste} (organise l'espace, les couleurs, les icônes des poubelles), un \textbf{présentateur} (prépare l'oral de 1 minute).
\end{enumerate}

\noindent \textbf{Étape 2 -- Conception du contenu textuel (25 min)}
\begin{enumerate}
    \item \textbf{Slogan principal~:} Choisissez ou créez un slogan percutant de 8 à 12 caractères. Exemples~: \zh{垃圾分类，从我做起} / \zh{保护环境，人人有责} / \zh{变废为宝，绿色未来}. Écrivez-le en grand au centre/haut de l'affiche avec pinyin.
    \item \textbf{Les 4 catégories de poubelles~:} Dessinez ou collez 4 icônes de poubelles de couleurs différentes (Bleu, Rouge, Vert, Gris/Noir). Pour chacune, écrivez \textbf{en caractères + pinyin}~:
    \begin{itemize}
        \item Nom de la catégorie (ex. \zh{可回收物  kě huíshōu wù}).
        \item 3 exemples d'objets vus en classe (ex. \zh{塑料瓶  sùliào píng}, \zh{旧报纸  jiù bàozhǐ}, \zh{金属罐  jīnshǔ guàn}).
        \item La couleur de la poubelle (ex. \zh{蓝色  lán sè}).
    \end{itemize}
    \item \textbf{Le processus de recyclage (3 phrases)~:} Rédigez \textbf{exactement 3 phrases} reliées par \cle{先}, \cle{然后}, \cle{最后} pour décrire ce qui arrive aux déchets après le tri. Utilisez le vocabulaire de la leçon (\zh{投放}, \zh{收集}, \zh{运输}, \zh{处理}, \zh{循环利用}, \zh{变成}).
    \begin{itemize}
        \item Phrase 1 (Début) : \zh{先\ldots} (Action du citoyen / collecte).
        \item Phrase 2 (Milieu) : \zh{然后\ldots} (Transport / usine).
        \item Phrase 3 (Fin) : \zh{最后\ldots} (Transformation / nouveau produit).
    \end{itemize}
    \item \textbf{Signature~:} \zh{Terminale A -- Groupe [Numéro] -- Lycée de Nkolbisson}.
\end{enumerate}

\noindent \textbf{Étape 3 -- Mise en page et révision (15 min)}
\begin{enumerate}
    \item Le maquettiste assemble le tout sur une feuille A3 (ou A2) fournie. Vérifiez l'alignement, la lisibilité des caractères (taille $\ge$ 1.5 cm pour le corps de texte), l'absence de fautes de traits.
    \item Le recherchiste relit \textbf{chaque caractère} et \textbf{chaque pinyin} (attention aux tons 3 et 2, aux rétroflexes \zh{zh, ch, sh, r}, aux voyelles composées).
    \item Le présentateur répète son texte oral (voir Ressources).
\end{enumerate}

\noindent \textbf{Étape 4 -- Présentation orale et vote (20 min)}
\begin{enumerate}
    \item Chaque présentateur dispose de \textbf{1 minute chrono} pour présenter l'affiche en chinois (pas de lecture mot à mot, notes autorisées sur fiche bristol).
    \item Structure de l'oral suggérée~: Salutation $\to$ Annonce du slogan $\to$ Explication rapide des 4 poubelles $\to$ Lecture des 3 phrases du processus $\to$ Conclusion engageante ($\to$ Merci).
    \item La classe vote à bulletin secret sur 3 critères~: \textbf{Langue} (correction, vocabulaire), \textbf{Contenu} (clarté du processus), \textbf{Visuel} (esthétique, lisibilité). L'affiche gagnante est affichée dans le couloir.
\end{enumerate}

\courssub{Ressources à mobiliser}
\noindent Rappel synthétique des outils linguistiques vus dans la leçon 9 et indispensables pour réussir l'affiche.

\begin{propriete}[Structure séquentielle : \cle{先\ldots 然后\ldots 最后\ldots}]
Cette structure permet de décrire un processus chronologique en trois étapes minimum. L'ordre des mots est strict~: \textbf{Sujet + 先 + Verbe 1 (+ Complément) , 然后 + Verbe 2 (+ Complément) , 最后 + Verbe 3 (+ Complément)}.
\begin{center}
\begin{tabularx}{\linewidth}{|X|X|X|}
\hline
\rowcolor{popgreenL} \textbf{Structure} & \textbf{Exemple (Caractères + Pinyin)} & \textbf{Traduction} \\ \hline
Sujet + \cle{先} + V1 + O , \cle{然后} + V2 + O , \cle{最后} + V3 + O & \zh{我们 先 把垃圾 分类投放 ， 然后 车子 收集运输 ， 最后 工厂 循环利用 。} & Nous trions d'abord les déchets, ensuite le camion les collecte et les transporte, enfin l'usine les recycle. \\ \hline
Sujet + \cle{先} + V1 , \cle{然后} + V2 , \cle{最后} + V3 & \zh{塑料瓶 先 被 投放 ， 然后 被 运输 ， 最后 变成 新产品 。} & La bouteille est d'abord jetée, puis transportée, enfin transformée en nouveau produit. (Passif \cle{被}) \\ \hline
\end{tabularx}
\end{center}
\noindent \textbf{Attention~:} Ne pas inverser l'ordre. \cle{然后} ne peut pas commencer la phrase. La virgule chinoise \zh{，} sépare les propositions.
\end{propriete}

\begin{propriete}[Structure \cle{把} (Disposition) : Manipulation d'objet connu]
Pour dire que l'on \textbf{fait quelque chose à un objet précis} (le déchet) et qu'on indique le résultat ou la destination~: \textbf{Sujet + 把 + Objet + Verbe + Complément de résultat/direction/lieu}.
\begin{center}
\begin{tabularx}{\linewidth}{|X|X|X|}
\hline
\rowcolor{poporangeL} \textbf{Structure} & \textbf{Exemple} & \textbf{Traduction} \\ \hline
Sujet + \cle{把} + Objet + \cle{动词} + \cle{方向/结果} & \zh{请 把 塑料瓶 投放 进 蓝色 垃圾桶 。} & Jetez la bouteille en plastique dans la poubelle bleue. \\ \hline
Sujet + \cle{把} + Objet + \cle{动词} + \cle{好/完/成} & \zh{我 已经 把 垃圾 分类 好 了 。} & J'ai déjà fini de trier les déchets. \\ \hline
Sujet + \cle{把} + Objet + \cle{动词} + \cle{到} + Lieu & \zh{环卫车 把 垃圾 运输 到 处理厂 了 。} & Le camion a transporté les déchets à l'usine. \\ \hline
\end{tabularx}
\end{center}
\noindent \textbf{Condition~:} L'objet doit être connu/défini (le, ce, mon, un quantificateur + classifieur). On ne dit pas \zh{*把一个垃圾} mais \zh{把这袋垃圾} ou \zh{把垃圾} (générique connu).
\end{propriete}

\begin{propriete}[Expression de l'initiative : \cle{从\ldots 做起} / \cle{从\ldots 开始}]
\begin{itemize}
    \item \zh{从我做起} (Cóng wǒ zuòqǐ) : « En commençant par moi » / « Ça commence par moi ». Structure figée, très utilisée pour les slogans civiques.
    \item \zh{从现在开始} (Cóng xiànzài kāishǐ) : « À partir de maintenant ».
    \item \zh{从小事做起} (Cóng xiǎoshì zuòqǐ) : « En commençant par les petites choses ».
\end{itemize}
\end{propriete}

\begin{methode}[Présentation orale de 1 minute -- Structure type]
\begin{enumerate}
    \item \textbf{Salutation formelle} (5 s) : \zh{各位同学，老师们，大家好。} (Gèwèi tóngxué, lǎoshīmen, dàjiā hǎo.)
    \item \textbf{Accroche + Slogan} (10 s) : \zh{我们组的海报主题是：「垃圾分类，从我做起」。} (Wǒmen zǔ de hǎibào zhǔtí shì: ``Lājī fēnlèi, cóng wǒ zuòqǐ''.)
    \item \textbf{Les 4 poubelles} (25 s) : Pointez l'affiche. \zh{这是四个垃圾桶：蓝色是可回收物，比如塑料瓶、纸；红色是有害垃圾，比如电池；绿色是厨余垃圾，比如剩菜；灰色是其他垃圾。} (Zhè shì sì gè lājītǒng: lán sè shì kě huíshōu wù, bǐrú sùliào píng, zhǐ; hóng sè shì yǒuhài lājī, bǐrú diànchí; lǜ sè shì chúyú lājī, bǐrú shèngcài; huī sè shì qítā lājī.)
    \item \textbf{Le processus (les 3 phrases)} (15 s) : Lisez vos 3 phrases \cle{先/然后/最后} distinctement.
    \item \textbf{Conclusion + Appel à l'action} (5 s) : \zh{保护环境，人人有责。请大家行动起来！谢谢。} (Bǎohù huánjìng, rénrén yǒu zé. Qǐng dàjiā xíngdòng qǐlái! Xièxie.)
\end{enumerate}
\end{methode}

\begin{attention}[Erreurs fréquentes à éviter pour l'affiche]
\begin{itemize}
    \item \textbf{Tons~:} \zh{垃圾} (lājī, 1-1) pas *lǎjǐ ; \zh{分类} (fēnlèi, 1-4) pas *fēnlěi ; \zh{处理} (chǔlǐ, 3-3 $\to$ 2-3 par sandhi) ; \zh{循环} (xúnhuán, 2-2) ; \zh{厨余} (chúyú, 2-2).
    \item \textbf{Classificateurs~:} Un déchet = \zh{一\cle{件} 垃圾} (pièce) ou \zh{一\cle{袋} 垃圾} (sac) ou \zh{一\cle{堆} 垃圾} (tas). Pas de \zh{个} pour \zh{垃圾} (masse). Une bouteille = \zh{一个\cle{个} 瓶子} ou \zh{一\cle{只} 瓶子}.
    \item \textbf{Ordre des mots (Processus)~:} Ne pas écrire \zh{最后 先 然后}. Respecter la chronologie logique.
    \item \textbf{\cle{把} sans complément~:} \zh{*我把垃圾分类。} (Incomplet) $\to$ \zh{我把垃圾分类\cle{好了} / \cle{投放到蓝桶里了}。}
    \item \textbf{Caractères proches~:} \zh{运} (transporter, radical \cle{辶}) $\neq$ \zh{过} (passer, radical \cle{辶}) ; \zh{处} (traiter, radical \cle{夂}) $\neq$ \zh{备} (préparer, radical \cle{亻}) ; \zh{粉} (poudre, radical \cle{米}) $\neq$ \zh{分} (diviser, radical \cle{刀}).
    \item \textbf{Pinyin sur l'affiche~:} Écrire le pinyin \textbf{sous} chaque mot chinois (pas syllabe par syllabe), en plus petit, aligné à gauche du mot. Ex. \zh{可回收物} \textit{kě huíshōu wù}.
\end{itemize}
\end{attention}

\courssub{Proposition de production et corrigé commenté}
\begin{exoresolu}[Affiche modèle -- Groupe 3 (Terminale A)]
\textbf{Énoncé~:} Production de l'affiche finale (contenu textuel) avec commentaires linguistiques.

\tcblower

\textbf{1. En-tête / Slogan principal (Grand format, centré)}
\begin{center}
\zh{\textbf{垃圾分类，从我做起}} \\[0.2em]
\textit{Lājī fēnlèi, cóng wǒ zuòqǐ} \\
« Tri des déchets, ça commence par moi »
\end{center}

\textbf{2. Zone centrale -- Les 4 poubelles (Tableau visuel 2x2)}

\begin{center}
\begin{tabular}{|c|c|}
\hline
\shortstack{\zh{\textbf{蓝色桶 \quad 可回收物}}\\ \zh{\textit{Lán sè tǒng \quad kě huíshōu wù}}} & \shortstack{\zh{\textbf{红色桶 \quad 有害垃圾}}\\ \zh{\textit{Hóng sè tǒng \quad yǒuhài lājī}}} \\
\zh{塑料瓶 (sùliào píng)} & \zh{废电池 (fèi diànchí)} \\
\zh{废纸/报纸 (fèi zhǐ / bàozhǐ)} & \zh{过期药品 (guòqī yàopǐn)} \\
\zh{金属罐 (jīnshǔ guàn)} & \zh{废灯管 (fèi dēngguǎn)} \\
\hline
\shortstack{\zh{\textbf{绿色桶 \quad 厨余垃圾}}\\ \zh{\textit{Lǜ sè tǒng \quad chúyú lājī}}} & \shortstack{\zh{\textbf{灰色桶 \quad 其他垃圾}}\\ \zh{\textit{Huī sè tǒng \quad qítā lājī}}} \\
\zh{剩饭菜 (shèng fàncài)} & \zh{卫生纸 (wèishēng zhǐ)} \\
\zh{果皮果核 (guǒpí guǒhé)} & \zh{烟蒂 (yāndī)} \\
\zh{菜叶 (càiyè)} & \zh{碎陶瓷 (suì táocí)} \\
\hline
\end{tabular}
\end{center}

\textbf{3. Zone basse -- Le processus de recyclage (3 phrases \cle{先/然后/最后})}
\begin{enumerate}
    \item \zh{\textbf{第一步：} 我们 \cle{先} \cle{把} 生活垃圾 \cle{分类投放} 到 对应 的 垃圾桶 里 。} \\
    \textit{Wǒmen \textbf{xiān} \textbf{bǎ} shēnghuó lājī \textbf{fēnlèi tóufàng} dào duìyìng de lājītǒng lǐ.} \\
    \textit{Étape 1 : Nous trions d'abord les déchets ménagers et les jeterons dans les poubelles correspondantes.}
    \item \zh{\textbf{第二步：} 环卫车 \cle{然后} \cle{把} 分类好 的 垃圾 \cle{收集运输} 到 处理厂 。} \\
    \textit{Huánwèi chē \textbf{ránhòu} \textbf{bǎ} fēnlèi hǎo de lājī \textbf{shōují yùnshū} dào chǔlǐ chǎng.} \\
    \textit{Étape 2 : Ensuite, le camion de voirie collecte et transporte les déchets bien triés vers l'usine de traitement.}
    \item \zh{\textbf{第三步：} 工厂 \cle{最后} \cle{把} 可回收物 \cle{处理} \cle{成} 新 原料 ， \cle{循环利用} \cle{变成} 新 产品 。} \\
    \textit{Gōngchǎng \textbf{zuìhòu} \textbf{bǎ} kě huíshōu wù \textbf{chǔlǐ} \textbf{chéng} xīn yuánliào, \textbf{xúnhuán lìyòng} \textbf{biànchéng} xīn chǎnpǐn.} \\
    \textit{Étape 3 : Enfin, l'usine traite les recyclables en nouvelles matières premières, les recycle pour en faire de nouveaux produits.}
\end{enumerate}

\textbf{4. Pied de page}
\zh{Terminale A -- Groupe 3 -- Lycée Bilingue de Nkolbisson -- 2025}

\noindent \textbf{Commentaire détaillé des choix linguistiques (Corrigé commenté)~:}

\begin{itemize}
    \item \textbf{Slogan~:} \zh{垃圾分类，从我做起} est la structure standard des campagnes officielles chinoises (ex. Shanghai 2019). Elle utilise la structure \cle{从\ldots 做起} vue en grammaire. Court, rythmé (4+4 syllabes), facile à mémoriser.
    \item \textbf{Vocabulaire poubelles~:} Utilisation exhaustive du lexique de la leçon 9 (\zh{可回收物, 有害垃圾, 厨余垃圾, 其他垃圾}). Les exemples sont des mots de fréquence élevée (HSK 2-3) : \zh{塑料瓶, 废纸, 金属罐, 废电池, 过期药品, 剩饭菜, 果皮, 卫生纸}. L'association \textbf{Couleur + Nom catégorie} est la norme visuelle chinoise (标准色).
    \item \textbf{Phrase 1 (Sujet \zh{我们})~:} Utilisation de \cle{先} pour marquer le point de départ. Structure \cle{把} (\zh{把 生活垃圾 分类投放 到\ldots 里}) : l'objet \zh{生活垃圾} est connu/défini, l'action \zh{分类投放} (verbe composé V+V directionnel/résultatif) est achevée, le complément de lieu \zh{到对应的垃圾桶里} précise la destination. \textbf{Point de grammaire~:} \zh{分类投放} = \zh{分类} (manière) + \zh{投放} (action).
    \item \textbf{Phrase 2 (Sujet \zh{环卫车})~:} \cle{然后} assure la transition. Nouveau \cle{把} : \zh{把 分类好 的 垃圾 收集运输 到 处理厂}. Ici \zh{分类好} est un complément d'état (verbe + \cle{好}) qui qualifie \zh{垃圾} (les déchets \textbf{bien triés}). \zh{收集运输} est une série verbale (collecter ET transporter). \zh{到处理厂} indique le terminus.
    \item \textbf{Phrase 3 (Sujet \zh{工厂})~:} \cle{最后} marque la fin. Double structure \cle{把} imbriquée logiquement~: 1. \zh{把 可回收物 处理 成 新 原料} (Traiter l'objet $\to$ résultat \cle{成}). 2. \zh{循环利用 变成 新 产品} (Le recyclage $\to$ transformation \cle{变成}). L'usage de \zh{原料} (matière première) et \zh{产品} (produit) montre la maîtrise du lexique technique de la leçon.
    \item \textbf{Cohérence textuelle~:} La chaîne de référence est claire : \zh{我们} (citoyens) $\to$ \zh{环卫车} (service public) $\to$ \zh{工厂} (industrie). L'objet \zh{垃圾} $\to$ \zh{分类好 的 垃圾} $\to$ \zh{可回收物} $\to$ \zh{新原料} $\to$ \zh{新产品} suit la transformation physique réelle.
    \item \textbf{Registre~:} Langage standard, écrit, adapté à l'affichage public (pas de familier, pas de contractions orales comme \zh{咱们}).
\end{itemize}

\textbf{Script oral type pour le présentateur (1 minute)~:}
\zh{各位同学，老师们，大家好。我们组的海报主题是：「垃圾分类，从我做起」。请看中间的四个垃圾桶：蓝色桶装可回收物，比如塑料瓶、废纸、金属罐；红色桶装有害垃圾，比如废电池、过期药品；绿色桶装厨余垃圾，比如剩饭菜、果皮；灰色桶装其他垃圾，比如卫生纸、烟蒂。下面是回收流程：第一步，我们先把生活垃圾分类投放到对应的垃圾桶里；第二步，环卫车然后把分类好的垃圾收集运输到处理厂；第三步，工厂最后把可回收物处理成新原料，循环利用变成新产品。保护环境，人人有责。请大家行动起来！谢谢。}
\end{exoresolu}

\courssub{Bilan de l'activité}
\noindent À l'issue de cette activité, l'enseignant mènera une brève méta-cognition collective (5 min) en posant les questions suivantes à la classe~:
\begin{enumerate}
    \item \textbf{Sur le lexique~:} « Quel mot de la leçon 9 n'avez-vous \textbf{pas} réussi à caser naturellement dans l'affiche ? Pourquoi ? » (Ex. \zh{降解} dégrader, \zh{焚烧} incinérer, \zh{填埋} enfouir, \zh{粉碎} broyer, \zh{熔融} fondre -- mots du texte support mais moins utiles pour une affiche grand public simplifiée).
    \item \textbf{Sur la grammaire~:} « Dans la phrase 3 du modèle, pourquoi utilise-t-on \cle{成} après \cle{处理} et \cle{变成} pour la seconde partie ? » (Différence entre \cle{处理成} « transformer \textbf{en} (matériau/état intermédiaire) » -- résultat de fabrication industrielle -- et \cle{变成} « devenir \textbf{en} (produit fini) » -- changement d'état final, plus large).
    \item \textbf{Sur la prononciation~:} « Quel ton a posé le plus de problème à l'oral ? » (Souvent le ton 3 de \zh{处} chǔ / \zh{厨} chú / \zh{垃} lǎ sandhi, ou le ton 2 de \zh{回} huí / \zh{循} xún, ou la neutralisation de \zh{里} lǐ).
    \item \textbf{Sur la culture~:} « La Chine impose le tri à Shanghai depuis 2019 avec amendes. Au Cameroun, HYSACAM gère la collecte. Selon vous, quelle étape du processus \cle{先/然后/最后} est la plus fragile chez nous actuellement ? » (Souvent : le \zh{分类投放} à la source -- étape 1 -- manque d'habitude/poubelles ; et le \zh{循环利用} industriel -- étape 3 -- manque d'usines locales de transformation).
\end{enumerate}

\noindent L'enseignant collecte les affiches, note la production écrite (grille : Lexique /10, Grammaire /10, Caractères/Pinyin /5, Visuel /5) et la production orale (grille : Prononciation /10, Fluidité /5, Contenu /5, Interaction /5). L'affiche lauréate est plastifiée et affichée : elle devient une trace matérielle de la compétence « Écrire en chinois pour agir dans la cité », objectif terminal du module 2.

\newpage

\coursec{Méthodes et exercices résolus}

\coursec{Leçon 9 — Vocabulaire et compréhension de texte : processus du recyclage}

\courssub{Méthode 1 : Lire et comprendre un texte descriptif et séquentiel sur un processus technique}

\begin{methode}[Démarche de lecture compréhensive : identifier les étapes d'un processus]
\textbf{Objectif :} Extraire l'information principale, repérer la structure chronologique et comprendre le vocabulaire technique d'un texte expositif sur le recyclage.

\textbf{Étapes :}
\begin{enumerate}
    \item \textbf{Observation paratextuelle :} Lire le titre (\zh{垃圾分类与回收利用的过程}), repérer les mots-clés en gras (\zh{分类}, \zh{回收}, \zh{处理}, \zh{再生}) et les connecteurs logiques (\zh{首先}, \zh{然后}, \zh{最后}).
    \item \textbf{Lecture globale (skimming) :} Lire une première fois sans s'arrêter au vocabulaire inconnu. Identifier le sujet (le recyclage des déchets) et le plan global (collecte $\to$ tri $\to$ traitement $\to$ fabrication).
    \item \textbf{Lecture détaillée (scanning) :} Relever les verbes d'action (\zh{投放}, \zh{运送}, \zh{粉碎}, \zh{熔化}, \zh{制成}) et les compléments d'objet. Souligner les indicateurs de séquence temporelle.
    \item \textbf{Inférence lexicale :} Deviner le sens des mots inconnus grâce aux radicaux (ex: \zh{塑} clé \zh{扌} main/action + \zh{谷} phonétique $\to$ plastique ; \zh{熔} clé \zh{火} feu $\to$ fondre) et au contexte.
    \item \textbf{Synthèse :} Reconstituer le schéma du processus en 4-5 étapes en français, puis en chinois avec les connecteurs logiques.
\end{enumerate}

\textbf{Structures à retenir :}
\begin{itemize}
    \item \zh{首先……，其次……，然后……，最后……} (D'abord..., ensuite..., puis..., enfin...) pour ordonner un processus.
    \item \zh{被 (bèi)} passif fréquent dans les descriptions techniques : \zh{垃圾被分类} (Les déchets sont triés).
    \item \zh{把 (bǎ)} pour focaliser sur la manipulation de l'objet : \zh{工人把塑料瓶粉碎} (Les ouvriers broient les bouteilles en plastique).
\end{itemize}
\end{methode}

\begin{popfigure}
\begin{tikzpicture}[
    node distance=1.2cm and 0.5cm,
    box/.style={rectangle, rounded corners, draw=popblue, thick, fill=popblueL, text width=2.8cm, align=center, font=\small},
    arrow/.style={-Stealth, thick, popdark}
]
\node[box, align=center] (start){Déchets ménagers\\\zh{生活垃圾}};
\node[box, below=of start, align=center] (sort){Tri à la source\\4 catégories\\\zh{分类投放}};
\node[box, below=of sort, align=center] (collect){Collecte \& Transport\\\zh{收集运送}};
\node[box, below=of collect, align=center] (center){Centre de tri\\\zh{回收站 / 分拣}};
\node[box, below left=of center, xshift=-1cm, align=center] (paper){Papier\\\zh{纸张 : 浆化$\to$脱墨$\to$造纸}};
\node[box, below=of center, align=center] (plastic){Plastique\\\zh{塑料 : 清洗$\to$粉碎$\to$熔化$\to$拉丝}};
\node[box, below right=of center, xshift=1cm, align=center] (metalglass){Métal / Verre\\\zh{金属/玻璃 : 熔炼/熔化$\to$铸造/重塑}};
\node[box, below=of plastic, yshift=-0.5cm, align=center] (factory){Usines de fabrication\\\zh{制成新制品}};
\node[box, below=of factory, align=center] (end){Marché / Boucle fermée\\\zh{市场循环 / 闭环}};

\draw[arrow] (start) -- (sort);
\draw[arrow] (sort) -- (collect);
\draw[arrow] (collect) -- (center);
\draw[arrow] (center) -- (paper);
\draw[arrow] (center) -- (plastic);
\draw[arrow] (center) -- (metalglass);
\draw[arrow] (paper) -- (factory);
\draw[arrow] (plastic) -- (factory);
\draw[arrow] (metalglass) -- (factory);
\draw[arrow] (factory) -- (end);
\draw[arrow, poporange] (end.west) .. controls +(-2,0) and +(-2,0) .. (start.west) node[midway, left, font=\tiny, poporange] {\zh{循环经济}};
\end{tikzpicture}
\legende{Schéma du processus de recyclage : de la collecte à la réintégration au marché (économie circulaire).}
\end{popfigure}

\begin{exoresolu}[Compréhension globale et détaillée du texte support]
\textbf{Énoncé :} Lisez le texte ci-dessous et répondez aux questions en chinois (phrases complètes) ou en français selon la consigne.

\begin{textebox}[Processus du recyclage des déchets ménagers]
垃圾分类与回收利用的过程

在喀麦隆的雅温得和杜阿拉，垃圾分类正在逐步推广。居民首先要把生活垃圾分为四类：可回收物、厨余垃圾、有害垃圾和其他垃圾。\textbf{首先}，环卫工人把分类好的可回收物（如纸张、塑料瓶、金属罐、玻璃瓶）收集起来，运送到回收站。\textbf{其次}，在回收站，工作人员把这些物品进行二次精细分拣。\textbf{然后}，不同材料进入不同的处理流程：纸张被浆化、脱墨、烘干，重新造纸；塑料被清洗、粉碎、熔化、拉丝，制成塑料颗粒或新制品；金属被熔炼、铸造；玻璃被粉碎、熔化、重塑。\textbf{最后}，这些再生原料被送往工厂，制成新的纸张、塑料制品、金属构件或玻璃器皿，重新进入市场循环。厨余垃圾则被送往堆肥厂或沼气站，变成有机肥料或清洁能源。通过这个闭环过程，我们减少了填埋量，节约了自然资源，保护了环境。

\textit{Pinyin :}
Lājī fēnlèi yǔ huíshōu lìyòng de guòchéng

Zài Kāmàilóng de Yǎwēndé hé Dūā'ěr, lājī fēnlèi zhèngzài zhúbù tuīguǎng. Jūmín shǒuxiān yào bǎ shēnghuó lājī fēn wéi sì lèi: kě huíshōu wù, chúyú lājī, yǒuhài lājī hé qítā lājī. \textbf{Shǒuxiān}, huánwèi gōngrén bǎ fēnlèi hǎo de kě huíshōu wù (rú zhǐzhāng, sùliáo píng, jīnshǔ guàn, bōlí píng) shōují qǐlái, yùnsòng dào huíshōuzhàn. \textbf{Qícì}, zài huíshōuzhàn, gōngzuò rényuán bǎ zhèxiē wùpǐn jìnxíng èrcì jīngxì fēnlǎn. \textbf{Ránhòu}, bùtóng cáiliào jìnrù bùtóng de chǔlǐ liúchéng: zhǐzhāng bèi jiānghuà, tuōmò, hōnggān, chóngxīn zàozhǐ; sùliào bèi qǐxǐ, fěnsuì, rónghuà, lāsī, zhìchéng sùliào kēlì huò xīn zhìpǐn; jīnshǔ bèi róngliàn, zhùzào; bōlí bèi fěnsuì, rónghuà, chóngsù. \textbf{Zuìhòu}, zhèxiē zàishēng yuánliào bèi sòngwǎng gōngchǎng, zhìchéng xīn de zhǐzhāng, sùliào zhìpǐn, jīnshǔ gòujiàn huò bōlí qìwǎn, chóngxīn jìnrù shìchǎng xúnhuán. Chúyú lājī zé bèi sòngwǎng duīféi chǎng huò zhàoqì zhàn, biànchéng yǒujī fèiliào huò qīngjié néngyuán. Tōngguò zhège bìhuán guòchéng, wǒmen jiǎnshǎole tiántáng liàng, jiéyuēle zìrán zīyuán, bǎohùle huánjìng.

\textit{Traduction :}
Le processus de tri et de recyclage des déchets

À Yaoundé et Douala au Cameroun, le tri des déchets est progressivement promu. Les habitants doivent d'abord diviser les déchets ménagers en quatre catégories : recyclables, déchets de cuisine, déchets dangereux et autres déchets. \textbf{D'abord}, les agents d'hygiène collectent les recyclables triés (papier, bouteilles en plastique, canettes métalliques, bouteilles en verre) et les transportent vers le centre de recyclage. \textbf{Ensuite}, au centre, le personnel effectue un tri fin secondaire. \textbf{Puis}, les différents matériaux entrent dans différentes filières de traitement : le papier est pulpé, désencré, séché, pour refaire du papier ; le plastique est lavé, broyé, fondu, filé, pour devenir des granulés ou nouveaux produits ; le métal est fondu, coulé ; le verre est broyé, fondu, remodelé. \textbf{Enfin}, ces matières premières régénérées sont envoyées aux usines pour fabriquer du nouveau papier, produits plastiques, pièces métalliques ou ustensiles en verre, réintégrant le circuit du marché. Les déchets de cuisine sont envoyés aux usines de compostage ou stations de biogaz, devenant engrais organique ou énergie propre. Par ce processus en boucle fermée, nous réduisons l'enfouissement, économisons les ressources naturelles, protégeons l'environnement.
\end{textebox}

\textbf{Questions :}
\begin{enumerate}
    \item Combien de catégories de déchets les habitants doivent-ils trier au départ ? Citez-les en chinois.
    \item Quel connecteur logique introduit l'étape du transport vers le centre de recyclage ?
    \item Transformez la phrase \zh{工人把塑料瓶粉碎} à la voix passive avec \zh{被}.
    \item Selon le texte, que deviennent les déchets de cuisine (\zh{厨余垃圾}) ? Répondez en français.
    \item Trouvez dans le texte le verbe signifiant « broyer / réduire en poudre » et écrivez son pinyin avec tons.
    \item Expliquez en français l'expression \zh{闭环过程} (bìhuán guòchéng) dans ce contexte.
\end{enumerate}

\tcblower

\textbf{Solution détaillée :}

\begin{enumerate}
    \item \textbf{Quatre catégories.} Le texte dit : \zh{居民首先要把生活垃圾分为四类：可回收物 (kě huíshōu wù)、厨余垃圾 (chúyú lājī)、有害垃圾 (yǒuhài lājī) 和其他垃圾 (qítā lājī)。}
    \item \textbf{\zh{首先} (shǒuxiān)}. C'est le premier connecteur de la séquence narrative qui introduit l'action des éboueurs. \zh{其次} (qícì) introduit l'étape suivante (tri fin), \zh{然后} (ránhòu) le traitement matière par matière, \zh{最后} (zuìhòu) la fabrication finale.
    \item \textbf{Structure \zh{把} $\to$ \zh{被} :} Sujet (objet manipulé) + \zh{被} + Agent + Verbe (+ Complément de résultat).
    \zh{塑料瓶被工人粉碎了。} (Sùliáo píng bèi gōngrén fěnsuì le.) \textit{Note : l'aspect accompli \zh{了} est naturel ici car l'action est terminée dans le processus.}
    \item Les déchets de cuisine sont envoyés à l'usine de compostage (\zh{堆肥厂}) ou à la station de biogaz (\zh{沼气站}) pour être transformés en engrais organique (\zh{有机肥料}) ou en énergie propre (\zh{清洁能源}).
    \item \textbf{\zh{粉碎} (fěn suì)}. Trouvé dans : \zh{塑料被清洗、粉碎、熔化、拉丝} et \zh{玻璃被粉碎、熔化、重塑}. Radicaux : \zh{米} (riz/éclats) + \zh{雹} (grêle/phonétique) pour \zh{粉} ; \zh{石} (pierre) + \zh{卒} (phonétique) pour \zh{碎}.
    \item \zh{闭环过程} (Processus en boucle fermée / Économie circulaire) : Cela signifie que les déchets ne sont pas une fin en soi (enfouissement), mais redeviennent des ressources (matières premières secondaires) pour fabriquer de nouveaux produits, qui redeviendront des déchets, etc. Le cycle est « fermé » : Production $\to$ Consommation $\to$ Déchet $\to$ Recyclage $\to$ Matière première $\to$ Production.
\end{enumerate}

\textbf{Conclusion :} La maîtrise des connecteurs logiques (\zh{首先/其次/然后/最后}) et des structures \zh{把}/\zh{被} est indispensable pour décrire un processus technique avec précision.
\end{exoresolu}

\courssub{Méthode 2 : Maîtriser le vocabulaire thématique : classification, matériaux, verbes d'action technique}

\begin{methode}[Apprentissage lexical structuré par champs sémantiques et radicaux]
\textbf{Objectif :} Mémoriser durablement le lexique du recyclage en l'organisant par catégories grammaticales et sémantiques, et en analysant la composition des caractères (clés/radicaux).

\textbf{Étapes :}
\begin{enumerate}
    \item \textbf{Classer par nature grammaticale :} Séparer les Noms (matériaux, lieux, machines), Verbes (actions humaines, transformations industrielles), Adjectifs/Adverbes (qualité, séquence).
    \item \textbf{Grouper par champ sémantique :}
        \begin{itemize}
            \item \textit{Catégories de déchets :} \zh{可回收物}, \zh{厨余垃圾}, \zh{有害垃圾}, \zh{其他垃圾}.
            \item \textit{Matériaux :} \zh{纸张}, \zh{塑料}, \zh{金属}, \zh{玻璃}.
            \item \textit{Verbes de transformation industrielle (chaîne opératoire) :} \zh{分拣} $\to$ \zh{清洗} $\to$ \zh{粉碎/破碎} $\to$ \zh{熔化/熔炼} $\to$ \zh{拉丝/铸造/重塑} $\to$ \zh{制成}.
            \item \textit{Lieux :} \zh{回收站}, \zh{堆肥厂}, \zh{沼气站}, \zh{填埋场}.
        \end{itemize}
    \item \textbf{Analyser les radicaux (clés) pour deviner le sens :}
        \begin{itemize}
            \item \zh{塑} (sù, plastique) : clé \zh{扌} (main/action) + phonétique \zh{谷}. Action de pétrir la matière.
            \item \zh{熔} (róng, fondre) : clé \zh{火} (feu/heat) + phonétique \zh{荣}. Action du feu sur le métal/verre/plastique.
            \item \zh{堆} (duī, tas/tasser) : clé \zh{木} (bois/matière) + phonétique \zh{隹}. Accumulation de matière organique.
            \item \zh{沼} (zhào, marais/biogaz) : clé \zh{氵} (eau/liquide) + phonétique \zh{召}. Liquide dans un trou (méthanisation).
        \end{itemize}
    \item \textbf{Travailler les collocations (associations fixes) :} \zh{进行分类}, \zh{投放垃圾}, \zh{节约资源}, \zh{保护环境}, \zh{再生资源}.
    \item \textbf{Entraînement écriture :} Respecter l'ordre des traits (haut $\to$ bas, gauche $\to$ droite, horizontal avant vertical, cadre avant contenu). Ex: \zh{回} (6 traits : 口 $\to$ 口 $\to$ 囗).
\end{enumerate}
\end{methode}

\begin{center}
\begin{tabularx}{\linewidth}{|X|X|X|X|}
\hline
\rowcolor{popblueL} \textbf{Caractères} & \textbf{Pinyin} & \textbf{Nature} & \textbf{Traduction} \\
\hline
垃圾 & lājī & Nom & Déchet, ordures \\
分类 & fēnlèi & Verbe / Nom & Trier / Classification \\
回收 & huíshōu & Verbe & Recycler, récupérer \\
利用 & lìyòng & Verbe & Utiliser, valoriser \\
可回收物 & kě huíshōu wù & Nom & Déchets recyclables / Matières récupérables \\
厨余垃圾 & chúyú lājī & Nom & Déchets de cuisine (épluchures, restes) \\
有害垃圾 & yǒuhài lājī & Nom & Déchets dangereux (piles, médicaments) \\
其他垃圾 & qítā lājī & Nom & Autres déchets (résidus ultimes) \\
环卫工人 & huánwèi gōngrén & Nom & Agent d'hygiène / Éboueur \\
回收站 & huíshōuzhàn & Nom & Centre de collecte / Déchetterie \\
分拣 & fēnlǎn & Verbe & Trier (finement), séparer \\
纸张 & zhǐzhāng & Nom & Papier \\
塑料 & sùliào & Nom & Plastique \\
金属 & jīnshǔ & Nom & Métal \\
玻璃 & bōlí & Nom & Verre \\
浆化 & jiānghuà & Verbe & Pulpifier (papier) \\
脱墨 & tuōmò & Verbe & Désencrer \\
粉碎 & fěnsuì & Verbe & Broyer, pulvériser \\
熔化 & rónghuà & Verbe & Fondre (verre, plastique) \\
熔炼 & róngliàn & Verbe & Fondre / Métallurgie (métal) \\
铸造 & zhùzào & Verbe & Couler (métal en moule) \\
拉丝 & lāsī & Verbe & Filer / Extruder (plastique) \\
颗粒 & kēlì & Nom & Granulés (matière première plastique) \\
再生原料 & zàishēng yuánliào & Nom & Matières premières secondaires / régénérées \\
堆肥厂 & duīféi chǎng & Nom & Usine de compostage \\
沼气站 & zhàoqì zhàn & Nom & Station de biogaz / méthanisation \\
有机肥料 & yǒujī fèiliào & Nom & Engrais organique \\
清洁能源 & qīngjié néngyuán & Nom & Énergie propre / renouvelable \\
填埋 & tiántáng & Verbe & Enfouir (décharge) \\
闭环 & bìhuán & Adj / Nom & Boucle fermée / Circulaire \\
节约 & jiéyuē & Verbe & Économiser \\
\hline
\end{tabularx}
\end{center}

\begin{exoresolu}[Exercices lexicaux : classification, radicaux, collocations]
\textbf{Énoncé :}
\begin{enumerate}
    \item \textbf{Classement :} Rangez les mots suivants dans le bon tableau : \zh{塑料, 分拣, 熔炼, 堆肥厂, 金属, 环卫工人, 纸张, 清洗, 回收站, 铸造}.
    \item \textbf{Radicaux :} Associez chaque caractère à sa clé (radical) et déduisez le sens général lié à cette clé.
    \begin{itemize}
        \item \zh{熔} (róng) \quad \zh{堆} (duī) \quad \zh{沼} (zhào) \quad \zh{粉} (fěn)
    \end{itemize}
    \item \textbf{Collocations :} Complétez avec le verbe qui colle le mieux (un seul usage par verbe) : \zh{进行 / 投放 / 节约 / 保护 / 制成}.
    \begin{enumerate}
        \item 居民要按时 \_\_\_\_\_\_ 垃圾。
        \item 工厂用再生料 \_\_\_\_\_\_ 新产品。
        \item 我们要 \_\_\_\_\_\_ 自然资源。
        \item 学校 \_\_\_\_\_\_ 了垃圾分类宣传活动。
        \item 大家一起 \_\_\_\_\_\_ 环境。
    \end{enumerate}
    \item \textbf{Ordre des traits :} Écrivez l'ordre des traits pour le caractère \zh{回} (huí) et \zh{收} (shōu).
    \item \textbf{Production :} Écrivez une phrase en chinois (caractères + pinyin) pour décrire ce qui arrive aux bouteilles en plastique après le tri fin, en utilisant \zh{被} et au moins deux verbes de transformation du texte (\zh{清洗, 粉碎, 熔化, 拉丝}).
\end{enumerate}

\tcblower

\textbf{Solution détaillée :}

\begin{enumerate}
    \item \textbf{Classement :}
    \begin{center}
    \begin{tabular}{|l|l|l|}
    \hline
    \rowcolor{popblueL} \textbf{Matériaux (Noms)} & \textbf{Verbes d'action / Traitement} & \textbf{Lieux / Acteurs (Noms)} \\
    \hline
    \zh{塑料} (plastique) & \zh{分拣} (trier finement) & \zh{堆肥厂} (usine compostage) \\
    \zh{金属} (métal) & \zh{熔炼} (fondre métal) & \zh{环卫工人} (éboueur) \\
    \zh{纸张} (papier) & \zh{清洗} (laver) & \zh{回收站} (centre tri) \\
    \hline
    \multicolumn{3}{|l|}{\textit{Note : \zh{铸造} (couler métal) est aussi un verbe de traitement.}} \\
    \hline
    \end{tabular}
    \end{center}

    \item \textbf{Radicaux :}
    \begin{center}
    \begin{tabular}{|c|c|c|c|}
    \hline
    \rowcolor{popblueL} \textbf{Caractère} & \textbf{Clé (Radical)} & \textbf{Sens de la clé} & \textbf{Lien avec le sens du caractère} \\
    \hline
    \zh{熔} & \zh{火} (huǒ, feu) & Feu, chaleur, haute température & Fondre nécessite le feu/chaud. \\
    \zh{堆} & \zh{木} (mù, bois/arbre) & Bois, matière végétale, empilement & Un tas de bois/feuilles/matière. \\
    \zh{沼} & \zh{氵} (shuǐ, eau) & Eau, liquide, rivière & Marais, biogaz = liquide/gaz dans l'eau/sol. \\
    \zh{粉} & \zh{米} (mǐ, riz) & Riz, grain, poudre fine & Poudre fine comme de la farine de riz. \\
    \hline
    \end{tabular}
    \end{center}

    \item \textbf{Collocations :}
    \begin{enumerate}
        \item \zh{投放} (tóufàng) : « déposer / jeter (au bon endroit) ». \zh{居民要按时 \textbf{投放} 垃圾。}
        \item \zh{制成} (zhìchéng) : « fabriquer / transformer en ». \zh{工厂用再生料 \textbf{制成} 新产品。}
        \item \zh{节约} (jiéyuē) : « économiser ». \zh{我们要 \textbf{节约} 自然资源。}
        \item \zh{进行} (jìnxíng) : « mener / réaliser (activité) ». \zh{学校 \textbf{进行} 了垃圾分类宣传活动。}
        \item \zh{保护} (bǎohù) : « protéger ». \zh{大家一起 \textbf{保护} 环境。}
    \end{enumerate}

    \item \textbf{Ordre des traits :}
    \begin{itemize}
        \item \zh{回} (huí, 6 traits) : 1. \zh{丨} (verticale) 2. \zh{(trait brisé)} (horizontale + crochet bas) 3. \zh{丨} (verticale milieu) 4. \zh{一} (horizontale milieu) 5. \zh{一} (horizontale bas) 6. \zh{(trait brisé)} (fermeture cadre bas + horizontale finale). \textit{Règle : cadre \zh{囗} s'écrit en 3 traits (gauche-haut, haut-droite, bas-fermeture), contenu avant fermeture finale.}
        \item \zh{收} (shōu, 6 traits) : 1. \zh{丨} (verticale gauche) 2. \zh{(trait brisé)} (horizontale + crochet) 3. \zh{丿} (pié gauche) 4. \zh{丶} (point) 5. \zh{一} (horizontale) 6. \zh{(trait brisé)} (verticale + crochet bas droite). \textit{Clé \zh{攵} (pū, action main/bâton) à droite : traits 3 à 6.}
    \end{itemize}

    \item \textbf{Production :}
    \zh{经过二次分拣后，塑料瓶被清洗、粉碎、熔化、拉丝，制成塑料颗粒。}
    (Jīngguò èrcì fēnlǎn hòu, sùliáo píng bèi qǐxǐ, fěnsuì, rónghuà, lāsī, zhìchéng sùliào kēlì.)
    \textit{Traduction : Après le tri secondaire, les bouteilles en plastique sont lavées, broyées, fondues, filées, et transformées en granulés de plastique.}
    \textit{Justification :} Utilisation de \zh{被} pour la voix passive (focus sur l'objet \zh{塑料瓶}), série de verbes (\zh{清洗、粉碎、熔化、拉丝}) marquant les étapes techniques, \zh{制成} pour le résultat final.
\end{enumerate}
\end{exoresolu}

\courssub{Méthode 3 : Répondre aux questions de compréhension : localiser l'information et reformuler}

\begin{methode}[Techniques de réponse aux questions de compréhension écrite (Bac)]
\textbf{Objectif :} Répondre avec précision, en phrases complètes (si exigé), en reprenant le vocabulaire du texte et en évitant le hors-sujet.

\textbf{Étapes :}
\begin{enumerate}
    \item \textbf{Analyser la question :} Identifier le mot-interrogatif (\zh{谁} qui, \zh{什么} quoi, \zh{怎么} comment, \zh{为什么} pourquoi, \zh{哪里} où, \zh{多少} combien) et le temps/aspect visé.
    \item \textbf{Localiser l'information dans le texte :} Scanner les mots-clés de la question dans le texte (synonymes, antonymes). Souligner la phrase source.
    \item \textbf{Choisir la langue de réponse :} Si « Répondez en chinois » $\to$ Phrase complète en caractères + pinyin. Si « Répondez en français » $\to$ Phrase claire, vocabulaire précis.
    \item \textbf{Reformuler (ne pas copier-coller aveuglément) :} Adapter la structure. Ex: Texte: \zh{工人把垃圾运走} $\to$ Question: \zh{垃圾怎么处理的?} $\to$ Réponse: \zh{垃圾被工人运走了。} (Passif \zh{被} plus naturel pour focus sur l'objet).
    \item \textbf{Vérifier :} Accord temps/aspect (\zh{了}, \zh{过}, \zh{着}), ordre des mots (Complément de lieu/temps AVANT le verbe), classificateurs.
\end{enumerate}

\textbf{Réflexes « Bac » :}
\begin{itemize}
    \item Question \zh{为什么} (Pourquoi) $\to$ Réponse commence par \zh{因为} (yīnwèi) / \zh{由于} (yóuyú).
    \item Question \zh{怎么} / \zh{怎样} (Comment) $\to$ Réponse décrit la manière/le processus (\zh{首先...然后...}).
    \item Question \zh{变成什么了} (Devenir quoi) $\to$ Verbe \zh{变成} (biànchéng) ou \zh{制成} (zhìchéng).
\end{itemize}
\end{methode}

\begin{exoresolu}[Questions de compréhension type examen : localisation et reformulation]
\textbf{Énoncé :} Reprenez le texte « Processus du recyclage ». Répondez \textbf{en chinois} (caractères + pinyin) aux questions 1-3, et \textbf{en français} aux questions 4-5.

\begin{enumerate}
    \item 谁负责收集可回收物？ (Qui collecte les recyclables ?)
    \item 可回收物被运送到哪里？ (Où sont transportés les recyclables ?)
    \item 纸张在回收站经过哪三个步骤才能重新造纸？ (Quelles sont les 3 étapes pour le papier ?)
    \item Pourquoi le texte mentionne-t-il Yaoundé et Douala spécifiquement ? (Contexte camerounais).
    \item Quel est le bénéfice environnemental final mentionné dans la dernière phrase ? Citez les trois éléments.
\end{enumerate}

\tcblower

\textbf{Solution détaillée :}

\begin{enumerate}
    \item \textbf{Localisation :} Phrase 3 : \zh{环卫工人把分类好的可回收物...收集起来}.\par
    \textbf{Réponse :} \zh{环卫工人负责收集可回收物。} (Huánwèi gōngrén fùzé shōují kě huíshōu wù.)\par
    \textit{Analyse :} \zh{谁} appelle le sujet. \zh{负责} (être en charge de) structure la réponse. \zh{把} structure disparait car le focus est sur l'agent.
    \item \textbf{Localisation :} Phrase 3 : \zh{运送到回收站}.\par
    \textbf{Réponse :} \zh{可回收物被运送到回收站。} (Kě huíshōu wù bèi yùnsòng dào huíshōuzhàn.)\par
    \textit{Analyse :} Question \zh{哪里} (où). Complément de lieu \zh{到回收站}. Voix passive \zh{被} naturelle car l'objet \zh{可回收物} est le sujet de la phrase réponse.
    \item \textbf{Localisation :} Phrase 5 : \zh{纸张被浆化、脱墨、烘干，重新造纸}.\par
    \textbf{Réponse :} \zh{纸张被浆化、脱墨、烘干，然后重新造纸。} (Zhǐzhāng bèi jiānghuà, tuōmò, hōnggān, ránhòu chóngxīn zàozhǐ.)\par
    \textit{Analyse :} Question \zh{哪三个步骤} (quelles 3 étapes). Liste les 3 V + \zh{重新造纸} (résultat). Utilisation de \zh{被} (passif) et \zh{然后} (liaison).
    \item \textbf{Contexte :} Yaoundé (capitale politique) et Douala (capitale économique, port) sont les deux plus grandes villes du Cameroun, générant le plus de déchets. Le texte ancre le processus dans la réalité locale pour montrer que le tri \zh{正在逐步推广} (est en cours de promotion progressive).
    \item \textbf{Synthèse finale :} Dernière phrase : \zh{减少了填埋量，节约了自然资源，保护了环境}.\par
    \textbf{Réponse :} Les trois bénéfices sont : 1. Réduire la quantité de déchets enfouis (réduire l'enfouissement). 2. Économiser les ressources naturelles. 3. Protéger l'environnement.
\end{enumerate}

\textbf{Conclusion :} La réponse en chinois exige une manipulation syntaxique (passif \zh{被}, ordre des compléments) et non une simple copie. La réponse en français exige une synthèse précise du sens global.
\end{exoresolu}

\courssub{Méthode 4 : Décrire un processus technique à l'oral et à l'écrit (Production guidée)}

\begin{methode}[Structurer une description de processus : du schéma à la phrase]
\textbf{Objectif :} Produire un exposé oral ou un paragraphe écrit cohérent décrivant le cycle de vie d'un déchet (ex: une canette, un sac plastique, un journal).

\textbf{Étapes :}
\begin{enumerate}
    \item \textbf{Choisir l'objet :} Un déchet unique (ex: \zh{一个易拉罐} - une canette, \zh{一张旧报纸} - un vieux journal).
    \item \textbf{Définir les 5 étapes clés (Schéma) :}
        \begin{center}
        \zh{使用} $\to$ \zh{分类投放} $\to$ \zh{收集运输} $\to$ \zh{工厂处理} $\to$ \zh{再生产品}
        \end{center}
    \item \textbf{Sélectionner les verbes précis par matériau :}
        \begin{itemize}
            \item Métal (\zh{金属/易拉罐}) : \zh{熔炼} (fondre), \zh{铸造} (couler), \zh{轧制} (laminer).
            \item Papier (\zh{纸张/报纸}) : \zh{浆化} (pulpifier), \zh{脱墨} (désencrer), \zh{造纸} (fabriquer papier).
            \item Plastique (\zh{塑料/瓶子}) : \zh{清洗} (laver), \zh{粉碎} (broyer), \zh{熔化} (fondre), \zh{拉丝/注塑} (filer/mouler par injection).
        \end{itemize}
    \item \textbf{Rédiger / Parler en utilisant la structure « Sujet + Connecteur + Verbe (+ Complément) » :}
        \zh{首先，我把空罐子投放到可回收箱。其次，环卫工人把它运到回收站。然后，罐子被熔炼、铸造，变成新的铝材料。最后，它被制成新的易拉罐，回到超市。}
    \item \textbf{Vérifier les points de grammaire critiques :}
        \begin{itemize}
            \item \zh{把} structure pour les actions humaines sur objet précis (\zh{我把罐子投放...}).
            \item \zh{被} structure pour les étapes industrielles où l'agent est inconnu/secondaire (\zh{罐子被熔炼...}).
            \item Connecteurs : \zh{首先、其次、然后、最后}.
            \item Aspect \zh{了} pour actions terminées, \zh{着} pour états résultants si besoin.
        \end{itemize}
\end{enumerate}
\end{methode}

\begin{exoresolu}[Production écrite : « La vie d'une bouteille en plastique »]
\textbf{Énoncé :} Rédigez un court paragraphe (60-80 caractères chinois) pour décrire le processus de recyclage d'une bouteille en plastique (\zh{一个塑料瓶}), depuis le geste du citoyen jusqu'au nouveau produit. Utilisez \textbf{au moins 5 verbes du vocabulaire de la leçon}, les \textbf{4 connecteurs logiques} (\zh{首先、其次、然后、最后}), et \textbf{les deux structures \zh{把} et \zh{被}}.

\tcblower

\textbf{Proposition de rédaction (Modèle) :}

\begin{textebox}[Production élève type]
首先，我喝完水，\textbf{把}空塑料瓶\textbf{投放}到可回收箱。\\
其次，环卫工人\textbf{把}它\textbf{收集}起来，\textbf{运送}到回收站。\\
然后，工作人员\textbf{把}瓶子\textbf{清洗}、\textbf{粉碎}，瓶子\textbf{被熔化}、\textbf{拉丝}，\textbf{制成}塑料颗粒。\\
最后，颗粒\textbf{被制成}新的塑料制品，\textbf{重新进入}市场循环。

\textit{Pinyin :}
Shǒuxiān, wǒ hē wán shuǐ, \textbf{bǎ} kōng sùliáo píng \textbf{tóufàng} dào kě huíshōu xiāng.
Qícì, huánwèi gōngrén \textbf{bǎ} tā \textbf{shōují} qǐlái, \textbf{yùnsòng} dào huíshōuzhàn.
Ránhòu, gōngzuò rényuán \textbf{bǎ} píngzi \textbf{qǐxǐ}、\textbf{fěnsuì}, píngzi \textbf{bèi rónghuà}、\textbf{lāsī}, \textbf{zhìchéng} sùliào kēlì.
Zuìhòu, kēlì \textbf{bèi zhìchéng} xīn de sùliào zhìpǐn, \textbf{chóngxīn jìnrù} shìchǎng xúnhuán.

\textit{Traduction :}
D'abord, j'ai fini de boire l'eau, j'ai \textbf{jeté} la bouteille vide dans la poubelle recyclable.
Ensuite, l'éboueur \textbf{l'a collectée} et \textbf{transportée} au centre de tri.
Puis, le personnel \textbf{a lavé, broyé} la bouteille, la bouteille \textbf{a été fondue, filée}, \textbf{transformée en} granulés.
Enfin, les granulés \textbf{ont été transformés en} nouveaux produits plastiques, \textbf{réintégrant} le circuit du marché.
\end{textebox}

\textbf{Analyse grammaticale du modèle (Justification des points) :}
\begin{itemize}
    \item \textbf{Connecteurs :} \zh{首先、其次、然后、最后} (4/4 présents, bien placés en tête de phrase).
    \item \textbf{Structure \zh{把} (3 occurrences) :} \zh{把空塑料瓶投放}, \zh{把它收集起来}, \zh{把瓶子清洗、粉碎}. Sujet agent (Je, Ouvrier, Personnel) + \zh{把} + Objet + Verbe + Complément. Correct.
    \item \textbf{Structure \zh{被} (2 occurrences) :} \zh{瓶子被熔化、拉丝}, \zh{颗粒被制成}. Sujet patient (Bouteille, Granulés) + \zh{被} + Verbe. Agent omis (contexte usine). Correct.
    \item \textbf{Vocabulaire ciblé (6 verbes) :} \zh{投放, 收集, 运送, 清洗, 粉碎, 熔化, 拉丝, 制成, 重新进入}. (Largement au-dessus du minimum 5).
    \item \textbf{Complément de direction :} \zh{投放到...箱}, \zh{运送到...站}, \zh{进入...循环}. Structure Verbe + \zh{到/进} + Lieu correcte.
    \item \textbf{Particule \zh{了} :} Absente ici car description de processus général/habituel (aspect non accompli spécifique). Si c'était un récit passé : \zh{投放了, 收集了}.
\end{itemize}
\end{exoresolu}

\courssub{Méthode 5 : Traduction Thème/Version : gérer les faux-amis et l'ordre des mots}

\begin{methode}[Traduire « Environnement / Recyclage » : Français $\leftrightarrow$ Chinois]
\textbf{Objectif :} Éviter le calque syntaxique français (SVO rigide, prépositions « à/de » partout) pour produire du chinois naturel (Thème-Verbe, Compléments pré-verbaux, Sérialisation verbale).

\textbf{Principaux pièges Francophone $\to$ Chinois (Thème) :}
\begin{enumerate}
    \item \textbf{Ordre des compléments :} FR « Je jette la bouteille \textbf{dans la poubelle} » $\neq$ \zh{我扔瓶子在垃圾桶}. \textbf{CH :} \zh{我把瓶子扔\textbf{进}垃圾桶} (Verbe + Particule directionnelle + Lieu) OU \zh{我把瓶子\textbf{投放到}垃圾桶}.
    \item \textbf{Passif / Agent inconnu :} FR « Les bouteilles \textbf{sont lavées} » $\neq$ \zh{瓶子是洗}. \textbf{CH :} \zh{瓶子\textbf{被}洗\textbf{了}} / \zh{瓶子\textbf{被}清洗\textbf{了}} / \zh{洗瓶子} (Sujet patient sans marqueur passif formel, contexte clair).
    \item \textbf{Noms d'action vs Verbes :} FR « La \textbf{réduction} des déchets » $\to$ \zh{减少垃圾} (Verbe + Objet) ou \zh{垃圾减量} (Nom composé), pas \zh{减少的垃圾}.
    \item \textbf{Classificateurs obligatoires :} FR « Trois bouteilles » $\to$ \zh{三\textbf{个}瓶子} / \zh{三\textbf{只}瓶子} (pas \zh{三瓶子}).
    \item \textbf{\zh{了} vs \zh{过} :} FR « J'ai trié (hier/terminé) » $\to$ \zh{我分类\textbf{了}}. FR « J'ai déjà trié (expérience de vie) » $\to$ \zh{我分类\textbf{过}}.
\end{enumerate}

\textbf{Principaux pièges Chinois $\to$ Français (Version) :}
\begin{enumerate}
    \item \zh{把} structure : Ne pas traduire par « prendre ». \zh{我把垃圾扔了} $\to$ « J'ai jeté les ordures » (pas « J'ai pris les ordures pour jeter »).
    \item \zh{被} structure : Traduire par passif « être ... » ou « se faire ... » ou voix active indéfinie « on ... ».
    \item \zh{可回收物} : « Les matières recyclables / Les déchets recyclables » (pas « les objets pouvant être récupérés » trop lourd).
    \item \zh{厨余垃圾} : « Déchets de cuisine / Biodéchets / Déchets alimentaires » (pas « déchets de cuisine restants »).
    \item \zh{闭环} : « Boucle fermée / Économie circulaire » (terme technique).
\end{enumerate}
\end{methode}

\begin{exoresolu}[Thème et Version corrigés]
\textbf{Énoncé :}
\textbf{A. Thème (Français $\to$ Chinois) :} Traduisez en caractères + pinyin.
\begin{enumerate}
    \item Les habitants doivent trier les déchets en quatre catégories.
    \item Les bouteilles en verre sont envoyées à l'usine pour être fondues.
    \item Nous devons économiser les ressources et protéger l'environnement.
\end{enumerate}

\textbf{B. Version (Chinois $\to$ Français) :} Traduisez en français naturel.
\begin{enumerate}
    \item 塑料瓶被粉碎后，变成了塑料颗粒。
    \item 厨余垃圾可以变成有机肥料。
    \item 通过垃圾分类，我们减少了填埋量。
\end{enumerate}

\tcblower

\textbf{Solution détaillée :}

\textbf{A. Thème :}
\begin{enumerate}
    \item \zh{居民必须把垃圾分为四类。} (Jūmín bìxū bǎ lājī fēn wéi sì lèi.)
    \textit{Justification :} « Doivent » = \zh{必须} (bìxū) ou \zh{要} (yào). « Trier en 4 catégories » = \zh{把} structure + \zh{分为} (diviser en) + Nombre + Classif \zh{类}. \textit{Erreur fréquente :} \zh{分成四类} (aussi correct) vs \zh{分为四类} (plus formel pour classification). \zh{垃圾分四类} (sans \zh{把}) possible mais moins précis sur l'action volontaire.
    \item \zh{玻璃瓶被送到工厂熔化。} (Bōlí píng bèi sòng dào gōngchǎng rónghuà.)
    \textit{Justification :} Passif \zh{被} indispensable (sujet = bouteille). « Envoyées à l'usine » = \zh{送到工厂} (Verbe directionnel \zh{送} + \zh{到} + Lieu). « Pour être fondues » = but final, en chinois simple sérialisation verbale \zh{熔化} après le lieu. Pas de « pour » (\zh{为了}) nécessaire ici.
    \item \zh{我们要节约资源，保护环境。} (Wǒmen yào jiéyuē zīyuán, bǎohù huánjìng.)
    \textit{Justification :} « Doit » = \zh{要} / \zh{应该}. Coordination de deux VO (Verbe+Objet) : \zh{节约资源} / \zh{保护环境}. Pas de « et » (\zh{和}) entre verbes, virgule suffit.
\end{enumerate}

\textbf{B. Version :}
\begin{enumerate}
    \item \textbf{Après avoir été broyées, les bouteilles en plastique sont devenues des granulés de plastique.}
    \textit{Analyse :} \zh{被粉碎后} = « Après avoir été broyées » (Passif + \zh{后} = antériorité). \zh{变成了} = « sont devenues / se sont transformées en » (changement d'état \zh{了}).
    \item \textbf{Les déchets de cuisine peuvent se transformer en engrais organique. / Les biodéchets peuvent devenir de l'engrais organique.}
    \textit{Analyse :} \zh{可以} (peuvent). \zh{变成} (devenir/se transformer en). \zh{有机肥料} (engrais organique). Éviter « être changés en » (lourd).
    \item \textbf{Gracias au tri des déchets, nous avons réduit la quantité enfouie. / Le tri des déchets a permis de réduire l'enfouissement.}
    \textit{Analyse :} \zh{通过...} = « Grâce à / Au moyen de / Par ». \zh{垃圾分类} (le tri des déchets). \zh{减少了} (ont réduit / a réduit - accompli). \zh{填埋量} (quantité enfouie / volume d'enfouissement).
\end{enumerate}
\end{exoresolu}

\courssub{Méthode 6 : Éléments socioculturels : Gestion des déchets Cameroun vs Chine (Comparaison factuelle)}

\begin{methode}[Analyser et comparer les politiques environnementales : Cameroun / Chine]
\textbf{Objectif :} Situer le vocabulaire appris dans une réalité socioculturelle concrète pour l'oral (exposé/débat) et l'écrit (argumentation).

\textbf{Étapes :}
\begin{enumerate}
    \item \textbf{Identifier les faits clés (Chiffres approximatifs, lois, initiatives) :}
        \begin{itemize}
            \item \textbf{Chine :} Politique « Zéro déchet » (\zh{无废城市}) pilote depuis 2019. Tri obligatoire à Shanghai (2019) puis 46 villes clés. Taux de traitement des déchets ménagers \textgreater 99\%. Industrie du recyclage massive (\zh{再生资源}).
            \item \textbf{Cameroun :} Loi 2012/001 sur l'environnement. Décret 2014 sur la gestion des déchets. Initiatives locales (Yaoundé, Douala) : \zh{垃圾分类} (tri) en phase pilote, ONG (ex: \textit{Green Youth}, \textit{COPED}), secteur informel majeur (\zh{拾荒者} - récupérateurs informels). Défis : infrastructures de traitement (\zh{处理设施}) insuffisantes, sensibilisation.
        \end{itemize}
    \item \textbf{Structurer la comparaison (Tableau mental) :}
        \begin{center}
        \begin{tabular}{|l|l|l|}
        \hline
        \rowcolor{popblueL} \textbf{Critère} & \textbf{Chine} & \textbf{Cameroun} \\
        \hline
        Cadre légal & Strict, national, sanctions & Existe, application locale variable \\
        Traitement & Industrialisé, incinération, recyclage high-tech & Majoritairement décharge (\zh{填埋}), compostage naissant \\
        Acteurs & Entreprises d'État, privé, \zh{环卫工人} statutaires & Communes, ONG, \textbf{secteur informel} crucial \\
        Vocabulaire clé & \zh{垃圾分类强制化}, \zh{无废城市} & \zh{分类试点}, \zh{回收商}, \zh{拾荒者} \\
        \hline
        \end{tabular}
        \end{center}
    \item \textbf{Préparer l'argumentaire (Débat/Oral) :} Utiliser \zh{虽然……，但是……} (Bien que..., cependant...), \zh{相比之下} (En comparaison), \zh{值得借鉴} (Vaut la peine d'être emprunté comme modèle).
    \item \textbf{Respect et nuance :} Éviter les généralités (« La Chine est propre, le Cameroun est sale »). Parler de « stades de développement différents », « densité urbaine », « moyens financiers », « rôle du secteur informel au Cameroun comme maillon essentiel manquant en Chine ».
\end{enumerate}
\end{methode}

\begin{exoresolu}[Expression orale guidée : Exposé comparatif 1-2 min]
\textbf{Énoncé :} Préparez un exposé oral de 90 secondes : « La gestion des déchets au Cameroun et en Chine : défis et perspectives ». Utilisez au moins 8 mots du vocabulaire de la leçon et la structure \zh{虽然……，但是……}.

\tcblower

\textbf{Proposition de script (Modèle oral) :}

\begin{textebox}[Script expos oral : Gestion déchets Cameroun / Chine]
各位同学，大家好。今天我谈谈中喀垃圾管理的比较。\\
中国推行\textbf{垃圾分类强制化}已有几年，\textbf{像}上海、北京等大城市\textbf{回收利用率}很高，\textbf{处理设施}先进。\textbf{虽然}喀麦隆也有环保法律，\textbf{但是}在雅温得和杜阿拉，\textbf{分类投放}还在\textbf{试点阶段}。\\
很多\textbf{可回收物}靠\textbf{拾荒者}收集，送到\textbf{回收站}。\textbf{厨余垃圾}大多直接\textbf{填埋}，缺乏\textbf{堆肥厂}和\textbf{沼气站}。\\
\textbf{相比之下}，中国的\textbf{无废城市}模式\textbf{值得我们借鉴}。喀麦隆应加强\textbf{基础设施}建设，提高居民\textbf{环保意识}，\textbf{减少填埋}，\textbf{实现资源循环}。谢谢大家。

\textit{Pinyin :}
Gèwèi tóngxué, dàjiā hǎo. Jīntiān wǒ tántan Zhōng-Kā lājī guǎnlǐ de bǐjiào.
Zhōngguó tuīxíng \textbf{lājī fēnlèi qiángzhìhuà} yǐyǒu jǐ nián, \textbf{xiàng} Shànghǎi, Běijīng děng dà chéngshì \textbf{huíshōu lìyònglǜ} hěn gāo, \textbf{chǔlǐ shèshī} xiānxian. \textbf{Suīrán} Kāmàilóng yě yǒu huánbǎo fǎlǜ, \textbf{dànshì} zài Yǎwēndé hé Dūā'ěr, \textbf{fēnlèi tóufàng} hái zài \textbf{shìdiǎn jiēduàn}.
Hěnduō \textbf{kě huíshōu wù} kào \textbf{shíhuāngzhě} shōují, sòng dào \textbf{huíshōuzhàn}. \textbf{Chúyú lājī} duō zhíjiē \textbf{tiántáng}, quēfá \textbf{duīféi chǎng} hé \textbf{zhàoqì zhàn}.
\textbf{Xiāngbǐ zhīxià}, Zhōngguó de \textbf{wúfèi chéngshì} móshì \textbf{zhíde wǒmen jièjiàn}. Kāmàilóng yīng jiāqiáng \textbf{jīchǔ shèshī} jiànshè, tígāo jūmín \textbf{huánbǎo yìshí}, \textbf{jiǎnshǎo tiántáng}, \textbf{shíxiàn zīyuán xúnhuán}. Xièxie dàjiā.

\textit{Traduction :}
Chers camarades, bonjour. Aujourd'hui je compare la gestion des déchets Sino-Camerounaise.
La Chine applique le tri obligatoire depuis quelques années, \textbf{comme} Shanghai, Pékin... taux de recyclage élevé, installations avancées. \textbf{Bien que} le Cameroun ait des lois environnementales, \textbf{mais} à Yaoundé/Douala, le tri est encore au \textbf{stade pilote}.
Beaucoup de \textbf{recyclables} dépendent des \textbf{récupérateurs informels}, envoyés au \textbf{centre de tri}. \textbf{Déchets cuisine} mostly direct \textbf{enfouissement}, manque \textbf{usines compostage} et \textbf{biogaz}.
\textbf{En comparaison}, le modèle chinois \textbf{villes sans déchets} \textbf{vaut la peine d'être emprunté}. Le Cameroun doit renforcer \textbf{infrastructures}, sensibiliser, \textbf{réduire enfouissement}, \textbf{réaliser économie circulaire}. Merci.
\end{textebox}

\textbf{Vocabulaire de la leçon utilisé (9 items) :} \zh{垃圾分类, 回收利用, 处理, 可回收物, 回收站, 厨余垃圾, 填埋, 堆肥厂, 沼气站, 资源循环 (循环)}. \textit{Structure \zh{虽然...但是...} utilisée 1 fois. Connecteurs \zh{相比之下} (en comparaison), \zh{像} (comme).}
\end{exoresolu}

\begin{attention}[Les 5 erreurs qui coûtent des points au Bac - Leçon 9 : Recyclage]
\begin{enumerate}
    \item \textbf{Tons \& Homophones : \zh{熔} (róng, 2ème ton, feu \zh{火} = fondre) $\neq$ \zh{荣} (róng, 2ème ton, gloire) $\neq$ \zh{容} (róng, 2ème ton, contenir/tolérer). \zh{粉} (fěn, 3ème ton, poudre) $\neq$ \zh{份} (fèn, 4ème ton, part/portion). \zh{堆} (duī, 1er ton, tas) $\neq$ \zh{对} (duì, 4ème ton, correct/vers). \textit{Astuce :} Associer le ton au radical (\zh{火} $\to$ action vive/soudaine souvent 2ème ou 4ème ; \zh{米} $\to$ matière fine souvent 3ème).}
    \item \textbf{Caractères confondus (Graphies proches) :} \zh{分} (fēn, diviser, 8 traits, haut \zh{八}) vs \zh{份} (fèn, part, 6 traits, haut \zh{亻}). \zh{类} (lèi, catégorie, 18 traits, clé \zh{米}) vs \zh{米} (mǐ, riz, 6 traits). \zh{回} (huí, retour, 6 traits, cadre \zh{囗}) vs \zh{囗} (wéi, enclos, 3 traits). \zh{收} (shōu, recevoir, 6 traits, clé \zh{攵}) vs \zh{改} (gǎi, changer, 7 traits, clé \zh{攵} + \zh{己}). \textit{Règle :} Écrire en grand, dire l'ordre des traits à voix haute.
    \item \textbf{Ordre des mots : Complément de lieu AVANT le verbe.} \textbf{Faux :} \zh{工人运送垃圾到回收站。} \textbf{Juste :} \zh{工人把垃圾运送\textbf{到回收站}。} ou \zh{工人\textbf{把垃圾运到回收站去}。} Le lieu de destination suit le verbe (ou \zh{把} OBJ V \zh{到} LIEU). Ne jamais mettre le lieu à la fin comme en français.
    \item \textbf{Place des compléments (Temps/Manière) :} \textbf{Faux :} \zh{我昨天分类垃圾。} \textbf{Juste :} \zh{我\textbf{昨天}把垃圾分类了。} (Temps avant \zh{把}/Verbe). \textbf{Faux :} \zh{塑料被熔化快。} \textbf{Juste :} \zh{塑料\textbf{很快}被熔化。} (Manière/Adverbe avant \zh{被}/Verbe).
    \item \textbf{Mots-outils \& Classificateurs :} Oublier \zh{了} (accompli) dans un récit de processus passé (\zh{分类了}). Confondre \zh{个} (général) et \zh{类} (catégorie) : \zh{四\textbf{类}垃圾} (4 catégories) vs \zh{四\textbf{个}垃圾桶} (4 poubelles). Calque « \zh*{把垃圾分类到四类} » $\to$ \textbf{Juste :} \zh{把垃圾\textbf{分成/分为}四类}. Verbe \zh{分} appelle \zh{成/为} pour le résultat de la division.
\end{enumerate}
\end{attention}

\newpage

\coursec{Exercices}

\courssub{Je vérifie mes connaissances}

\begin{exercice}[QCM : Vocabulaire du recyclage]
Chaque phrase ne comporte qu'une seule réponse correcte. Recopiez la phrase complète en choisissant la bonne option.
\begin{enumerate}
    \item 我们应该把垃圾进行 \zh{分类} ，以便 \zh{回收} 利用。
    \begin{itemize}
        \item A. 处理
        \item B. 分类
        \item C. 清洗
        \item D. 加工
    \end{itemize}
    \item 塑料瓶属于 \zh{可回收垃圾} ，不能随意丢弃。
    \begin{itemize}
        \item A. 有害垃圾
        \item B. 厨余垃圾
        \item C. 可回收垃圾
        \item D. 其他垃圾
    \end{itemize}
    \item 回收一吨废纸可以 \zh{节约} 很多木材和水资源。
    \begin{itemize}
        \item A. 浪费
        \item B. 污染
        \item C. 保护
        \item D. 节约
    \end{itemize}
    \item 经过 \zh{清洗} 和加工，旧塑料可以变成新的 \zh{产品} 。
    \begin{itemize}
        \item A. 收集
        \item B. 清洗
        \item C. 丢弃
        \item D. 埋葬
    \end{itemize}
    \item 保护 \zh{环境} 是我们每一个人的责任。
    \begin{itemize}
        \item A. 资源
        \item B. 环境
        \item C. 工厂
        \item D. 过程
    \end{itemize}
\end{enumerate}
\end{exercice}

\begin{exercice}[Vrai ou Faux : Compréhension du processus]
Lisez les affirmations suivantes. Si l'affirmation est correcte, écrivez \textbf{V} (真) ; si elle est fausse, écrivez \textbf{F} (假) et corrigez-la en chinois (caractères + pinyin).
\begin{enumerate}
    \item Le processus de recyclage commence par l'étape de \zh{加工} (jiāgōng).
    \item Le \zh{分类} (fēnlèi) consiste à séparer les déchets selon leur matière (papier, plastique, verre).
    \item Les bouteilles en plastique sales peuvent être envoyées directement à l'usine sans être lavées.
    \item Le recyclage permet de \zh{节约} (jiéyuē) les ressources naturelles et de \zh{保护} (bǎohù) l'environnement.
    \item La dernière étape du processus est la \zh{收集} (shōují) des déchets.
\end{enumerate}
\end{exercice}

\begin{exercice}[Vocabulaire : Associer le mot à sa définition]
Reliez chaque mot du vocabulaire (colonne de gauche) à sa définition correcte (colonne de droite) en écrivant le numéro correspondant.
\begin{center}
\begin{tabularx}{\linewidth}{|X|X|}
\hline
\rowcolor{popblueL} \textbf{Mot (Caractères / Pinyin)} & \textbf{Définition} \\ \hline
1. 回收  & A. Matière transparente, fragile, recyclable à l'infini (verre). \\ \hline
2. 分类  & B. Action de rassembler des déchets dispersés. \\ \hline
3. 资源  & C. Traiter des déchets usagés pour en faire de nouvelles matières premières. \\ \hline
4. 塑料  & D. Matières premières fournies par la nature (bois, eau, minerais). \\ \hline
5. 玻璃  & E. Matière synthétique dérivée du pétrole (plastique). \\ \hline
6. 收集  & F. Séparer les déchets en catégories différentes. \\ \hline
\end{tabularx}
\end{center}
\end{exercice}

\begin{exercice}[Pinyin et Tons : Associer aux caractères]
Écrivez le pinyin complet avec les marques de tons (ex: \zh{huíshōu}) pour chaque caractère ou mot ci-dessous. Attention aux changements de ton (沙 hi) pour \zh{一} et \zh{不} si applicable.
\begin{enumerate}
    \item 垃圾
    \item 环境
    \item 纸
    \item 保护
    \item 清洗
    \item 再利用
    \item 过程
    \item 步骤
\end{enumerate}
\end{exercice}

\courssub{Je m'entraîne}

\begin{exercice}[Traduction : Version (Chinois $\to$ Français)]
Traduisez les phrases suivantes en français naturel. Respectez le temps et l'aspect.
\begin{enumerate}
    \item 我们每天都要把垃圾分类。
    \item 塑料瓶经过清洗和加工，可以变成新产品。
    \item 回收废纸能够节约森林资源。
    \item 请不要把电池扔进普通垃圾桶，电池是有害垃圾。
    \item 保护环境，从我做起，从现在做起。
\end{enumerate}
\end{exercice}

\begin{exercice}[Transformation : Mettre au passif avec \zh{被}  ou \zh{把} ]
Transformez les phrases actives en utilisant la structure \zh{被}  (passif formel/écrit) ou \zh{把}  (disposition active) selon l'indication. Conservez le sens.
\begin{enumerate}
    \item (用 \zh{被}) 工人们收集了垃圾。 $\to$ \dots
    \item (用 \zh{把}) 我们分类了垃圾。 $\to$ \dots
    \item (用 \zh{被}) 工厂加工了塑料。 $\to$ \dots
    \item (用 \zh{把}) 他们清洗了瓶子。 $\to$ \dots
    \item (用 \zh{被}) 大家保护了环境。 $\to$ \dots
\end{enumerate}
\end{exercice}

\begin{exercice}[Texte à trous : Les étapes du recyclage]
Complétez le texte suivant avec les mots de la liste ci-dessous. Écrivez les \textbf{caractères chinois} correspondants. Chaque mot s'utilise une seule fois.
\zh{收集、分类、清洗、加工、再利用、资源、环境、节约}

\begin{textebox}[Processus de recyclage]
回收的过程通常分为几个 \zh{步骤} 。首先是 \zh{\_\_\_\_\_\_} ，把可回收垃圾集中起来；然后是 \zh{\_\_\_\_\_\_} ，把纸、塑料、玻璃分开；接着是 \zh{\_\_\_\_\_\_} ，去除污垢；之后送到工厂 \zh{\_\_\_\_\_\_} 成新原料；最后是 \zh{\_\_\_\_\_\_} ，制成新 \zh{产品} 。这样可以 \zh{\_\_\_\_\_\_} 自然 \zh{\_\_\_\_\_\_} ，也能 \zh{保护} \zh{\_\_\_\_\_\_} 。
\end{textebox}
\end{exercice}

\begin{exercice}[Remise en ordre : Chronologie et connecteurs]
Remettez les cinq étapes du recyclage dans l'ordre logique (1 à 5). Ensuite, rédigez \textbf{une phrase complète en chinois par étape} en utilisant les connecteurs temporels \zh{首先} 、\zh{然后} 、\zh{接着} 、\zh{之后} 、\zh{最后} .
\begin{center}
\begin{tabular}{|c|c|}
\hline
\rowcolor{popblueL} \textbf{Étape (mots mélangés)} & \textbf{Ordre \& Phrase} \\ \hline
加工 & 1. \dots \\ \hline
收集 & 2. \dots \\ \hline
再利用 & 3. \dots \\ \hline
分类 & 4. \dots \\ \hline
清洗 & 5. \dots \\ \hline
\end{tabular}
\end{center}
\end{exercice}

\courssub{Je me prépare au Bac}

\begin{exercice}[Type Bac — Compréhension écrite : Le recyclage du papier]
Lisez le texte inédit ci-dessous (environ 8 lignes). Répondez aux 5 questions \textbf{en chinois} (caractères + pinyin). Les réponses doivent être des phrases complètes.

\begin{textebox}[Le recyclage du papier : un geste pour la forêt]
纸是我们生活中常用的东西。造纸需要砍伐很多树木，消耗大量的水和电。如果我们把用过的纸扔掉，不仅浪费资源，还会污染环境。其实，废纸可以回收再利用。回收纸的过程是这样的：首先，收集废纸；然后，分类，把报纸、纸箱、办公纸分开；接着，送到造纸厂，用水清洗，去掉墨水和胶水，变成纸浆；最后，把纸浆压干、压平，就做成了再生纸。回收一吨废纸，可以节约木材17棵，节约水100吨，减少污染。保护森林，从回收一张纸开始。
\par
\textbf{Pinyin :} Zhǐ shì wǒmen shēnghuó zhōng chángyòng de dōngxi. Zàozhǐ xūyào kǎnfá hěnduō shùmù, xiāohào dàliàng de shuǐ hé diàn. Rúguǒ wǒmen bǎ yòngguò de zhǐ rēngdiào, bùjǐn làngfèi zīyuán, hái huì wūrǎn huánjìng. Qíshí, fèizhǐ kěyǐ huíshōu zàilìyòng. Huíshōu zhǐ de guòchéng shì zhèyàng de: shǒuxiān, shōují fèizhǐ; ránhòu, fēnlèi, bǎ bàozhǐ, zhǐxiāng, bàngōng zhǐ fēnkāi; jiēzhe, sòngdào zàozhǐ chǎng, yòng shuǐ qīngxǐ, qùdiào mòshuǐ hé jiāoshuǐ, biànchéng zhǐjiāng; zuìhòu, bǎ zhǐjiāng yāgān, yāpíng, jiù zuòchéng le zàishēng zhǐ. Huíshōu yī dūn fèizhǐ, kěyǐ jiéyuē mùcái 17 kē, jiéyuē shuǐ 100 dūn, jiǎnshǎo wūrǎn. Bǎohù sēnlín, cóng huíshōu yī zhāng zhǐ kāishǐ.
\par
\textbf{Traduction :} Le papier est un objet courant dans notre vie. La fabrication du papier nécessite d'abattre beaucoup d'arbres, consommant beaucoup d'eau et d'électricité. Si nous jetons le papier usagé, non seulement nous gaspillons des ressources, mais nous polluons aussi l'environnement. En fait, le papier usagé peut être recyclé. Le processus est le suivant : d'abord, collecter le papier usagé ; ensuite, trier, séparer journaux, cartons, papier de bureau ; puis, l'envoyer à la papeterie, le laver à l'eau, enlever l'encre et la colle, en faire de la pâte à papier ; enfin, presser la pâte pour la sécher et l'aplatir, on obtient du papier recyclé. Recycler une tonne de papier usagé permet d'économiser 17 arbres, 100 tonnes d'eau, et de réduire la pollution. Protéger la forêt commence par recycler une feuille de papier.
\end{textebox}

\begin{enumerate}
    \item 造纸需要消耗哪些自然资源？ (Citez-en trois d'après le texte).
    \item 如果乱扔废纸，会有什么后果？ (Deux conséquences).
    \item 请用中文概括回收纸的四个主要步骤 (用 \zh{首先、然后、接着、最后} 连接).
    \item 回收一吨废纸能节约多少棵树木和多少吨水？
    \item 根据课文最后一句，你认为“从我做起”具体可以做什么？(Donnez un exemple concret personnel).
\end{enumerate}
\end{exercice}

\begin{exercice}[Type Bac — Traduction : Version et Thème]
\begin{enumerate}
    \item \textbf{Version (Chinois $\to$ Français) :} Traduisez en français soigné :
    \zh{塑料、纸和玻璃被分类、清洗，再送到工厂加工，变成新的产品。}
    \item \textbf{Thème (Français $\to$ Chinois) :} Traduisez en chinois (caractères + pinyin). Utilisez le vocabulaire de la leçon (\zh{回收、塑料瓶、为了、保护、环境}).
    \emph{« Nous devons recycler les bouteilles en plastique pour protéger l'environnement. »}
\end{enumerate}
\end{exercice}

\begin{exercice}[Type Bac — Expression écrite : Lettre de sensibilisation]
\textbf{Consigne :} Vous êtes \zh{李华} (Lǐ Huá), délégué écolo de votre lycée à Yaoundé. Vous écrivez une lettre ouverte (environ \textbf{80 à 100 caractères chinois}) à vos camarades pour les inciter à participer à la semaine du recyclage (\zh{回收周}).

\textbf{Points obligatoires à inclure :}
\begin{itemize}
    \item L'importance du tri (\zh{分类}) et du recyclage (\zh{回收}) pour \zh{节约资源} et \zh{保护环境}.
    \item Citer au moins 3 types de déchets recyclables (\zh{纸、塑料瓶、玻璃、金属}...).
    \item Appeler à l'action concrète : \zh{把垃圾扔进对应的垃圾桶} / \zh{参加学校的回收活动}.
    \item Formules de politesse épistolaire adaptées (\zh{亲爱的同学们：} ... \zh{李华} \zh{202x年x月x日}).
\end{itemize}

\begin{textebox}[Votre production]
(Espace réservé pour la rédaction du candidat)
\end{textebox}
\end{exercice}

\newpage

\coursec{Corrigés des exercices}

\begin{corrige}[Exercice 1 — QCM : Vocabulaire du recyclage]
\begin{enumerate}
    \item \textbf{B. 分类} — 我们应该把垃圾进行分类，以便回收利用。\\ \emph{Wǒmen yīnggāi bǎ lājī jìnxíng fēnlèi, yǐbiàn huíshōu lìyòng.}\\ « Nous devons trier les déchets afin de pouvoir les recycler et les réutiliser. » Justification : \zh{分类} (trier) est l'action préalable au recyclage ; \zh{处理} (traiter) serait possible grammaticalement mais ne correspond pas au sens « afin de recycler ».
    \item \textbf{C. 可回收垃圾} — 塑料瓶属于可回收垃圾，不能随意丢弃。\\ \emph{Sùliào píng shǔyú kě huíshōu lājī, bù néng suíyì diūqì.}\\ « Les bouteilles en plastique appartiennent aux déchets recyclables, on ne peut pas les jeter n'importe où. »
    \item \textbf{D. 节约} — 回收一吨废纸可以节约很多木材和水资源。\\ \emph{Huíshōu yī dūn fèizhǐ kěyǐ jiéyuē hěnduō mùcái hé shuǐ zīyuán.}\\ « Recycler une tonne de papier usagé permet d'économiser beaucoup de bois et de ressources en eau. » \zh{浪费} (gaspiller) est l'antonyme : sens opposé.
    \item \textbf{B. 清洗} — 经过清洗和加工，旧塑料可以变成新的产品。\\ \emph{Jīngguò qīngxǐ hé jiāgōng, jiù sùliào kěyǐ biànchéng xīn de chǎnpǐn.}\\ « Après lavage et transformation, le vieux plastique peut devenir de nouveaux produits. »
    \item \textbf{B. 环境} — 保护环境是我们每一个人的责任。\\ \emph{Bǎohù huánjìng shì wǒmen měi yí ge rén de zérèn.}\\ « Protéger l'environnement est la responsabilité de chacun d'entre nous. »
\end{enumerate}
\end{corrige}

\begin{corrige}[Exercice 2 — Vrai ou Faux : Compréhension du processus]
\begin{enumerate}
    \item \textbf{F} (假). Correction : 回收的过程首先从收集开始。\\ \emph{Huíshōu de guòchéng shǒuxiān cóng shōují kāishǐ.}\\ « Le processus de recyclage commence d'abord par la collecte. » L'\zh{加工} (transformation) est une étape ultérieure, en usine.
    \item \textbf{V} (真). Le \zh{分类} consiste bien à séparer les déchets selon leur matière : papier, plastique, verre, métal.
    \item \textbf{F} (假). Correction : 脏的塑料瓶要先清洗干净，才能送到工厂加工。\\ \emph{Zāng de sùliào píng yào xiān qīngxǐ gānjìng, cái néng sòngdào gōngchǎng jiāgōng.}\\ « Les bouteilles en plastique sales doivent d'abord être lavées proprement avant de pouvoir être envoyées à l'usine pour transformation. » Remarquez la structure \zh{先……才……} (d'abord\dots seulement ensuite\dots).
    \item \textbf{V} (真). C'est exactement le double bénéfice du recyclage : \zh{节约资源} (économiser les ressources) et \zh{保护环境} (protéger l'environnement).
    \item \textbf{F} (假). Correction : 回收过程的最后一步是再利用，制成新产品。\\ \emph{Huíshōu guòchéng de zuìhòu yī bù shì zàilìyòng, zhìchéng xīn chǎnpǐn.}\\ « La dernière étape du processus de recyclage est la réutilisation, pour fabriquer de nouveaux produits. » La \zh{收集} est la première étape, pas la dernière.
\end{enumerate}
\end{corrige}

\begin{corrige}[Exercice 3 — Vocabulaire : Associer le mot à sa définition]
\begin{center}
\begin{tabularx}{\linewidth}{|c|X|X|}
\hline
\rowcolor{popblueL} \textbf{Mot} & \textbf{Bonne définition} & \textbf{Rappel} \\ \hline
1. 回收 \emph{huíshōu} & \textbf{C} & Traiter des déchets usagés pour en faire de nouvelles matières premières. \\ \hline
2. 分类 \emph{fēnlèi} & \textbf{F} & Séparer les déchets en catégories différentes. \\ \hline
3. 资源 \emph{zīyuán} & \textbf{D} & Matières premières fournies par la nature (bois, eau, minerais). \\ \hline
4. 塑料 \emph{sùliào} & \textbf{E} & Matière synthétique dérivée du pétrole (plastique). \\ \hline
5. 玻璃 \emph{bōli} & \textbf{A} & Matière transparente, fragile, recyclable à l'infini (verre). \\ \hline
6. 收集 \emph{shōují} & \textbf{B} & Action de rassembler des déchets dispersés. \\ \hline
\end{tabularx}
\end{center}
Réponses : 1-C ; 2-F ; 3-D ; 4-E ; 5-A ; 6-B.\\
\textbf{Point de vigilance :} ne pas confondre \zh{回收} (recycler, récupérer) et \zh{收集} (collecter, ramasser) : \zh{收集} est la première étape matérielle, \zh{回收} désigne tout le processus de valorisation.
\end{corrige}

\begin{corrige}[Exercice 4 — Pinyin et Tons : Associer aux caractères]
\begin{enumerate}
    \item 垃圾 — \emph{lājī} (1\textsuperscript{er} ton + 1\textsuperscript{er} ton)
    \item 环境 — \emph{huánjìng} (2\textsuperscript{e} ton + 4\textsuperscript{e} ton)
    \item 纸 — \emph{zhǐ} (3\textsuperscript{e} ton ; attention, \zh{zh-} rétroflexe, pas « z »)
    \item 保护 — \emph{bǎohù} (3\textsuperscript{e} ton + 4\textsuperscript{e} ton)
    \item 清洗 — \emph{qīngxǐ} (1\textsuperscript{er} ton + 3\textsuperscript{e} ton)
    \item 再利用 — \emph{zàilìyòng} (4\textsuperscript{e} + 4\textsuperscript{e} + 4\textsuperscript{e} ton) ; on écrit aussi \zh{再使用}. Pas de changement de ton ici.
    \item 过程 — \emph{guòchéng} (4\textsuperscript{e} ton + 2\textsuperscript{e} ton)
    \item 步骤 — \emph{bùzhòu} : attention à la règle de \zh{不} ! Devant un 4\textsuperscript{e} ton (\zh{骤} \emph{zhòu}), \zh{不} passe au 2\textsuperscript{e} ton : on prononce \emph{búzhòu}, mais on écrit le pinyin lexical \emph{bùzhòu}.
\end{enumerate}
\textbf{Point de vigilance :} deux 3\textsuperscript{e} tons qui se suivent entraînent un changement de ton (le premier devient 2\textsuperscript{e} ton à l'oral), mais aucun mot de cette liste ne présente ce cas. En revanche, retenez la règle de \zh{不} : \zh{不} + 4\textsuperscript{e} ton $\to$ \emph{bú}.
\end{corrige}

\begin{corrige}[Exercice 5 — Traduction : Version (Chinois $\to$ Français)]
\begin{enumerate}
    \item 我们每天都要把垃圾分类。\\ \emph{Wǒmen měitiān dōu yào bǎ lājī fēnlèi.}\\ « Nous devons trier les déchets tous les jours. » Structure \zh{把} + objet + verbe : \zh{把垃圾分类} = « trier les déchets » (litt. « prendre les déchets et les trier »).
    \item 塑料瓶经过清洗和加工，可以变成新产品。\\ \emph{Sùliào píng jīngguò qīngxǐ hé jiāgōng, kěyǐ biànchéng xīn chǎnpǐn.}\\ « Après être passées par le lavage et la transformation, les bouteilles en plastique peuvent devenir de nouveaux produits. » \zh{经过} = « après avoir traversé / par le biais de ».
    \item 回收废纸能够节约森林资源。\\ \emph{Huíshōu fèizhǐ nénggòu jiéyuē sēnlín zīyuán.}\\ « Recycler le papier usagé permet d'économiser les ressources forestières. »
    \item 请不要把电池扔进普通垃圾桶，电池是有害垃圾。\\ \emph{Qǐng búyào bǎ diànchí rēng jìn pǔtōng lājī tǒng, diànchí shì yǒuhài lājī.}\\ « S'il vous plaît, ne jetez pas les piles dans la poubelle ordinaire : les piles sont des déchets dangereux. » Notez \zh{不要} prononcé \emph{búyào} (changement de ton de \zh{不}).
    \item 保护环境，从我做起，从现在做起。\\ \emph{Bǎohù huánjìng, cóng wǒ zuòqǐ, cóng xiànzài zuòqǐ.}\\ « Protéger l'environnement : cela commence par moi, cela commence maintenant. » Slogan classique : \zh{从……做起} = « commencer à agir à partir de\dots ».
\end{enumerate}
\end{corrige}

\begin{corrige}[Exercice 6 — Transformation : passif avec 被 / disposition avec 把]
\begin{enumerate}
    \item 垃圾被工人们收集了。\\ \emph{Lājī bèi gōngrénmen shōují le.}\\ « Les déchets ont été collectés par les ouvriers. » Structure : Objet + \zh{被} + agent + verbe + \zh{了}.
    \item 我们把垃圾分类了。\\ \emph{Wǒmen bǎ lājī fēnlèi le.}\\ « Nous avons trié les déchets. » Structure : Sujet + \zh{把} + objet + verbe + complément (\zh{了}).
    \item 塑料被工厂加工了。\\ \emph{Sùliào bèi gōngchǎng jiāgōng le.}\\ « Le plastique a été transformé par l'usine. »
    \item 他们把瓶子清洗了。\\ \emph{Tāmen bǎ píngzi qīngxǐ le.}\\ « Ils ont lavé les bouteilles. » Variante plus naturelle avec résultatif : \zh{他们把瓶子清洗干净了。} (\emph{Tāmen bǎ píngzi qīngxǐ gānjìng le.}) « Ils ont lavé les bouteilles jusqu'à ce qu'elles soient propres. »
    \item 环境被大家保护了。\\ \emph{Huánjìng bèi dàjiā bǎohù le.}\\ « L'environnement a été protégé par tous. »
\end{enumerate}
\textbf{Point de vigilance :} avec \zh{把} et \zh{被}, le verbe ne peut JAMAIS rester seul : il faut un élément après lui (\zh{了}, un complément de résultat, de direction, etc.). On ne dit pas \zh{把垃圾分}. De plus, l'objet de \zh{把}/\zh{被} doit être défini (connu des interlocuteurs).
\end{corrige}

\begin{corrige}[Exercice 7 — Texte à trous : Les étapes du recyclage]
Texte complété (les mots attendus sont en gras) :

回收的过程通常分为几个步骤。首先是\textbf{收集}，把可回收垃圾集中起来；然后是\textbf{分类}，把纸、塑料、玻璃分开；接着是\textbf{清洗}，去除污垢；之后送到工厂\textbf{加工}成新原料；最后是\textbf{再利用}，制成新产品。这样可以\textbf{节约}自然\textbf{资源}，也能保护\textbf{环境}。

\emph{Huíshōu de guòchéng tōngcháng fēn wéi jǐ ge bùzhòu. Shǒuxiān shì shōují, bǎ kě huíshōu lājī jízhōng qǐlai; ránhòu shì fēnlèi, bǎ zhǐ, sùliào, bōli fēnkāi; jiēzhe shì qīngxǐ, qùchú wūgòu; zhīhòu sòngdào gōngchǎng jiāgōng chéng xīn yuánliào; zuìhòu shì zàilìyòng, zhìchéng xīn chǎnpǐn. Zhèyàng kěyǐ jiéyuē zìrán zīyuán, yě néng bǎohù huánjìng.}

« Le processus de recyclage se divise généralement en plusieurs étapes. D'abord la collecte, qui rassemble les déchets recyclables ; ensuite le tri, qui sépare le papier, le plastique et le verre ; puis le lavage, qui enlève les saletés ; après quoi on les envoie à l'usine pour être transformés en nouvelles matières premières ; enfin la réutilisation, qui fabrique de nouveaux produits. Ainsi, on peut économiser les ressources naturelles et aussi protéger l'environnement. »

\textbf{Point de vigilance :} l'ordre logique est imposé par le texte lui-même (collecte $\to$ tri $\to$ lavage $\to$ transformation $\to$ réutilisation). \zh{节约资源} et \zh{保护环境} sont des collocations fixes à mémoriser ensemble.
\end{corrige}

\begin{corrige}[Exercice 8 — Remise en ordre : Chronologie et connecteurs]
Ordre correct : 1. 收集 $\to$ 2. 分类 $\to$ 3. 清洗 $\to$ 4. 加工 $\to$ 5. 再利用.

Phrases modèles (une par étape, avec connecteur) :
\begin{enumerate}
    \item 首先，我们收集可回收的垃圾。\\ \emph{Shǒuxiān, wǒmen shōují kě huíshōu de lājī.}\\ « D'abord, nous collectons les déchets recyclables. »
    \item 然后，我们把垃圾分类，把纸、塑料和玻璃分开。\\ \emph{Ránhòu, wǒmen bǎ lājī fēnlèi, bǎ zhǐ, sùliào hé bōli fēnkāi.}\\ « Ensuite, nous trions les déchets, en séparant le papier, le plastique et le verre. »
    \item 接着，我们清洗脏的瓶子和盒子。\\ \emph{Jiēzhe, wǒmen qīngxǐ zāng de píngzi hé hézi.}\\ « Puis, nous lavons les bouteilles et les boîtes sales. »
    \item 之后，工厂把这些材料加工成新的原料。\\ \emph{Zhīhòu, gōngchǎng bǎ zhèxiē cáiliào jiāgōng chéng xīn de yuánliào.}\\ « Après cela, l'usine transforme ces matériaux en nouvelles matières premières. »
    \item 最后，我们再利用这些原料，制成新的产品。\\ \emph{Zuìhòu, wǒmen zàilìyòng zhèxiē yuánliào, zhìchéng xīn de chǎnpǐn.}\\ « Enfin, nous réutilisons ces matières premières pour fabriquer de nouveaux produits. »
\end{enumerate}
\textbf{Point de vigilance :} les connecteurs se placent en tête de phrase, suivis d'une virgule chinoise \zh{，}. Ne pas confondre \zh{之后} (après cela) et \zh{最后} (finalement) : \zh{最后} clôt obligatoirement la série.
\end{corrige}

\begin{corrige}[Exercice 9 — Type Bac : Compréhension écrite (Le recyclage du papier)]
\textbf{Barème indicatif : 5 points (1 point par question : 0,5 pour le sens, 0,5 pour la correction de la langue).}

\begin{enumerate}
    \item 造纸需要消耗树木（木材）、水和电。\\ \emph{Zàozhǐ xūyào xiāohào shùmù (mùcái), shuǐ hé diàn.}\\ « La fabrication du papier consomme des arbres (du bois), de l'eau et de l'électricité. »
    \item 如果乱扔废纸，不仅浪费资源，还会污染环境。\\ \emph{Rúguǒ luàn rēng fèizhǐ, bùjǐn làngfèi zīyuán, hái huì wūrǎn huánjìng.}\\ « Si l'on jette le papier usagé n'importe où, non seulement on gaspille des ressources, mais on pollue aussi l'environnement. » La corrélation \zh{不仅……还……} est attendue pour citer les deux conséquences.
    \item 回收纸的步骤是：首先，收集废纸；然后，分类，把报纸、纸箱、办公纸分开；接着，送到造纸厂清洗，去掉墨水和胶水，变成纸浆；最后，把纸浆压干、压平，做成再生纸。\\ \emph{Huíshōu zhǐ de bùzhòu shì: shǒuxiān, shōují fèizhǐ; ránhòu, fēnlèi, bǎ bàozhǐ, zhǐxiāng, bàngōng zhǐ fēnkāi; jiēzhe, sòngdào zàozhǐ chǎng qīngxǐ, qùdiào mòshuǐ hé jiāoshuǐ, biànchéng zhǐjiāng; zuìhòu, bǎ zhǐjiāng yāgān, yāpíng, zuòchéng zàishēng zhǐ.}\\ « Les étapes : d'abord collecter le papier usagé ; ensuite trier (journaux, cartons, papier de bureau) ; puis envoyer à la papeterie pour lavage, enlever l'encre et la colle, obtenir de la pâte ; enfin presser et aplatir la pâte pour obtenir du papier recyclé. »
    \item 回收一吨废纸可以节约17棵树木和100吨水。\\ \emph{Huíshōu yī dūn fèizhǐ kěyǐ jiéyuē 17 kē shùmù hé 100 dūn shuǐ.}\\ « Recycler une tonne de papier usagé permet d'économiser 17 arbres et 100 tonnes d'eau. » Attention au classificateur : \zh{一棵} pour les arbres.
    \item 例：从我做起，我可以把用过的纸放进可回收垃圾桶，不浪费纸。\\ \emph{Lì: cóng wǒ zuòqǐ, wǒ kěyǐ bǎ yòngguò de zhǐ fàng jìn kě huíshōu lājī tǒng, bù làngfèi zhǐ.}\\ « Exemple : en commençant par moi, je peux mettre le papier usagé dans la poubelle de recyclage et ne pas gaspiller le papier. » Toute réponse personnelle cohérente utilisant le vocabulaire de la leçon est acceptée.
\end{enumerate}
\textbf{Points de vigilance :} répondre par des phrases complètes (sujet + verbe), reprendre les mots du texte sans les recopier mot à mot sur plusieurs lignes, soigner les tons dans le pinyin et ne pas oublier le classificateur \zh{棵} devant \zh{树木}.
\end{corrige}

\begin{corrige}[Exercice 10 — Type Bac : Traduction (Version et Thème)]
\textbf{Barème indicatif : 4 points (2 points par traduction : 1 pour le sens, 1 pour la langue).}

\begin{enumerate}
    \item \textbf{Version :} 塑料、纸和玻璃被分类、清洗，再送到工厂加工，变成新的产品。\\ \emph{Sùliào, zhǐ hé bōli bèi fēnlèi, qīngxǐ, zài sòngdào gōngchǎng jiāgōng, biànchéng xīn de chǎnpǐn.}\\ Traduction : « Le plastique, le papier et le verre sont triés, lavés, puis envoyés à l'usine pour être transformés et devenir de nouveaux produits. » Justification : \zh{被} marque le passif (les trois matériaux subissent les actions) ; \zh{再} exprime ici « puis, ensuite » (action suivante), et non « encore une fois ».
    \item \textbf{Thème :} « Nous devons recycler les bouteilles en plastique pour protéger l'environnement. »\\ 我们应该回收塑料瓶，为了保护环境。\\ \emph{Wǒmen yīnggāi huíshōu sùliào píng, wèile bǎohù huánjìng.}\\ Variante tout aussi correcte : 为了保护环境，我们应该回收塑料瓶。\\ \emph{Wèile bǎohù huánjìng, wǒmen yīnggāi huíshōu sùliào píng.}\\ « Pour protéger l'environnement, nous devons recycler les bouteilles en plastique. » Structure : \zh{为了} + but ; \zh{应该} + verbe pour le devoir.
\end{enumerate}
\textbf{Points de vigilance :} \zh{塑料瓶} est un composé inséparable (\zh{塑料} + \zh{瓶}) ; ne pas traduire « pour » par \zh{给} ici — le but se rend par \zh{为了}. En version, rendre le passif \zh{被} par « sont triés, lavés » et non par une tournure active qui changerait le sens.
\end{corrige}

\begin{corrige}[Exercice 11 — Type Bac : Expression écrite (Lettre de sensibilisation)]
\textbf{Barème indicatif : 6 points.} Respect de la forme épistolaire (1 pt) ; nombre de caractères 80--100 (1 pt) ; présence des 4 points obligatoires (2 pts) ; correction grammaticale et vocabulaire de la leçon (1,5 pt) ; clarté et cohérence (0,5 pt).

\textbf{Proposition de rédaction modèle :}

亲爱的同学们：

大家好！下周是我们学校的回收周。回收和分类非常重要：它们可以节约资源，也能保护环境。纸、塑料瓶、玻璃和金属都可以回收再利用。请大家把垃圾扔进对应的垃圾桶，并积极参加学校的回收活动。保护环境，从我做起，从现在做起！

李华\\
2024年5月10日

\emph{Qīn'ài de tóngxuémen:}

\emph{Dàjiā hǎo! Xià zhōu shì wǒmen xuéxiào de huíshōu zhōu. Huíshōu hé fēnlèi fēicháng zhòngyào: tāmen kěyǐ jiéyuē zīyuán, yě néng bǎohù huánjìng. Zhǐ, sùliào píng, bōli hé jīnshǔ dōu kěyǐ huíshōu zàilìyòng. Qǐng dàjiā bǎ lājī rēng jìn duìyìng de lājī tǒng, bìng jījí cānjiā xuéxiào de huíshōu huódòng. Bǎohù huánjìng, cóng wǒ zuòqǐ, cóng xiànzài zuòqǐ!}

\emph{Lǐ Huá}\\\emph{2024 nián 5 yuè 10 rì}

« Chers camarades, bonjour ! La semaine prochaine, c'est la semaine du recyclage de notre lycée. Le recyclage et le tri sont très importants : ils permettent d'économiser les ressources et de protéger l'environnement. Le papier, les bouteilles en plastique, le verre et le métal peuvent tous être recyclés et réutilisés. Je vous demande de jeter vos déchets dans les poubelles correspondantes et de participer activement aux activités de recyclage de l'école. Protéger l'environnement commence par moi, commence maintenant ! »

\textbf{Points de vigilance :}
\begin{itemize}
    \item La formule d'appel \zh{亲爱的同学们：} se termine par deux-points chinois \zh{：} et la salutation \zh{大家好！} est en retrait.
    \item La signature \zh{李华} et la date (format \zh{年} + \zh{月} + \zh{日}) se placent en bas à droite.
    \item Citer au moins trois déchets reliés par \zh{、} (virgule d'énumération chinoise), jamais par une virgule occidentale.
    \item Employer les collocations de la leçon : \zh{节约资源}, \zh{保护环境}, \zh{回收再利用}, \zh{把垃圾扔进垃圾桶}.
    \item Compter les caractères : la ponctuation ne compte généralement pas ; visez 85--95 caractères pour rester dans la fourchette.
\end{itemize}
\end{corrige}

\newpage

\coursec{Fiche bilan}

\begin{popfigure}
\begin{center}\tcbox[colback=popblue,colframe=popblue,coltext=white,arc=9pt,boxsep=4pt,left=6pt,right=6pt]{\bfseries\small Leçon 9 Le processus du recyclage}\end{center}
\vspace{-2pt}
\begin{tcbraster}[raster columns=3,raster equal height=rows,raster column skip=6pt,raster row skip=6pt,enhanced,blankest]
\begin{tcolorbox}[enhanced,colback=poporange,colframe=poporange,coltext=white,boxrule=0.9pt,arc=3pt,left=4pt,right=4pt,top=3pt,bottom=3pt,boxsep=1pt,halign=left]\bfseries\scriptsize Lexique clé 回收 垃圾 分类 处理\end{tcolorbox}
\begin{tcolorbox}[enhanced,colback=popgreen,colframe=popgreen,coltext=white,boxrule=0.9pt,arc=3pt,left=4pt,right=4pt,top=3pt,bottom=3pt,boxsep=1pt,halign=left]\bfseries\scriptsize Étapes du tri 分类 → 收集 → 处理\end{tcolorbox}
\begin{tcolorbox}[enhanced,colback=poppurple,colframe=poppurple,coltext=white,boxrule=0.9pt,arc=3pt,left=4pt,right=4pt,top=3pt,bottom=3pt,boxsep=1pt,halign=left]\bfseries\scriptsize Clés et traits 扌 土 纟 讠\end{tcolorbox}
\begin{tcolorbox}[enhanced,colback=poppink,colframe=poppink,coltext=white,boxrule=0.9pt,arc=3pt,left=4pt,right=4pt,top=3pt,bottom=3pt,boxsep=1pt,halign=left]\bfseries\scriptsize Structures 被 + verbe 把 + objet\end{tcolorbox}
\begin{tcolorbox}[enhanced,colback=popgold,colframe=popgold,coltext=white,boxrule=0.9pt,arc=3pt,left=4pt,right=4pt,top=3pt,bottom=3pt,boxsep=1pt,halign=left]\bfseries\scriptsize Compréhension lire, repérer, répondre\end{tcolorbox}
\begin{tcolorbox}[enhanced,colback=popblueL,colframe=popblueL,coltext=white,boxrule=0.9pt,arc=3pt,left=4pt,right=4pt,top=3pt,bottom=3pt,boxsep=1pt,halign=left]\bfseries\scriptsize Collocations 保护环境 减少污染\end{tcolorbox}
\begin{tcolorbox}[enhanced,colback=popgreenL,colframe=popgreenL,coltext=white,boxrule=0.9pt,arc=3pt,left=4pt,right=4pt,top=3pt,bottom=3pt,boxsep=1pt,halign=left]\bfseries\scriptsize Socio-culture Cameroun / Chine\end{tcolorbox}
\end{tcbraster}

\legende{Carte mentale de la leçon 9 : le processus du recyclage}
\end{popfigure}

\begin{bilan}
\textbf{1. Lexique clé à connaître par cœur}
\begin{center}
\begin{tabularx}{\linewidth}{|X|X|X|X|}
\hline
\rowcolor{popblueL} Caractères & Pinyin & Nature & Traduction \\
\hline
回收 & huíshōu & verbe & recycler, récupérer \\
\hline
垃圾 & lājī & nom & déchets, ordures \\
\hline
分类 & fēnlèi & verbe & trier, classer \\
\hline
处理 & chǔlǐ & verbe & traiter \\
\hline
污染 & wūrǎn & nom / verbe & pollution / polluer \\
\hline
环境 & huánjìng & nom & environnement \\
\hline
瓶子 & píngzi & nom & bouteille \\
\hline
塑料 & sùliào & nom & plastique \\
\hline
节约 & jiéyuē & verbe & économiser \\
\hline
保护 & bǎohù & verbe & protéger \\
\hline
\end{tabularx}
\end{center}

\textbf{2. Règles de grammaire à connaître par cœur}
\begin{center}
\begin{tabularx}{\linewidth}{|X|X|X|}
\hline
\rowcolor{popblueL} Structure & Exemple & Traduction \\
\hline
Sujet + 被 + agent + verbe & 垃圾被工人收集。Lājī bèi gōngrén shōují. & Les déchets sont collectés par les ouvriers. \\
\hline
Sujet + 把 + objet + verbe + complément & 我们把垃圾分类。Wǒmen bǎ lājī fēnlèi. & Nous trions les déchets. \\
\hline
先 …… 然后 …… 最后 …… & 先分类，然后收集，最后处理。 & D'abord trier, ensuite collecter, enfin traiter. \\
\hline
可以 + verbe (possibilité) & 塑料可以回收。Sùliào kěyǐ huíshōu. & Le plastique peut être recyclé. \\
\hline
应该 + verbe (devoir) & 我们应该保护环境。Wǒmen yīnggāi bǎohù huánjìng. & Nous devons protéger l'environnement. \\
\hline
\end{tabularx}
\end{center}

\textbf{3. Clés (radicaux) du thème}
\begin{itemize}
  \item \cle{扌} (clé de la main) : 扔 rēng « jeter », 拾 shí « ramasser » — traits : horizontal, vertical-crochet, puis relevé.
  \item \cle{土} (clé de la terre) : 垃 lā, 圾 jī — la clé 土 s'écrit à gauche, horizontal court, vertical, horizontal long.
  \item \cle{纟} (clé de la soie) : 纸 zhǐ « papier » — trois traits à gauche avant la partie droite.
  \item \cle{讠} (clé de la parole) : 说 shuō, 话 huà — point puis pli horizontal-vertical.
\end{itemize}

\textbf{4. Méthode express : comprendre un texte sur un processus}
\begin{enumerate}
  \item Lire une première fois sans s'arrêter : repérer le thème général.
  \item Souligner les mots d'étapes : \zh{先} xiān, \zh{然后} ránhòu, \zh{最后} zuìhòu.
  \item Repérer les verbes d'action : 分类, 收集, 处理, 回收.
  \item Répondre aux questions en recopiant la structure du texte, en phrases complètes.
  \item Vérifier les tons et l'ordre des mots avant de rendre la copie.
\end{enumerate}
\end{bilan}

\begin{aretenir}[Checklist avant le Bac]
\begin{itemize}
  \item[$\square$] Je sais lire et écrire les dix mots clés du thème (回收, 垃圾, 分类, 处理, 污染, 环境, 瓶子, 塑料, 节约, 保护).
  \item[$\square$] Je prononce correctement les tons de chaque mot (huíshōu : 2\textsuperscript{e} + 1\textsuperscript{er} ton).
  \item[$\square$] Je connais les étapes du recyclage dans l'ordre : 分类 → 收集 → 运输 → 处理 → 再利用.
  \item[$\square$] Je sais construire une phrase passive avec 被.
  \item[$\square$] Je sais construire une phrase avec 把 + objet + verbe.
  \item[$\square$] J'utilise correctement 先 …… 然后 …… 最后 …… pour décrire un processus.
  \item[$\square$] Je distingue 可以 (possibilité) et 应该 (devoir).
  \item[$\square$] Je reconnais les clés 扌, 土, 纟, 讠 et je sais les nommer en français.
  \item[$\square$] Je respecte l'ordre des traits : haut → bas, gauche → droite, horizontal avant vertical.
  \item[$\square$] Je sais répondre à une question de compréhension par une phrase complète en chinois.
  \item[$\square$] Je sais citer un exemple de recyclage au Cameroun (bouteilles en plastique à Douala) et en Chine (tri sélectif dans les grandes villes).
  \item[$\square$] Je vérifie la ponctuation chinoise pleine chasse ：，。？ dans ma production écrite.
\end{itemize}
\end{aretenir}

\findecours

\end{document}
